\documentclass[12pt, ukrainian]{article}
\usepackage[a4paper]{geometry}
\setlength\parindent{0mm}
\geometry{top=1.1cm,bottom=1.1cm,left=1.1cm,right=1.1cm}
\pagestyle{empty}
\usepackage[T2A]{fontenc}
\usepackage[utf8]{inputenc}
\usepackage[ukrainian]{babel}
\usepackage{graphicx}
\usepackage{caption}
\usepackage{verbatim, listings, clrscode3e, fancyvrb}
\DeclareCaptionFormat{nocaption}{}
\captionsetup{format=nocaption,aboveskip=0pt,belowskip=0pt}
\usepackage[Export]{adjustbox}
\adjustboxset{max size={0.9\linewidth}{0.9\paperheight}}
\def\code#1{\Verb!#1!}
\begin{document}
{{23.05.2025 Контрольна робота, варіант  \No \textbf { 1 }} група К1{\parbox {0.5 cm} \dotfill}  ПІБ {\parbox {2.5 cm} \dotfill}\par}
{
%\begin {large}
\begin{minipage}[t]{9.7cm}
               
{\begin{center}Завдання {1}\end{center}}
\vspace{-6mm}
Укажіть ТИП значення виразу.\par\par
\begin{verbatim}
{y:None for y in range(-8,3)}
\end{verbatim}

\vspace{15mm}
У завданнях 2, 3 вказати, що буде виведено за виконання. Повідомлення інтерпретатора про помилку позначати ERROR

               
{\begin{center}Завдання {2}\end{center}}
\vspace{-6mm}

\begin{verbatim}
try:
    t = {80:3, 74:2, 71:8, 55:6, 71:5}
    t[71] = 5
    for x in t.values():
        print(x, sep=' ', end=' ')
except:
    print('ERR')
\end{verbatim}


\end{minipage}
\vline \hspace{8mm}
\begin{minipage}[t]{8.0cm}                
               
{\begin{center}Завдання {3}\end{center}}
\vspace{-6mm}

\begin{verbatim}
c = 0

class C:
    c = 1
    
    def __init__(self):
        self.c = 2
        c = 3
        C.c = 4

obj = C()
obj.c = 5
try:
    print(c, end=' ')
    print(obj.c, end=' ')
    print(C.c, end=' ')
    print(obj.c, end=' ')
    print(obj._C__c)
except:
    print('ERR')
\end{verbatim}

\end{minipage}

               
{\begin{center}Завдання {4}\end{center}}
\vspace{-6mm}
Нехай змінна \kw{s} позначає послідовність цілих, по якій можна ітерувати.

Написати вираз, що перевіряє, чи правда, що послідовність  містить числа, більші за 5.\par

\vspace{10mm}
               
{\begin{center}Завдання {5}\end{center}}
\vspace{-6mm}
Змінні \kw{next} та \kw{prev} вузла списку позначають наступний та попередній за ним вузол відповідно. Починаючи з голови, у список записано послідовні цілі числа від~$-23$ до~$37$. Нехай змінна \kw{p} позначає вузол, що зберігає значення~$-7$.
Яке число зберігається у вузлі \kw{p.prev.next.prev.prev.prev} ?\par

               
{\begin{center}Завдання {6}\end{center}}
\vspace{-6mm}
Значенням змінної \kw{a} є цілочисловий масив типу $numpy.ndarray$ розмірів $4{\times}7$. На скільки збільшиться сума елементів масиву \kw{a} після виконання
\begin{verbatim}
a[0] += 3
\end{verbatim}


Якщо код некоректний, то в якості відповіді зазначити ERROR.\par

               
{\begin{center}Завдання {10}\end{center}}
\vspace{-6mm}

Написати функцію $f$, яка для файлу, заданого шляхом, розв’язує задачу:\\ Текстовий файл містить послідовність чисел $a_1$, $a_2$, $\ldots$, $a_t$, записану у форматі одне число на рядок. Знайти середнє арифметичне чисел $a_1$, $a_4$, $a_7$, $\ldots$ (починаючи з першого через два).

В усіх випадках розміри файлів такі, що їх вміст повністю завантажений у пам'ять бути не може. Якість коду також оцінюється.



\vspace{10mm}\par

\newpage
\begin{minipage}[t]{9.7cm}
               
{\begin{center}Завдання {7}\end{center}}
\vspace{-6mm}

\begin{center}
	\adjustimage{max size={0.8\linewidth}{0.8\paperheight}}
	{./15BinTree/trees_exp/126.png}
\end{center}
Указати послідовність міток вершин дерева, зображеного на рис., що утвориться в результаті обходу дерева в оберненому порядку. Мітки записувати через пробіл.\par Алгоритм починає роботу з кореня (на рис. це верхня вершина).\par

\end{minipage}
\vline \hspace{8mm}
\begin{minipage}[t]{8.0cm}
               
{\begin{center}Завдання {8}\end{center}}
\vspace{-6mm}
Для графа на рисунку з вершини 4 запущено алгоритм пошуку в ширину, що будує кістякове дерево. Списки суміжності вершин переглядаються алгоритмом за зростанням міток вершин.\par
\begin{center}
	\adjustimage{max size={0.9\linewidth}{0.9\paperheight}}
	{./16Graphs/Traversals/121.png}
\end{center}
\par Які з наведених ребер входять до складу отриманого кістякового дерева?\par
1) (0, 3)\par
2) (1, 3)\par
3) (4, 6)\par
4) (2, 4)\par
5) (0, 2)\par
6) (0, 4)\par
{\vskip 15pt} \par

\end{minipage}

               
{\begin{center}Завдання {9}\end{center}}
\vspace{-6mm}
Файл \code{'1.txt'} містить відомість з рядками, в яких через двокрапку записано поряд\-ковий номер (ціле), назву предмету (знаків пунктуації не містить), прізвище студента, ім'я студента, номер заліковки, оцінку в балах за 100 бальною шкалою.

Білі символи навколо роздільників не допускаються.

Приклад рядка файлу



\begin{verbatim}
13:algebra:surname:name:123456:77
\end{verbatim}


Нижче наведено код, який має в кожному рядку файлу переставити місцями назву предмету та номер заліковки.

Рядок з прикладу має бути замінений на такий



\begin{verbatim}
13:123456:surname:name:algebra:77
\end{verbatim}




Код:



\begin{verbatim}
regex = re.compile(r'XXX')
with open('1.txt', encoding='utf8') as f:
    lst = f.readlines()
for i in range(len(lst)):
    lst[i] = re.sub(regex, YYY, lst[i])
with open('1.txt', 'w', encoding='utf8') as f:
    f.writelines(lst)
\end{verbatim}





Який рядок треба записати на місце \kw{XXX} ?
Який рядок треба записати на місце \kw{YYY} ?\par

%\end {large}
}
\newpage

{{23.05.2025 Контрольна робота, варіант  \No \textbf { 2 }} група К1{\parbox {0.5 cm} \dotfill}  ПІБ {\parbox {2.5 cm} \dotfill}\par}
{
%\begin {large}
\begin{minipage}[t]{9.7cm}
               
{\begin{center}Завдання {1}\end{center}}
\vspace{-6mm}
Укажіть ТИП значення виразу.\par\par
\begin{verbatim}
[x for x in range(7)]
\end{verbatim}

\vspace{15mm}
У завданнях 2, 3 вказати, що буде виведено за виконання. Повідомлення інтерпретатора про помилку позначати ERROR

               
{\begin{center}Завдання {2}\end{center}}
\vspace{-6mm}

\begin{verbatim}
try:
    s = {87:4, 47:6, 88:4, 87:5}
    s[33] = 2
    for x in s.keys():
        print(x, sep=' ', end=' ')
except:
    print('ERR')
\end{verbatim}


\end{minipage}
\vline \hspace{8mm}
\begin{minipage}[t]{8.0cm}                
               
{\begin{center}Завдання {3}\end{center}}
\vspace{-6mm}

\begin{verbatim}
b = 0
m = 1

class B:
    b = 2
    
    def __init__(self):
        self.b = 3
        b = 4
        B.b = 5

obj = B()
try:
    print(b, end=' ')
    print(m, end=' ')
    print(B.b, end=' ')
    print(obj.b)
except:
    print('ERR')
\end{verbatim}

\end{minipage}

               
{\begin{center}Завдання {4}\end{center}}
\vspace{-6mm}
Нехай змінна \kw{s} позначає послідовність цілих, по якій можна ітерувати.

Написати вираз, що перевіряє, чи правда, що всі елементи послідовності не діляться на 8.\par

\vspace{10mm}
               
{\begin{center}Завдання {5}\end{center}}
\vspace{-6mm}
Змінні \kw{next} та \kw{prev} вузла списку позначають наступний та попередній за ним вузол відповідно. Починаючи з голови, у список записано послідовні цілі числа від~$-40$ до~$19$. Нехай змінна \kw{p} позначає вузол, що зберігає значення~$-9$.
Яке число зберігається у вузлі \kw{p.next.prev.next.next.next} ?\par

               
{\begin{center}Завдання {6}\end{center}}
\vspace{-6mm}
Значенням змінної \kw{a} є цілочисловий масив типу $numpy.ndarray$ розмірів $5{\times}6$. На скільки збільшиться сума елементів масиву \kw{a} після виконання
\begin{verbatim}
a.T[-1] += 1
\end{verbatim}


Якщо код некоректний, то в якості відповіді зазначити ERROR.\par

               
{\begin{center}Завдання {10}\end{center}}
\vspace{-6mm}

Написати функцію $f$, яка для файлу, заданого шляхом, розв’язує задачу:\\ Текстовий файл містить послідовність чисел $a_1$, $a_2$, $\ldots$, $a_t$, записану у форматі одне число на рядок. Знайти третє найбільше значення, що записане у файлі. Кожне значення, що записане у файлі, враховується стільки разів, скільки входить у файл. Для послідовності 2, 4, 12, 5, 2, 0 таким значенням буде 2.

В усіх випадках розміри файлів такі, що їх вміст повністю завантажений у пам'ять бути не може. Якість коду також оцінюється.



\vspace{10mm}\par

\newpage
\begin{minipage}[t]{9.7cm}
               
{\begin{center}Завдання {7}\end{center}}
\vspace{-6mm}

\begin{center}
	\adjustimage{max size={0.8\linewidth}{0.8\paperheight}}
	{./15BinTree/trees_exp/74.png}
\end{center}
Указати послідовність міток вершин дерева, зображеного на рис., що утвориться в результаті обходу дерева в прямому порядку. Мітки записувати через пробіл.\par Алгоритм починає роботу з кореня (на рис. це верхня вершина).\par

\end{minipage}
\vline \hspace{8mm}
\begin{minipage}[t]{8.0cm}
               
{\begin{center}Завдання {8}\end{center}}
\vspace{-6mm}
Для графа на рисунку з вершини 0 запущено алгоритм пошуку в глибину, що будує кістякове дерево. Списки суміжності вершин переглядаються алгоритмом за зростанням міток вершин.\par
\begin{center}
	\adjustimage{max size={0.9\linewidth}{0.9\paperheight}}
	{./16Graphs/Traversals/284.png}
\end{center}
\par Які з наведених ребер входять до складу отриманого кістякового дерева?\par
1) (0, 6)\par
2) (2, 4)\par
3) (3, 4)\par
4) (4, 6)\par
5) (1, 4)\par
6) (0, 5)\par
{\vskip 15pt} \par

\end{minipage}

               
{\begin{center}Завдання {9}\end{center}}
\vspace{-6mm}
Рядки журналу через кому містять тип події (ERROR, WARNING, INFO, STATUS), ім'я користувача (складається з латиниці), час події (фор\-мат як у прикладі), код події (ціле).

Білі символи навколо роздільників не допускаються.

Приклад такого рядка присвоєно змінній \kw{s}.



\begin{verbatim}
s = 'ERROR,username,2022-05-05 13:26,404'
\end{verbatim}


Нижче наведено код, за виконання якого значенням змінної \kw{out} має стати ім'я користувача.

Для наведеного прикладу значенням змінної \kw{out} має стати рядок \code{username}



\begin{verbatim}
res = re.fullmatch(r'XXX', s)
s = ''
out = YYY
\end{verbatim}


Код має працювати не тільки для конкретного рядка з прикладу, але й для кожного рядка з журналу.


Який рядок треба записати на місце \kw{XXX} ?
Який рядок треба записати на місце \kw{YYY} ?\par

%\end {large}
}
\newpage

{{23.05.2025 Контрольна робота, варіант  \No \textbf { 3 }} група К1{\parbox {0.5 cm} \dotfill}  ПІБ {\parbox {2.5 cm} \dotfill}\par}
{
%\begin {large}
\begin{minipage}[t]{9.7cm}
               
{\begin{center}Завдання {1}\end{center}}
\vspace{-6mm}
Укажіть ТИП значення виразу.\par\par
\begin{verbatim}
{(1,4), 8}
\end{verbatim}

\vspace{15mm}
У завданнях 2, 3 вказати, що буде виведено за виконання. Повідомлення інтерпретатора про помилку позначати ERROR

               
{\begin{center}Завдання {2}\end{center}}
\vspace{-6mm}

\begin{verbatim}
try:
    d = {38:2, 40:9, 48:1, 48:7}
    d[48] = 4
    for x, y in d.items():
        print(x, y, sep=' ', end=' ')
except:
    print('ERR')
\end{verbatim}


\end{minipage}
\vline \hspace{8mm}
\begin{minipage}[t]{8.0cm}                
               
{\begin{center}Завдання {3}\end{center}}
\vspace{-6mm}

\begin{verbatim}
__b = 0

class B:
    __b = 1
    
    def __init__(self):
        self.__b = 2
        __b = 3
        B.__b = 4

obj = B()
B.__b = 5
try:
    print(__b, end=' ')
    print(obj.__b, end=' ')
    print(B.__b, end=' ')
    print(obj.__b, end=' ')
    print(obj._B__b)
except:
    print('ERR')
\end{verbatim}

\end{minipage}

               
{\begin{center}Завдання {4}\end{center}}
\vspace{-6mm}
Нехай змінна \kw{s} позначає дуже довгу послідовність, по якій можна ітерувати.

Написати вираз-генератор, що задає підпослідовність її елементів, що не є рядками типу \kw{str}.

Обмеження: усі елементи шуканої підпослідовності одночасно в пам'яті розміщені бути не можуть.\par

\vspace{10mm}
               
{\begin{center}Завдання {5}\end{center}}
\vspace{-6mm}
Змінні \kw{next} та \kw{prev} вузла списку позначають наступний та попередній за ним вузол відповідно. Починаючи з голови, у список записано послідовні цілі числа від~$-65$ до~$21$. Нехай змінна \kw{p} позначає вузол, що зберігає значення~$-8$.
Яке число зберігається у вузлі \kw{p.next.prev.next.prev.prev} ?\par

               
{\begin{center}Завдання {6}\end{center}}
\vspace{-6mm}
Значенням змінної \kw{a} є цілочисловий масив типу $numpy.ndarray$ розмірів $7{\times}3$. На скільки збільшиться сума елементів масиву \kw{a} після виконання
\begin{verbatim}
a.T[1, -1] += 2
\end{verbatim}


Якщо код некоректний, то в якості відповіді зазначити ERROR.\par

               
{\begin{center}Завдання {10}\end{center}}
\vspace{-6mm}

Написати функцію $f$, яка для файлу, заданого шляхом, розв’язує задачу:\\ Текстовий файл містить послідовність чисел $a_1$, $a_2$, $\ldots$, $a_t$, записану у форматі одне число на рядок. Знайти третє найменше значення, що записане у файлі. Кожне значення, що записане у файлі, враховується стільки разів, скільки входить у файл. Для послідовності 2, 4, 12, 5, 2, 0 таким значенням буде 2.

В усіх випадках розміри файлів такі, що їх вміст повністю завантажений у пам'ять бути не може. Якість коду також оцінюється.



\vspace{10mm}\par

\newpage
\begin{minipage}[t]{9.7cm}
               
{\begin{center}Завдання {7}\end{center}}
\vspace{-6mm}

\begin{center}
	\adjustimage{max size={0.8\linewidth}{0.8\paperheight}}
	{./15BinTree/trees_exp/134.png}
\end{center}
Указати послідовність міток вершин дерева, зображеного на рис., що утвориться в результаті обходу дерева в оберненому порядку. Мітки записувати через пробіл.\par Алгоритм починає роботу з кореня (на рис. це верхня вершина).\par

\end{minipage}
\vline \hspace{8mm}
\begin{minipage}[t]{8.0cm}
               
{\begin{center}Завдання {8}\end{center}}
\vspace{-6mm}
Для графа на рисунку з вершини 0 запущено алгоритм пошуку в ширину, що будує кістякове дерево. Списки суміжності вершин переглядаються алгоритмом за зростанням міток вершин.\par
\begin{center}
	\adjustimage{max size={0.9\linewidth}{0.9\paperheight}}
	{./16Graphs/Traversals/329.png}
\end{center}
\par Які з наведених ребер входять до складу отриманого кістякового дерева?\par
1) (1, 5)\par
2) (1, 3)\par
3) (3, 5)\par
4) (2, 6)\par
5) (4, 5)\par
6) (0, 3)\par
{\vskip 15pt} \par

\end{minipage}

               
{\begin{center}Завдання {9}\end{center}}
\vspace{-6mm}
Файл \code{'1.txt'} містить журнал з рядками, в яких через двокрапку записано поряд\-ко\-вий номер запису в журналі (ціле), тип події (ERROR, WARNING, INFO, STATUS), ім'я користувача (складається з латиниці), час події (фор\-мат як у прикладі), опис події.

Білі символи навколо роздільників не допускаються.

Приклад рядка файлу



\begin{verbatim}
13:ERROR:username:2022-05-05 13-26:something happens
\end{verbatim}


Нижче наведено код, який має залишити у файлі в кожному рядку тільки тип події та її опис.

Рядок з прикладу має бути замінений на такий



\begin{verbatim}
ERROR:something happens
\end{verbatim}




Код:



\begin{verbatim}
regex = re.compile(r'XXX')
with open('1.txt', encoding='utf8') as f:
    lst = f.readlines()
for i in range(len(lst)):
    lst[i] = re.sub(regex, YYY, lst[i])
with open('1.txt', 'w', encoding='utf8') as f:
    f.writelines(lst)
\end{verbatim}







Який рядок треба записати на місце \kw{XXX} ?
Який рядок треба записати на місце \kw{YYY} ?\par

%\end {large}
}
\newpage

{{23.05.2025 Контрольна робота, варіант  \No \textbf { 4 }} група К1{\parbox {0.5 cm} \dotfill}  ПІБ {\parbox {2.5 cm} \dotfill}\par}
{
%\begin {large}
\begin{minipage}[t]{9.7cm}
               
{\begin{center}Завдання {1}\end{center}}
\vspace{-6mm}
Укажіть ТИП значення виразу.\par\par
\begin{verbatim}
{x:0 for x in range(7)}
\end{verbatim}

\vspace{15mm}
У завданнях 2, 3 вказати, що буде виведено за виконання. Повідомлення інтерпретатора про помилку позначати ERROR

               
{\begin{center}Завдання {2}\end{center}}
\vspace{-6mm}

\begin{verbatim}
try:
    s = {52:3, 48:2, 73:4, 41:6, 43:9}
    s[73] = 3
    for x in s.items():
        print(x, sep=' ', end=' ')
except:
    print('ERR')
\end{verbatim}


\end{minipage}
\vline \hspace{8mm}
\begin{minipage}[t]{8.0cm}                
               
{\begin{center}Завдання {3}\end{center}}
\vspace{-6mm}

\begin{verbatim}
a = 0
h = 1

class C:
    h = 2
    
    def __init__(self):
        global h
        self.h = 3
        h = 4
        C.h = 5

obj = C()
try:
    print(a, end=' ')
    print(h, end=' ')
    print(C.h, end=' ')
    print(obj.h)
except:
    print('ERR')
\end{verbatim}

\end{minipage}

               
{\begin{center}Завдання {4}\end{center}}
\vspace{-6mm}

Записати ВИРАЗ, значенням якого є список, що складається з кубів цілих чисел від 20 до 2015 (включно), що не діляться на 5, записаних в обернено\-му порядку.\par

\vspace{10mm}
               
{\begin{center}Завдання {5}\end{center}}
\vspace{-6mm}
Змінні \kw{next} та \kw{prev} вузла списку позначають наступний та попередній за ним вузол відповідно. Починаючи з голови, у список записано послідовні цілі числа від~$-39$ до~$42$. Нехай змінна \kw{p} позначає вузол, що зберігає значення~$5$.
Яке число зберігається у вузлі \kw{p.prev.next.prev.next.next} ?\par

               
{\begin{center}Завдання {6}\end{center}}
\vspace{-6mm}
Значенням змінної \kw{a} є цілочисловий масив типу $numpy.ndarray$ розмірів $3{\times}5$. На скільки збільшиться сума елементів масиву \kw{a} після виконання
\begin{verbatim}
a[-2] += 2
\end{verbatim}


Якщо код некоректний, то в якості відповіді зазначити ERROR.\par

               
{\begin{center}Завдання {10}\end{center}}
\vspace{-6mm}

Написати функцію $f$, яка для файлу, заданого шляхом, розв’язує задачу:\\ Текстовий файл містить послідовність чисел $a_1$, $a_2$, $\ldots$, $a_t$, записану у форматі одне число на рядок. Знайти третє найменше значення, що записане у файлі. Кожне значення, що записане у файлі, враховується тільки один раз. \par
Для послідовності 1, 4, 12, 5, 1, 0 таким значенням буде 4.

В усіх випадках розміри файлів такі, що їх вміст повністю завантажений у пам'ять бути не може. Якість коду також оцінюється.



\vspace{10mm}\par

\newpage
\begin{minipage}[t]{9.7cm}
               
{\begin{center}Завдання {7}\end{center}}
\vspace{-6mm}

\begin{center}
	\adjustimage{max size={0.8\linewidth}{0.8\paperheight}}
	{./15BinTree/trees_exp/170.png}
\end{center}
Указати послідовність міток вершин дерева, зображеного на рис., що утвориться в результаті обходу дерева в прямому порядку. Мітки записувати через пробіл.\par Алгоритм починає роботу з кореня (на рис. це верхня вершина).\par

\end{minipage}
\vline \hspace{8mm}
\begin{minipage}[t]{8.0cm}
               
{\begin{center}Завдання {8}\end{center}}
\vspace{-6mm}
Для графа на рисунку з вершини 4 запущено алгоритм пошуку в глибину, що будує кістякове дерево. Списки суміжності вершин переглядаються алгоритмом за зростанням міток вершин.\par
\begin{center}
	\adjustimage{max size={0.9\linewidth}{0.9\paperheight}}
	{./16Graphs/Traversals/12.png}
\end{center}
\par Які з наведених ребер входять до складу отриманого кістякового дерева?\par
1) (1, 5)\par
2) (4, 6)\par
3) (2, 6)\par
4) (5, 6)\par
5) (2, 4)\par
6) (3, 5)\par
{\vskip 15pt} \par

\end{minipage}

               
{\begin{center}Завдання {9}\end{center}}
\vspace{-6mm}
Файл \code{'1.txt'} містить журнал з рядками, в яких через кому записано поряд\-ко\-вий номер запису (ціле), ім'я користувача (складається з латиниці), час події (фор\-мат як у прикладі), тип події (ERROR, WARNING, INFO, STATUS), код події (ціле).

Білі символи навколо роздільників не допускаються.

Приклад рядка журналу



\begin{verbatim}
12,username,2022-05-05 13:26,ERROR,404
\end{verbatim}


Нижче наведено код, який шукає усі події, що відбулись\\ для користувача \code{anonymous}, та зберігає їх типи в змінній \kw{codes}.



\begin{verbatim}
codes = set()
regex = re.compile(r'XXX')
with open('1.txt', encoding='utf8') as f:
    for s in f:
        res = re.fullmatch(regex, s)
        if not res: continue
        s = ''
        codes.add(YYY)
\end{verbatim}



Який рядок треба записати на місце \kw{XXX} ?
Який рядок треба записати на місце \kw{YYY} ?\par

%\end {large}
}
\newpage

{{23.05.2025 Контрольна робота, варіант  \No \textbf { 5 }} група К1{\parbox {0.5 cm} \dotfill}  ПІБ {\parbox {2.5 cm} \dotfill}\par}
{
%\begin {large}
\begin{minipage}[t]{9.7cm}
               
{\begin{center}Завдання {1}\end{center}}
\vspace{-6mm}
Укажіть ТИП значення виразу.\par\par
\begin{verbatim}
{x for x in range(7)}
\end{verbatim}

\vspace{15mm}
У завданнях 2, 3 вказати, що буде виведено за виконання. Повідомлення інтерпретатора про помилку позначати ERROR

               
{\begin{center}Завдання {2}\end{center}}
\vspace{-6mm}

\begin{verbatim}
try:
    d = {23:2, 74:8, 44:6, 29:7, 29:1}
    d[23] = 3
    for x in d.items():
        print(*x, sep=' ', end=' ')
except:
    print('ERR')
\end{verbatim}


\end{minipage}
\vline \hspace{8mm}
\begin{minipage}[t]{8.0cm}                
               
{\begin{center}Завдання {3}\end{center}}
\vspace{-6mm}

\begin{verbatim}
d = 0
j = 1

class A:
    j = 2
    
    def __init__(self):
        self.d = 3
        j = 4
        A.j = 5

obj = A()
try:
    print(d, end=' ')
    print(j, end=' ')
    print(A.j, end=' ')
    print(obj.d)
except:
    print('ERR')
\end{verbatim}

\end{minipage}

               
{\begin{center}Завдання {4}\end{center}}
\vspace{-6mm}
Записати ВИРАЗ, значенням якого є словник, ключами якого є цілі числа від 12 до 2007 (включно), що діляться на 7, а ключам відповідають їх куби\par

\vspace{10mm}
               
{\begin{center}Завдання {5}\end{center}}
\vspace{-6mm}
Змінні \kw{next} та \kw{prev} вузла списку позначають наступний та попередній за ним вузол відповідно. Починаючи з голови, у список записано послідовні цілі числа від~$-19$ до~$40$. Нехай змінна \kw{p} позначає вузол, що зберігає значення~$9$.
Яке число зберігається у вузлі \kw{p.prev.next.next.next.next} ?\par

               
{\begin{center}Завдання {6}\end{center}}
\vspace{-6mm}
Значенням змінної \kw{a} є цілочисловий масив типу $numpy.ndarray$ розмірів $6{\times}4$. На скільки збільшиться сума елементів масиву \kw{a} після виконання
\begin{verbatim}
a[-3] += 1
\end{verbatim}


Якщо код некоректний, то в якості відповіді зазначити ERROR.\par

               
{\begin{center}Завдання {10}\end{center}}
\vspace{-6mm}

Написати функцію $f$, яка для файлу, заданого шляхом, розв’язує задачу:\\ Текстовий файл містить послідовність чисел $a_1$, $a_2$, $\ldots$, $a_t$, записану у форматі одне число на рядок. Знайти третє найбільше значення, що записане у файлі. Кожне значення, що записане у файлі, враховується тільки один раз. Для послідовності 1, 4, 12, 5, 1, 0 таким значенням буде 4.

В усіх випадках розміри файлів такі, що їх вміст повністю завантажений у пам'ять бути не може. Якість коду також оцінюється.



\vspace{10mm}\par

\newpage
\begin{minipage}[t]{9.7cm}
               
{\begin{center}Завдання {7}\end{center}}
\vspace{-6mm}

\begin{center}
	\adjustimage{max size={0.8\linewidth}{0.8\paperheight}}
	{./15BinTree/trees_exp/135.png}
\end{center}
Указати послідовність міток вершин дерева, зображеного на рис., що утвориться в результаті обходу дерева в прямому порядку. Мітки записувати через пробіл.\par Алгоритм починає роботу з кореня (на рис. це верхня вершина).\par

\end{minipage}
\vline \hspace{8mm}
\begin{minipage}[t]{8.0cm}
               
{\begin{center}Завдання {8}\end{center}}
\vspace{-6mm}
Для графа на рисунку з вершини 1 запущено алгоритм пошуку в ширину, що будує кістякове дерево. Списки суміжності вершин переглядаються алгоритмом за зростанням міток вершин.\par
\begin{center}
	\adjustimage{max size={0.9\linewidth}{0.9\paperheight}}
	{./16Graphs/Traversals/64.png}
\end{center}
\par Які з наведених ребер входять до складу отриманого кістякового дерева?\par
1) (1, 2)\par
2) (5, 6)\par
3) (0, 4)\par
4) (1, 3)\par
5) (1, 6)\par
6) (0, 1)\par
{\vskip 15pt} \par

\end{minipage}

               
{\begin{center}Завдання {9}\end{center}}
\vspace{-6mm}
Файл \code{'1.txt'} містить відомість з рядками, в яких через кому записа\-но поряд\-ковий номер (ціле), назву предмету (знаків пунктуації не містить), прізвище студента, ім'я студента, номер заліковки, оцінку в балах за 100 бальною шкалою.

Білі символи навколо роздільників не допускаються.

Приклад рядка файлу



\begin{verbatim}
13,algebra,surname,name,123456,77
\end{verbatim}


Нижче наведено код, який має залишити у файлі в кожному рядку тільки номер заліковки, назву предмету та оцінку в балах за 100-бальною шкалою (у заданій послі\-дов\-ності).

Рядок з прикладу має бути замінений на такий



\begin{verbatim}
123456,algebra,77
\end{verbatim}




Код



\begin{verbatim}
regex = re.compile(r'XXX')
with open('1.txt', encoding='utf8') as f:
    lst = f.readlines()
for i in range(len(lst)):
    lst[i] = re.sub(regex, YYY, lst[i])
with open('1.txt', 'w', encoding='utf8') as f:
    f.writelines(lst)
\end{verbatim}





Який рядок треба записати на місце \kw{XXX} ?
Який рядок треба записати на місце \kw{YYY} ?\par

%\end {large}
}
\newpage

{{23.05.2025 Контрольна робота, варіант  \No \textbf { 6 }} група К1{\parbox {0.5 cm} \dotfill}  ПІБ {\parbox {2.5 cm} \dotfill}\par}
{
%\begin {large}
\begin{minipage}[t]{9.7cm}
               
{\begin{center}Завдання {1}\end{center}}
\vspace{-6mm}
Укажіть ТИП значення виразу.\par\par
\begin{verbatim}
[{1,2}, 5]
\end{verbatim}

\vspace{15mm}
У завданнях 2, 3 вказати, що буде виведено за виконання. Повідомлення інтерпретатора про помилку позначати ERROR

               
{\begin{center}Завдання {2}\end{center}}
\vspace{-6mm}

\begin{verbatim}
try:
    t = {63:4, 50:2, 62:9, 62:2}
    t[62] = 9
    for x in t :
        print(x, sep=' ', end=' ')
except:
    print('ERR')
\end{verbatim}


\end{minipage}
\vline \hspace{8mm}
\begin{minipage}[t]{8.0cm}                
               
{\begin{center}Завдання {3}\end{center}}
\vspace{-6mm}

\begin{verbatim}
_d = 0

class A:
    _d = 1
    
    def __init__(self):
        self._d = 2
        _d = 3
        A._d = 4

obj = A()
A._d = 5
try:
    print(_d, end=' ')
    print(obj._d, end=' ')
    print(A._d, end=' ')
    print(obj._d, end=' ')
    print(obj._A__d)
except:
    print('ERR')
\end{verbatim}

\end{minipage}

               
{\begin{center}Завдання {4}\end{center}}
\vspace{-6mm}
Задано множину \kw{x}, що складається з цілих чисел. Записати ВИРАЗ, значенням якого є множина, що складається  з елементів множини \kw{x}, що не діляться на 7.\par

\vspace{10mm}
               
{\begin{center}Завдання {5}\end{center}}
\vspace{-6mm}
Змінні \kw{next} та \kw{prev} вузла списку позначають наступний та попередній за ним вузол відповідно. Починаючи з голови, у список записано послідовні цілі числа від~$-15$ до~$38$. Нехай змінна \kw{p} позначає вузол, що зберігає значення~$-1$.
Яке число зберігається у вузлі \kw{p.next.prev.prev.prev.prev} ?\par

               
{\begin{center}Завдання {6}\end{center}}
\vspace{-6mm}
Значенням змінної \kw{a} є цілочисловий масив типу $numpy.ndarray$ розмірів $5{\times}4$. На скільки збільшиться сума елементів масиву \kw{a} після виконання
\begin{verbatim}
a.T[2] += 3
\end{verbatim}


Якщо код некоректний, то в якості відповіді зазначити ERROR.\par

               
{\begin{center}Завдання {10}\end{center}}
\vspace{-6mm}

Написати функцію $f$, яка для файлу, заданого шляхом, розв’язує задачу:\\ Кожен з двох текстових файлів містить послідовність чисел, записану у форматі одне число на рядок.  Перевірити, чи задають вони одну ту саму послідовність.

В усіх випадках розміри файлів такі, що їх вміст повністю завантажений у пам'ять бути не може. Якість коду також оцінюється.



\vspace{10mm}\par

\newpage
\begin{minipage}[t]{9.7cm}
               
{\begin{center}Завдання {7}\end{center}}
\vspace{-6mm}

\begin{center}
	\adjustimage{max size={0.8\linewidth}{0.8\paperheight}}
	{./15BinTree/trees_exp/161.png}
\end{center}
Указати послідовність міток вершин дерева, зображеного на рис., що утвориться в результаті обходу дерева в оберненому порядку. Мітки записувати через пробіл.\par Алгоритм починає роботу з кореня (на рис. це верхня вершина).\par

\end{minipage}
\vline \hspace{8mm}
\begin{minipage}[t]{8.0cm}
               
{\begin{center}Завдання {8}\end{center}}
\vspace{-6mm}
Для графа на рисунку з вершини 1 запущено алгоритм пошуку в глибину, що будує кістякове дерево. Списки суміжності вершин переглядаються алгоритмом за зростанням міток вершин.\par
\begin{center}
	\adjustimage{max size={0.9\linewidth}{0.9\paperheight}}
	{./16Graphs/Traversals/256.png}
\end{center}
\par Які з наведених ребер входять до складу отриманого кістякового дерева?\par
1) (0, 4)\par
2) (1, 4)\par
3) (3, 5)\par
4) (0, 3)\par
5) (3, 4)\par
6) (1, 3)\par
{\vskip 15pt} \par

\end{minipage}

               
{\begin{center}Завдання {9}\end{center}}
\vspace{-6mm}
Файл \code{'1.txt'} містить відомість з рядками, в яких через кому записано поряд\-ковий номер (ціле), назву предмету (знаків пунктуації не містить), прізвище студента, ім'я студента, номер заліковки, оцінку в балах за 100 бальною шкалою.

Білі символи навколо роздільників не допускаються.

Приклад рядка файлу



\begin{verbatim}
13,algebra,surname,name,123456,77
\end{verbatim}


Нижче наведено код, який має в кожному рядку файлу переставити місцями назву предмету та номер заліковки.

Рядок з прикладу має бути замінений на такий



\begin{verbatim}
13,123456,surname,name,algebra,77
\end{verbatim}




Код



\begin{verbatim}
regex = re.compile(r'XXX')
with open('1.txt', encoding='utf8') as f:
    lst = f.readlines()
for i in range(len(lst)):
    lst[i] = re.sub(regex, YYY, lst[i])
with open('1.txt', 'w', encoding='utf8') as f:
    f.writelines(lst)
\end{verbatim}





Який рядок треба записати на місце \kw{XXX} ?
Який рядок треба записати на місце \kw{YYY} ?\par

%\end {large}
}
\newpage

{{23.05.2025 Контрольна робота, варіант  \No \textbf { 7 }} група К1{\parbox {0.5 cm} \dotfill}  ПІБ {\parbox {2.5 cm} \dotfill}\par}
{
%\begin {large}
\begin{minipage}[t]{9.7cm}
               
{\begin{center}Завдання {1}\end{center}}
\vspace{-6mm}
Укажіть ТИП значення виразу.\par\par
\begin{verbatim}
{y:y*y for y in range(-7,7)}
\end{verbatim}

\vspace{15mm}
У завданнях 2, 3 вказати, що буде виведено за виконання. Повідомлення інтерпретатора про помилку позначати ERROR

               
{\begin{center}Завдання {2}\end{center}}
\vspace{-6mm}

\begin{verbatim}
try:
    d = {25:6, 44:3, 55:7, 55:8}
    d[55] = 3
    for x in d.keys():
        print(x, sep=' ', end=' ')
except:
    print('ERR')
\end{verbatim}


\end{minipage}
\vline \hspace{8mm}
\begin{minipage}[t]{8.0cm}                
               
{\begin{center}Завдання {3}\end{center}}
\vspace{-6mm}

\begin{verbatim}
_a = 0

class D:
    _a = 1
    
    def __init__(self):
        self._a = 2
        _a = 3
        D._a = 4

obj = D()
obj._a = 5
try:
    print(_a, end=' ')
    print(obj._a, end=' ')
    print(D._a, end=' ')
    print(obj._a, end=' ')
    print(obj._D__a)
except:
    print('ERR')
\end{verbatim}

\end{minipage}

               
{\begin{center}Завдання {4}\end{center}}
\vspace{-6mm}

 Задано числову послідовність \kw{s} довжини 6000. Записати ВИРАЗ, значенням якого є список, що складається з останніх 2019 елементів послідовності.\par

\vspace{10mm}
               
{\begin{center}Завдання {5}\end{center}}
\vspace{-6mm}
Змінні \kw{next} та \kw{prev} вузла списку позначають наступний та попередній за ним вузол відповідно. Починаючи з голови, у список записано послідовні цілі числа від~$-16$ до~$27$. Нехай змінна \kw{p} позначає вузол, що зберігає значення~$6$.
Яке число зберігається у вузлі \kw{p.next.next.prev.prev.next} ?\par

               
{\begin{center}Завдання {6}\end{center}}
\vspace{-6mm}
Значенням змінної \kw{a} є цілочисловий масив типу $numpy.ndarray$ розмірів $6{\times}5$. На скільки збільшиться сума елементів масиву \kw{a} після виконання
\begin{verbatim}
a[-3] += 2
\end{verbatim}


Якщо код некоректний, то в якості відповіді зазначити ERROR.\par

               
{\begin{center}Завдання {10}\end{center}}
\vspace{-6mm}

Написати функцію $f$, яка для файлу, заданого шляхом, розв’язує задачу:\\ Текстовий файл містить послідовність чисел $a_1$, $a_2$, $\ldots$, $a_t$, записану у форматі одне число на рядок. Знайти третє найбільше значення, що записане у файлі. Кожне значення, що записане у файлі, враховується стільки разів, скільки входить у файл. Для послідовності 2, 4, 12, 5, 2, 0 таким значенням буде 2.

В усіх випадках розміри файлів такі, що їх вміст повністю завантажений у пам'ять бути не може. Якість коду також оцінюється.



\vspace{10mm}\par

\newpage
\begin{minipage}[t]{9.7cm}
               
{\begin{center}Завдання {7}\end{center}}
\vspace{-6mm}

\begin{center}
	\adjustimage{max size={0.8\linewidth}{0.8\paperheight}}
	{./15BinTree/trees_exp/81.png}
\end{center}
Указати послідовність міток вершин дерева, зображеного на рис., що утвориться в результаті обходу дерева в прямому порядку. Мітки записувати через пробіл.\par Алгоритм починає роботу з кореня (на рис. це верхня вершина).\par

\end{minipage}
\vline \hspace{8mm}
\begin{minipage}[t]{8.0cm}
               
{\begin{center}Завдання {8}\end{center}}
\vspace{-6mm}
Для графа на рисунку з вершини 6 запущено алгоритм пошуку в ширину, що будує кістякове дерево. Списки суміжності вершин переглядаються алгоритмом за зростанням міток вершин.\par
\begin{center}
	\adjustimage{max size={0.9\linewidth}{0.9\paperheight}}
	{./16Graphs/Traversals/286.png}
\end{center}
\par Які з наведених ребер входять до складу отриманого кістякового дерева?\par
1) (0, 1)\par
2) (0, 5)\par
3) (2, 5)\par
4) (4, 6)\par
5) (1, 6)\par
6) (1, 3)\par
{\vskip 15pt} \par

\end{minipage}

               
{\begin{center}Завдання {9}\end{center}}
\vspace{-6mm}
Файл \code{'1.txt'} містить журнал з рядками, в яких через кому записано поряд\-ко\-вий номер запису в журналі (ціле), код події (ціле), тип події (ERROR,\\ WARNING, INFO, STATUS), ім'я користувача (складається з латиниці), час події (фор\-мат як у прикладі), опис події.

Білі символи навколо роздільників не допускаються.

Приклад рядка файлу



\begin{verbatim}
13,404,ERROR,username,2022-05-05 13:26,something happens
\end{verbatim}


Нижче наведено код, який має в кожному рядку файли переставити місцями код та тип події.

Рядок з прикладу має бути замінений на такий



\begin{verbatim}
13,ERROR,404,username,2022-05-05 13:26,something happens
\end{verbatim}




Код:



\begin{verbatim}
regex = re.compile(r'XXX')
with open('1.txt', encoding='utf8') as f:
    lst = f.readlines()
for i in range(len(lst)):
    lst[i] = re.sub(regex, YYY, lst[i])
with open('1.txt', 'w', encoding='utf8') as f:
    f.writelines(lst)
\end{verbatim}





Який рядок треба записати на місце \kw{XXX} ?
Який рядок треба записати на місце \kw{YYY} ?\par

%\end {large}
}
\newpage

{{23.05.2025 Контрольна робота, варіант  \No \textbf { 8 }} група К1{\parbox {0.5 cm} \dotfill}  ПІБ {\parbox {2.5 cm} \dotfill}\par}
{
%\begin {large}
\begin{minipage}[t]{9.7cm}
               
{\begin{center}Завдання {1}\end{center}}
\vspace{-6mm}
Укажіть ТИП значення виразу.\par\par
\begin{verbatim}
{y:13 for y in range(-7,7)}
\end{verbatim}

\vspace{15mm}
У завданнях 2, 3 вказати, що буде виведено за виконання. Повідомлення інтерпретатора про помилку позначати ERROR

               
{\begin{center}Завдання {2}\end{center}}
\vspace{-6mm}

\begin{verbatim}
try:
    t = {23:2, 24:7, 48:8, 23:9}
    t[24] = 6
    for x, y in t.items():
        print(x, y, sep=' ', end=' ')
except:
    print('ERR')
\end{verbatim}


\end{minipage}
\vline \hspace{8mm}
\begin{minipage}[t]{8.0cm}                
               
{\begin{center}Завдання {3}\end{center}}
\vspace{-6mm}

\begin{verbatim}
a = 0

class D:
    a = 1
    
    def __init__(self):
        self.a = 2
        a = 3
        D.a = 4

obj = D()
obj.a = 5
try:
    print(a, end=' ')
    print(obj.a, end=' ')
    print(D.a, end=' ')
    print(obj.a, end=' ')
    print(obj._D__a)
except:
    print('ERR')
\end{verbatim}

\end{minipage}

               
{\begin{center}Завдання {4}\end{center}}
\vspace{-6mm}
Задано словник \kw{d} з цілими ключами та значеннями. Записати ВИРАЗ, значенням якого є множина, що складається  з тих значень, чиї ключі не діляться на 11.\par

\vspace{10mm}
               
{\begin{center}Завдання {5}\end{center}}
\vspace{-6mm}
Змінні \kw{next} та \kw{prev} вузла списку позначають наступний та попередній за ним вузол відповідно. Починаючи з голови, у список записано послідовні цілі числа від~$-50$ до~$16$. Нехай змінна \kw{p} позначає вузол, що зберігає значення~$-5$.
Яке число зберігається у вузлі \kw{p.prev.prev.next.next.prev} ?\par

               
{\begin{center}Завдання {6}\end{center}}
\vspace{-6mm}
Значенням змінної \kw{a} є цілочисловий масив типу $numpy.ndarray$ розмірів $3{\times}6$. На скільки збільшиться сума елементів масиву \kw{a} після виконання
\begin{verbatim}
a.T[0] += 1
\end{verbatim}


Якщо код некоректний, то в якості відповіді зазначити ERROR.\par

               
{\begin{center}Завдання {10}\end{center}}
\vspace{-6mm}

Написати функцію $f$, яка для файлу, заданого шляхом, розв’язує задачу:\\ Текстовий файл містить послідовність чисел $a_1$, $a_2$, $\ldots$, $a_t$, записану у форматі одне число на рядок. Знайти середнє арифметичне чисел $a_1$, $a_4$, $a_7$, $\ldots$ (починаючи з першого через два).

В усіх випадках розміри файлів такі, що їх вміст повністю завантажений у пам'ять бути не може. Якість коду також оцінюється.



\vspace{10mm}\par

\newpage
\begin{minipage}[t]{9.7cm}
               
{\begin{center}Завдання {7}\end{center}}
\vspace{-6mm}

\begin{center}
	\adjustimage{max size={0.8\linewidth}{0.8\paperheight}}
	{./15BinTree/trees_exp/7.png}
\end{center}
Указати послідовність міток вершин дерева, зображеного на рис., що утвориться в результаті обходу дерева в оберненому порядку. Мітки записувати через пробіл.\par Алгоритм починає роботу з кореня (на рис. це верхня вершина).\par

\end{minipage}
\vline \hspace{8mm}
\begin{minipage}[t]{8.0cm}
               
{\begin{center}Завдання {8}\end{center}}
\vspace{-6mm}
Для графа на рисунку з вершини 5 запущено алгоритм пошуку в ширину, що будує кістякове дерево. Списки суміжності вершин переглядаються алгоритмом за зростанням міток вершин.\par
\begin{center}
	\adjustimage{max size={0.9\linewidth}{0.9\paperheight}}
	{./16Graphs/Traversals/31.png}
\end{center}
\par Які з наведених ребер входять до складу отриманого кістякового дерева?\par
1) (0, 5)\par
2) (2, 4)\par
3) (0, 6)\par
4) (0, 4)\par
5) (3, 4)\par
6) (1, 5)\par
{\vskip 15pt} \par

\end{minipage}

               
{\begin{center}Завдання {9}\end{center}}
\vspace{-6mm}
Файл \code{'1.txt'} містить журнал з рядками, в яких через кому записано поряд\-ко\-вий номер запису в журналі (ціле), код події (ціле), тип події (ERROR, \\WARNING, INFO, STATUS), ім'я користувача (складається з латиниці), час події (фор\-мат як у прикладі), опис події.

Білі символи навколо роздільників не допускаються.

Приклад рядка файлу



\begin{verbatim}
13,404,ERROR,username,2022-05-05 13:26,something happens
\end{verbatim}


Нижче наведено код, який має залишити у файлі в кожному рядку тільки код події та ім'я користувача. Рядок з прикладу має бути замінений на такий



\begin{verbatim}
404,username
\end{verbatim}




Код:



\begin{verbatim}
regex = re.compile(r'XXX')
with open('1.txt', encoding='utf8') as f:
    lst = f.readlines()
for i in range(len(lst)):
    lst[i] = re.sub(regex, YYY, lst[i])
with open('1.txt', 'w', encoding='utf8') as f:
    f.writelines(lst)
\end{verbatim}





Який рядок треба записати на місце \kw{XXX} ?
Який рядок треба записати на місце \kw{YYY} ?\par

%\end {large}
}
\newpage

{{23.05.2025 Контрольна робота, варіант  \No \textbf { 9 }} група К1{\parbox {0.5 cm} \dotfill}  ПІБ {\parbox {2.5 cm} \dotfill}\par}
{
%\begin {large}
\begin{minipage}[t]{9.7cm}
               
{\begin{center}Завдання {1}\end{center}}
\vspace{-6mm}
Укажіть ТИП значення виразу.\par\par
\begin{verbatim}
(2, 7, -4)
\end{verbatim}

\vspace{15mm}
У завданнях 2, 3 вказати, що буде виведено за виконання. Повідомлення інтерпретатора про помилку позначати ERROR

               
{\begin{center}Завдання {2}\end{center}}
\vspace{-6mm}

\begin{verbatim}
try:
    s = {50:6, 18:8, 77:1, 80:8, 50:7}
    s[31] = 1
    for x in s.items():
        print(*x, sep=' ', end=' ')
except:
    print('ERR')
\end{verbatim}


\end{minipage}
\vline \hspace{8mm}
\begin{minipage}[t]{8.0cm}                
               
{\begin{center}Завдання {3}\end{center}}
\vspace{-6mm}

\begin{verbatim}
c = 0
k = 1

class D:
    c = 2
    
    def __init__(self):
        global c
        self.k = 3
        c = 4
        D.c = 5

obj = D()
try:
    print(c, end=' ')
    print(k, end=' ')
    print(D.c, end=' ')
    print(obj.k)
except:
    print('ERR')
\end{verbatim}

\end{minipage}

               
{\begin{center}Завдання {4}\end{center}}
\vspace{-6mm}
Нехай змінна \kw{s} позначає дуже довгу послідовність, по якій можна ітерувати.

Написати вираз-генератор, що задає підпослідовність її елементів, що є числами типу \kw{int}.

Обмеження: усі елементи шуканої підпослідовності одночасно в пам'яті розміщені бути не можуть.\par

\vspace{10mm}
               
{\begin{center}Завдання {5}\end{center}}
\vspace{-6mm}
Змінні \kw{next} та \kw{prev} вузла списку позначають наступний та попередній за ним вузол відповідно. Починаючи з голови, у список записано послідовні цілі числа від~$-46$ до~$24$. Нехай змінна \kw{p} позначає вузол, що зберігає значення~$-6$.
Яке число зберігається у вузлі \kw{p.next.next.prev.next.next} ?\par

               
{\begin{center}Завдання {6}\end{center}}
\vspace{-6mm}
Значенням змінної \kw{a} є цілочисловий масив типу $numpy.ndarray$ розмірів $4{\times}7$. На скільки збільшиться сума елементів масиву \kw{a} після виконання
\begin{verbatim}
a[-2] += 3
\end{verbatim}


Якщо код некоректний, то в якості відповіді зазначити ERROR.\par

               
{\begin{center}Завдання {10}\end{center}}
\vspace{-6mm}

Написати функцію $f$, яка для файлу, заданого шляхом, розв’язує задачу:\\ Текстовий файл містить послідовність чисел $a_1$, $a_2$, $\ldots$, $a_t$, записану у форматі одне число на рядок. Знайти третє найменше значення, що записане у файлі. Кожне значення, що записане у файлі, враховується тільки один раз. \par
Для послідовності 1, 4, 12, 5, 1, 0 таким значенням буде 4.

В усіх випадках розміри файлів такі, що їх вміст повністю завантажений у пам'ять бути не може. Якість коду також оцінюється.



\vspace{10mm}\par

\newpage
\begin{minipage}[t]{9.7cm}
               
{\begin{center}Завдання {7}\end{center}}
\vspace{-6mm}

\begin{center}
	\adjustimage{max size={0.8\linewidth}{0.8\paperheight}}
	{./15BinTree/trees_exp/102.png}
\end{center}
Указати послідовність міток вершин дерева, зображеного на рис., що утвориться в результаті обходу дерева в оберненому порядку. Мітки записувати через пробіл.\par Алгоритм починає роботу з кореня (на рис. це верхня вершина).\par

\end{minipage}
\vline \hspace{8mm}
\begin{minipage}[t]{8.0cm}
               
{\begin{center}Завдання {8}\end{center}}
\vspace{-6mm}
Для графа на рисунку з вершини 0 запущено алгоритм пошуку в глибину, що будує кістякове дерево. Списки суміжності вершин переглядаються алгоритмом за зростанням міток вершин.\par
\begin{center}
	\adjustimage{max size={0.9\linewidth}{0.9\paperheight}}
	{./16Graphs/Traversals/190.png}
\end{center}
\par Які з наведених ребер входять до складу отриманого кістякового дерева?\par
1) (2, 4)\par
2) (2, 5)\par
3) (1, 6)\par
4) (0, 4)\par
5) (4, 5)\par
6) (0, 3)\par
{\vskip 15pt} \par

\end{minipage}

               
{\begin{center}Завдання {9}\end{center}}
\vspace{-6mm}
Файл \code{'1.txt'} містить журнал з рядками, в яких через двокрапку записано поряд\-ко\-вий номер запису (ціле), ім'я користувача (складається з латиниці), час події (фор\-мат як у прикладі), тип події (ERROR, WARNING, INFO, STATUS), код події (ціле).

Білі символи навколо роздільників не допускаються.

Приклад рядка журналу



\begin{verbatim}
12:username:2022-05-05 13-26:ERROR:404
\end{verbatim}


Нижче наведено код, який шукає усі події, що відбулись\\ для користувача \code{anonymous}, та зберігає їх типи в змінній \kw{codes}.



\begin{verbatim}
codes = set()
regex = re.compile(r'XXX')
with open('1.txt', encoding='utf8') as f:
    for s in f:
        res = re.fullmatch(regex, s)
        if not res: continue
        s = ''
        codes.add(YYY)
\end{verbatim}



Який рядок треба записати на місце \kw{XXX} ?
Який рядок треба записати на місце \kw{YYY} ?\par

%\end {large}
}
\newpage

{{23.05.2025 Контрольна робота, варіант  \No \textbf { 10 }} група К1{\parbox {0.5 cm} \dotfill}  ПІБ {\parbox {2.5 cm} \dotfill}\par}
{
%\begin {large}
\begin{minipage}[t]{9.7cm}
               
{\begin{center}Завдання {1}\end{center}}
\vspace{-6mm}
Укажіть ТИП значення виразу.\par\par
\begin{verbatim}
[13,2,4]+[3,5]
\end{verbatim}

\vspace{15mm}
У завданнях 2, 3 вказати, що буде виведено за виконання. Повідомлення інтерпретатора про помилку позначати ERROR

               
{\begin{center}Завдання {2}\end{center}}
\vspace{-6mm}

\begin{verbatim}
try:
    s = {77:2, 81:3, 83:1, 70:1, 70:1}
    s[77] = 4
    for x in s :
        print(x, sep=' ', end=' ')
except:
    print('ERR')
\end{verbatim}


\end{minipage}
\vline \hspace{8mm}
\begin{minipage}[t]{8.0cm}                
               
{\begin{center}Завдання {3}\end{center}}
\vspace{-6mm}

\begin{verbatim}
__b = 0

class B:
    __b = 1
    
    def __init__(self):
        self.__b = 2
        __b = 3
        B.__b = 4

obj = B()
B.__b = 5
try:
    print(__b, end=' ')
    print(obj.__b, end=' ')
    print(B.__b, end=' ')
    print(obj.__b, end=' ')
    print(obj._B__b)
except:
    print('ERR')
\end{verbatim}

\end{minipage}

               
{\begin{center}Завдання {4}\end{center}}
\vspace{-6mm}
Нехай змінна \kw{s} позначає дуже довгу послідовність цілих, по якій можна ітерувати.

Написати вираз-генератор, що задає підпослідовність її елементів, що є від'єм\-ни\-ми.

Обмеження: усі елементи шуканої підпослідовності одночасно в пам'яті розміщені бути не можуть.\par

\vspace{10mm}
               
{\begin{center}Завдання {5}\end{center}}
\vspace{-6mm}
Змінні \kw{next} та \kw{prev} вузла списку позначають наступний та попередній за ним вузол відповідно. Починаючи з голови, у список записано послідовні цілі числа від~$-31$ до~$44$. Нехай змінна \kw{p} позначає вузол, що зберігає значення~$3$.
Яке число зберігається у вузлі \kw{p.prev.prev.next.prev.prev} ?\par

               
{\begin{center}Завдання {6}\end{center}}
\vspace{-6mm}
Значенням змінної \kw{a} є цілочисловий масив типу $numpy.ndarray$ розмірів $7{\times}3$. На скільки збільшиться сума елементів масиву \kw{a} після виконання
\begin{verbatim}
a.T[1] += 2
\end{verbatim}


Якщо код некоректний, то в якості відповіді зазначити ERROR.\par

               
{\begin{center}Завдання {10}\end{center}}
\vspace{-6mm}

Написати функцію $f$, яка для файлу, заданого шляхом, розв’язує задачу:\\ Текстовий файл містить послідовність чисел $a_1$, $a_2$, $\ldots$, $a_t$, записану у форматі одне число на рядок. Знайти третє найменше значення, що записане у файлі. Кожне значення, що записане у файлі, враховується стільки разів, скільки входить у файл. Для послідовності 2, 4, 12, 5, 2, 0 таким значенням буде 2.

В усіх випадках розміри файлів такі, що їх вміст повністю завантажений у пам'ять бути не може. Якість коду також оцінюється.



\vspace{10mm}\par

\newpage
\begin{minipage}[t]{9.7cm}
               
{\begin{center}Завдання {7}\end{center}}
\vspace{-6mm}

\begin{center}
	\adjustimage{max size={0.8\linewidth}{0.8\paperheight}}
	{./15BinTree/trees_exp/45.png}
\end{center}
Указати послідовність міток вершин дерева, зображеного на рис., що утвориться в результаті обходу дерева в прямому порядку. Мітки записувати через пробіл.\par Алгоритм починає роботу з кореня (на рис. це верхня вершина).\par

\end{minipage}
\vline \hspace{8mm}
\begin{minipage}[t]{8.0cm}
               
{\begin{center}Завдання {8}\end{center}}
\vspace{-6mm}
Для графа на рисунку з вершини 1 запущено алгоритм пошуку в ширину, що будує кістякове дерево. Списки суміжності вершин переглядаються алгоритмом за зростанням міток вершин.\par
\begin{center}
	\adjustimage{max size={0.9\linewidth}{0.9\paperheight}}
	{./16Graphs/Traversals/301.png}
\end{center}
\par Які з наведених ребер входять до складу отриманого кістякового дерева?\par
1) (3, 6)\par
2) (5, 6)\par
3) (0, 4)\par
4) (2, 3)\par
5) (0, 1)\par
6) (3, 5)\par
{\vskip 15pt} \par

\end{minipage}

               
{\begin{center}Завдання {9}\end{center}}
\vspace{-6mm}
Файл \code{'1.txt'} містить журнал з рядками, в яких через двокрапку записано поряд\-ко\-вий номер запису в журналі (ціле), код події (ціле), тип події (ERROR,\\ WARNING, INFO, STATUS), ім'я користувача (складається з латиниці), час події (фор\-мат як у прикладі), опис події.

Білі символи навколо роздільників не допускаються.

Приклад рядка файлу



\begin{verbatim}
13:404:ERROR:username:2022-05-05 13:26:something happens
\end{verbatim}


Нижче наведено код, який має в кожному рядку файли переставити місцями код та тип події.

Рядок з прикладу має бути замінений на такий



\begin{verbatim}
13:ERROR:404:username:2022-05-05 13:26:something happens
\end{verbatim}




Код:



\begin{verbatim}
regex = re.compile(r'XXX')
with open('1.txt', encoding='utf8') as f:
    lst = f.readlines()
for i in range(len(lst)):
    lst[i] = re.sub(regex, YYY, lst[i])
with open('1.txt', 'w', encoding='utf8') as f:
    f.writelines(lst)
\end{verbatim}





Який рядок треба записати на місце \kw{XXX} ?
Який рядок треба записати на місце \kw{YYY} ?\par

%\end {large}
}
\newpage

{{23.05.2025 Контрольна робота, варіант  \No \textbf { 11 }} група К1{\parbox {0.5 cm} \dotfill}  ПІБ {\parbox {2.5 cm} \dotfill}\par}
{
%\begin {large}
\begin{minipage}[t]{9.7cm}
               
{\begin{center}Завдання {1}\end{center}}
\vspace{-6mm}
Укажіть ТИП значення виразу.\par\par
\begin{verbatim}
{x*x:x for x in range(-7,7)}
\end{verbatim}

\vspace{15mm}
У завданнях 2, 3 вказати, що буде виведено за виконання. Повідомлення інтерпретатора про помилку позначати ERROR

               
{\begin{center}Завдання {2}\end{center}}
\vspace{-6mm}

\begin{verbatim}
try:
    t = {14:3, 78:4, 62:9, 87:7, 78:0}
    t[14] = 2
    for x in t.values():
        print(x, sep=' ', end=' ')
except:
    print('ERR')
\end{verbatim}


\end{minipage}
\vline \hspace{8mm}
\begin{minipage}[t]{8.0cm}                
               
{\begin{center}Завдання {3}\end{center}}
\vspace{-6mm}

\begin{verbatim}
a = 0
n = 1

class C:
    a = 2
    
    def __init__(self):
        global a
        self.n = 3
        a = 4
        C.a = 5

obj = C()
try:
    print(a, end=' ')
    print(n, end=' ')
    print(C.a, end=' ')
    print(obj.n)
except:
    print('ERR')
\end{verbatim}

\end{minipage}

               
{\begin{center}Завдання {4}\end{center}}
\vspace{-6mm}
Нехай змінна \kw{s} позначає послідовність цілих, по якій можна ітерувати.

Написати вираз, що перевіряє, чи правда, що всі елементи послідовності є точними квадратами.\par

\vspace{10mm}
               
{\begin{center}Завдання {5}\end{center}}
\vspace{-6mm}
Змінні \kw{next} та \kw{prev} вузла списку позначають наступний та попередній за ним вузол відповідно. Починаючи з голови, у список записано послідовні цілі числа від~$-69$ до~$25$. Нехай змінна \kw{p} позначає вузол, що зберігає значення~$8$.
Яке число зберігається у вузлі \kw{p.next.prev.next.prev.next} ?\par

               
{\begin{center}Завдання {6}\end{center}}
\vspace{-6mm}
Значенням змінної \kw{a} є цілочисловий масив типу $numpy.ndarray$ розмірів $5{\times}4$. На скільки збільшиться сума елементів масиву \kw{a} після виконання
\begin{verbatim}
a.T += 3
\end{verbatim}


Якщо код некоректний, то в якості відповіді зазначити ERROR.\par

               
{\begin{center}Завдання {10}\end{center}}
\vspace{-6mm}

Написати функцію $f$, яка для файлу, заданого шляхом, розв’язує задачу:\\ Текстовий файл містить послідовність чисел $a_1$, $a_2$, $\ldots$, $a_t$, записану у форматі одне число на рядок. Знайти третє найбільше значення, що записане у файлі. Кожне значення, що записане у файлі, враховується тільки один раз. Для послідовності 1, 4, 12, 5, 1, 0 таким значенням буде 4.

В усіх випадках розміри файлів такі, що їх вміст повністю завантажений у пам'ять бути не може. Якість коду також оцінюється.



\vspace{10mm}\par

\newpage
\begin{minipage}[t]{9.7cm}
               
{\begin{center}Завдання {7}\end{center}}
\vspace{-6mm}

\begin{center}
	\adjustimage{max size={0.8\linewidth}{0.8\paperheight}}
	{./15BinTree/trees_exp/70.png}
\end{center}
Указати послідовність міток вершин дерева, зображеного на рис., що утвориться в результаті обходу дерева в прямому порядку. Мітки записувати через пробіл.\par Алгоритм починає роботу з кореня (на рис. це верхня вершина).\par

\end{minipage}
\vline \hspace{8mm}
\begin{minipage}[t]{8.0cm}
               
{\begin{center}Завдання {8}\end{center}}
\vspace{-6mm}
Для графа на рисунку з вершини 4 запущено алгоритм пошуку в глибину, що будує кістякове дерево. Списки суміжності вершин переглядаються алгоритмом за зростанням міток вершин.\par
\begin{center}
	\adjustimage{max size={0.9\linewidth}{0.9\paperheight}}
	{./16Graphs/Traversals/6.png}
\end{center}
\par Які з наведених ребер входять до складу отриманого кістякового дерева?\par
1) (0, 1)\par
2) (2, 6)\par
3) (1, 2)\par
4) (4, 5)\par
5) (0, 6)\par
6) (1, 3)\par
{\vskip 15pt} \par

\end{minipage}

               
{\begin{center}Завдання {9}\end{center}}
\vspace{-6mm}
Файл \code{'1.txt'} містить журнал з рядками, в яких через двокрапку записано поряд\-ко\-вий номер запису в журналі (ціле), код події (ціле), тип події (ERROR, \\WARNING, INFO, STATUS), ім'я користувача (складається з латиниці), час події (фор\-мат як у прикладі), опис події.

Білі символи навколо роздільників не допускаються.

Приклад рядка файлу



\begin{verbatim}
13:404:ERROR:username:2022-05-05 13-26:something happens
\end{verbatim}


Нижче наведено код, який має залишити у файлі в кожному рядку тільки код події та ім'я користувача. Рядок з прикладу має бути замінений на такий



\begin{verbatim}
404:username
\end{verbatim}




Код:



\begin{verbatim}
regex = re.compile(r'XXX')
with open('1.txt', encoding='utf8') as f:
    lst = f.readlines()
for i in range(len(lst)):
    lst[i] = re.sub(regex, YYY, lst[i])
with open('1.txt', 'w', encoding='utf8') as f:
    f.writelines(lst)
\end{verbatim}





Який рядок треба записати на місце \kw{XXX} ?
Який рядок треба записати на місце \kw{YYY} ?\par

%\end {large}
}
\newpage

{{23.05.2025 Контрольна робота, варіант  \No \textbf { 12 }} група К1{\parbox {0.5 cm} \dotfill}  ПІБ {\parbox {2.5 cm} \dotfill}\par}
{
%\begin {large}
\begin{minipage}[t]{9.7cm}
               
{\begin{center}Завдання {1}\end{center}}
\vspace{-6mm}
Укажіть ТИП значення виразу.\par\par
\begin{verbatim}
{x:x for x in range(7)}
\end{verbatim}

\vspace{15mm}
У завданнях 2, 3 вказати, що буде виведено за виконання. Повідомлення інтерпретатора про помилку позначати ERROR

               
{\begin{center}Завдання {2}\end{center}}
\vspace{-6mm}

\begin{verbatim}
try:
    d = {55:2, 71:2, 73:4, 12:4, 73:8}
    d[18] = 8
    for x in d.items():
        print(x, sep=' ', end=' ')
except:
    print('ERR')
\end{verbatim}


\end{minipage}
\vline \hspace{8mm}
\begin{minipage}[t]{8.0cm}                
               
{\begin{center}Завдання {3}\end{center}}
\vspace{-6mm}

\begin{verbatim}
_c = 0

class A:
    _c = 1
    
    def __init__(self):
        self._c = 2
        _c = 3
        A._c = 4

obj = A()
A._c = 5
try:
    print(_c, end=' ')
    print(obj._c, end=' ')
    print(A._c, end=' ')
    print(obj._c, end=' ')
    print(obj._A__c)
except:
    print('ERR')
\end{verbatim}

\end{minipage}

               
{\begin{center}Завдання {4}\end{center}}
\vspace{-6mm}
Нехай змінна \kw{s} позначає дуже довгу послідовність цілих, по якій можна ітерувати.

Написати вираз-генератор, що задає підпослідовність її елементів, що діляться на 3.

Обмеження: усі елементи шуканої підпослідовності одночасно в пам'яті розміщені бути не можуть.\par

\vspace{10mm}
               
{\begin{center}Завдання {5}\end{center}}
\vspace{-6mm}
Змінні \kw{next} та \kw{prev} вузла списку позначають наступний та попередній за ним вузол відповідно. Починаючи з голови, у список записано послідовні цілі числа від~$-20$ до~$46$. Нехай змінна \kw{p} позначає вузол, що зберігає значення~$4$.
Яке число зберігається у вузлі \kw{p.prev.next.prev.next.prev} ?\par

               
{\begin{center}Завдання {6}\end{center}}
\vspace{-6mm}
Значенням змінної \kw{a} є цілочисловий масив типу $numpy.ndarray$ розмірів $6{\times}3$. На скільки збільшиться сума елементів масиву \kw{a} після виконання
\begin{verbatim}
a[2] += 1
\end{verbatim}


Якщо код некоректний, то в якості відповіді зазначити ERROR.\par

               
{\begin{center}Завдання {10}\end{center}}
\vspace{-6mm}

Написати функцію $f$, яка для файлу, заданого шляхом, розв’язує задачу:\\ Кожен з двох текстових файлів містить послідовність чисел, записану у форматі одне число на рядок.  Перевірити, чи задають вони одну ту саму послідовність.

В усіх випадках розміри файлів такі, що їх вміст повністю завантажений у пам'ять бути не може. Якість коду також оцінюється.



\vspace{10mm}\par

\newpage
\begin{minipage}[t]{9.7cm}
               
{\begin{center}Завдання {7}\end{center}}
\vspace{-6mm}

\begin{center}
	\adjustimage{max size={0.8\linewidth}{0.8\paperheight}}
	{./15BinTree/trees_exp/115.png}
\end{center}
Указати послідовність міток вершин дерева, зображеного на рис., що утвориться в результаті обходу дерева в оберненому порядку. Мітки записувати через пробіл.\par Алгоритм починає роботу з кореня (на рис. це верхня вершина).\par

\end{minipage}
\vline \hspace{8mm}
\begin{minipage}[t]{8.0cm}
               
{\begin{center}Завдання {8}\end{center}}
\vspace{-6mm}
Для графа на рисунку з вершини 3 запущено алгоритм пошуку в ширину, що будує кістякове дерево. Списки суміжності вершин переглядаються алгоритмом за зростанням міток вершин.\par
\begin{center}
	\adjustimage{max size={0.9\linewidth}{0.9\paperheight}}
	{./16Graphs/Traversals/309.png}
\end{center}
\par Які з наведених ребер входять до складу отриманого кістякового дерева?\par
1) (2, 4)\par
2) (2, 5)\par
3) (1, 4)\par
4) (0, 4)\par
5) (0, 3)\par
6) (2, 6)\par
{\vskip 15pt} \par

\end{minipage}

               
{\begin{center}Завдання {9}\end{center}}
\vspace{-6mm}
Рядки журналу через кому містять ім'я користувача (складається з латиниці), час події (фор\-мат як у прикладі), тип події (ERROR, WARNING, INFO, STATUS), код події (ціле).

Білі символи навколо роздільників не допускаються.

Приклад такого рядка присвоєно змінній \kw{s}.



\begin{verbatim}
s = 'username,2022-05-05 13:26,ERROR,404'
\end{verbatim}


Нижче наведено код, за виконання якого значенням змінної \kw{out} має стати код події.

Для наведеного прикладу значенням змінної \kw{out} має стати рядок \code{404}



\begin{verbatim}
res = re.fullmatch(r'XXX', s)
s = ''
out = YYY
\end{verbatim}


Код має працювати не тільки для конкретного рядка з прикладу, але й для кожного рядка з журналу.


Який рядок треба записати на місце \kw{XXX} ?
Який рядок треба записати на місце \kw{YYY} ?\par

%\end {large}
}
\newpage

{{23.05.2025 Контрольна робота, варіант  \No \textbf { 13 }} група К1{\parbox {0.5 cm} \dotfill}  ПІБ {\parbox {2.5 cm} \dotfill}\par}
{
%\begin {large}
\begin{minipage}[t]{9.7cm}
               
{\begin{center}Завдання {1}\end{center}}
\vspace{-6mm}
Укажіть ТИП значення виразу.\par\par
\begin{verbatim}
(11,2,3)
\end{verbatim}

\vspace{15mm}
У завданнях 2, 3 вказати, що буде виведено за виконання. Повідомлення інтерпретатора про помилку позначати ERROR

               
{\begin{center}Завдання {2}\end{center}}
\vspace{-6mm}

\begin{verbatim}
try:
    s = {43:0, 70:6, 48:2, 70:7}
    s[70] = 7
    for x in s.values():
        print(x, sep=' ', end=' ')
except:
    print('ERR')
\end{verbatim}


\end{minipage}
\vline \hspace{8mm}
\begin{minipage}[t]{8.0cm}                
               
{\begin{center}Завдання {3}\end{center}}
\vspace{-6mm}

\begin{verbatim}
c = 0
n = 1

class A:
    n = 2
    
    def __init__(self):
        self.c = 3
        n = 4
        A.n = 5

obj = A()
try:
    print(c, end=' ')
    print(n, end=' ')
    print(A.n, end=' ')
    print(obj.c)
except:
    print('ERR')
\end{verbatim}

\end{minipage}

               
{\begin{center}Завдання {4}\end{center}}
\vspace{-6mm}
Нехай змінна \kw{s} позначає дуже довгу послідовність цілих, по якій можна ітерувати.

Написати вираз-генератор, що задає підпослідовність її елементів, що більші за задане число num.

Обмеження: усі елементи шуканої підпослідовності одночасно в пам'яті розміщені бути не можуть.\par

\vspace{10mm}
               
{\begin{center}Завдання {5}\end{center}}
\vspace{-6mm}
Змінні \kw{next} та \kw{prev} вузла списку позначають наступний та попередній за ним вузол відповідно. Починаючи з голови, у список записано послідовні цілі числа від~$-30$ до~$28$. Нехай змінна \kw{p} позначає вузол, що зберігає значення~$7$.
Яке число зберігається у вузлі \kw{p.next.prev.prev.prev.prev} ?\par

               
{\begin{center}Завдання {6}\end{center}}
\vspace{-6mm}
Значенням змінної \kw{a} є цілочисловий масив типу $numpy.ndarray$ розмірів $3{\times}6$. На скільки збільшиться сума елементів масиву \kw{a} після виконання
\begin{verbatim}
a.T[-1] += 2
\end{verbatim}


Якщо код некоректний, то в якості відповіді зазначити ERROR.\par

               
{\begin{center}Завдання {10}\end{center}}
\vspace{-6mm}

Написати функцію $f$, яка для файлу, заданого шляхом, розв’язує задачу:\\ Текстовий файл містить послідовність чисел $a_1$, $a_2$, $\ldots$, $a_t$, записану у форматі одне число на рядок. Знайти третє найбільше значення, що записане у файлі. Кожне значення, що записане у файлі, враховується тільки один раз. Для послідовності 1, 4, 12, 5, 1, 0 таким значенням буде 4.

В усіх випадках розміри файлів такі, що їх вміст повністю завантажений у пам'ять бути не може. Якість коду також оцінюється.



\vspace{10mm}\par

\newpage
\begin{minipage}[t]{9.7cm}
               
{\begin{center}Завдання {7}\end{center}}
\vspace{-6mm}

\begin{center}
	\adjustimage{max size={0.8\linewidth}{0.8\paperheight}}
	{./15BinTree/trees_exp/13.png}
\end{center}
Указати послідовність міток вершин дерева, зображеного на рис., що утвориться в результаті обходу дерева в оберненому порядку. Мітки записувати через пробіл.\par Алгоритм починає роботу з кореня (на рис. це верхня вершина).\par

\end{minipage}
\vline \hspace{8mm}
\begin{minipage}[t]{8.0cm}
               
{\begin{center}Завдання {8}\end{center}}
\vspace{-6mm}
Для графа на рисунку з вершини 5 запущено алгоритм пошуку в глибину, що будує кістякове дерево. Списки суміжності вершин переглядаються алгоритмом за зростанням міток вершин.\par
\begin{center}
	\adjustimage{max size={0.9\linewidth}{0.9\paperheight}}
	{./16Graphs/Traversals/129.png}
\end{center}
\par Які з наведених ребер входять до складу отриманого кістякового дерева?\par
1) (3, 5)\par
2) (0, 3)\par
3) (4, 6)\par
4) (0, 1)\par
5) (2, 5)\par
6) (4, 5)\par
{\vskip 15pt} \par

\end{minipage}

               
{\begin{center}Завдання {9}\end{center}}
\vspace{-6mm}
Файл \code{'1.txt'} містить відомість з рядками, в яких через двокрапку записано поряд\-ковий номер (ціле), назву предмету (знаків пунктуації не містить), прізвище студента, ім'я студента, номер заліковки, оцінку в балах за 100 бальною шкалою.

Білі символи навколо роздільників не допускаються.

Приклад рядка файлу



\begin{verbatim}
13:algebra:surname:name:123456:77
\end{verbatim}


Нижче наведено код, який шукає усі назви предметів, що зустрічаються у файлі, та зберігає їх у змінній \id{codes}.



\begin{verbatim}
codes = set()
regex = re.compile(r'XXX')
with open('1.txt', encoding='utf8') as f:
    for s in f:
        res = re.fullmatch(regex, s)
        if not res: continue
        s = ''
        codes.add(YYY)
\end{verbatim}



Який рядок треба записати на місце \kw{XXX} ?
Який рядок треба записати на місце \kw{YYY} ?\par

%\end {large}
}
\newpage

{{23.05.2025 Контрольна робота, варіант  \No \textbf { 14 }} група К1{\parbox {0.5 cm} \dotfill}  ПІБ {\parbox {2.5 cm} \dotfill}\par}
{
%\begin {large}
\begin{minipage}[t]{9.7cm}
               
{\begin{center}Завдання {1}\end{center}}
\vspace{-6mm}
Укажіть ТИП значення виразу.\par\par
\begin{verbatim}
{x for x in range(7) if x%3==1}
\end{verbatim}

\vspace{15mm}
У завданнях 2, 3 вказати, що буде виведено за виконання. Повідомлення інтерпретатора про помилку позначати ERROR

               
{\begin{center}Завдання {2}\end{center}}
\vspace{-6mm}

\begin{verbatim}
try:
    d = {77:9, 16:9, 84:5, 16:9}
    d[58] = 5
    for x in d :
        print(x, sep=' ', end=' ')
except:
    print('ERR')
\end{verbatim}


\end{minipage}
\vline \hspace{8mm}
\begin{minipage}[t]{8.0cm}                
               
{\begin{center}Завдання {3}\end{center}}
\vspace{-6mm}

\begin{verbatim}
d = 0
m = 1

class D:
    m = 2
    
    def __init__(self):
        global m
        self.d = 3
        m = 4
        D.m = 5

obj = D()
try:
    print(d, end=' ')
    print(m, end=' ')
    print(D.m, end=' ')
    print(obj.d)
except:
    print('ERR')
\end{verbatim}

\end{minipage}

               
{\begin{center}Завдання {4}\end{center}}
\vspace{-6mm}
Нехай змінна \kw{s} позначає послідовність цілих, по якій можна ітерувати.

Написати вираз, що перевіряє, чи правда, що всі елементи послідовності діляться на 11.\par

\vspace{10mm}
               
{\begin{center}Завдання {5}\end{center}}
\vspace{-6mm}
Змінні \kw{next} та \kw{prev} вузла списку позначають наступний та попередній за ним вузол відповідно. Починаючи з голови, у список записано послідовні цілі числа від~$-44$ до~$33$. Нехай змінна \kw{p} позначає вузол, що зберігає значення~$2$.
Яке число зберігається у вузлі \kw{p.prev.next.next.next.next} ?\par

               
{\begin{center}Завдання {6}\end{center}}
\vspace{-6mm}
Значенням змінної \kw{a} є цілочисловий масив типу $numpy.ndarray$ розмірів $7{\times}7$. На скільки збільшиться сума елементів масиву \kw{a} після виконання
\begin{verbatim}
a[0] += 3
\end{verbatim}


Якщо код некоректний, то в якості відповіді зазначити ERROR.\par

               
{\begin{center}Завдання {10}\end{center}}
\vspace{-6mm}

Написати функцію $f$, яка для файлу, заданого шляхом, розв’язує задачу:\\ Текстовий файл містить послідовність чисел $a_1$, $a_2$, $\ldots$, $a_t$, записану у форматі одне число на рядок. Знайти третє найбільше значення, що записане у файлі. Кожне значення, що записане у файлі, враховується стільки разів, скільки входить у файл. Для послідовності 2, 4, 12, 5, 2, 0 таким значенням буде 2.

В усіх випадках розміри файлів такі, що їх вміст повністю завантажений у пам'ять бути не може. Якість коду також оцінюється.



\vspace{10mm}\par

\newpage
\begin{minipage}[t]{9.7cm}
               
{\begin{center}Завдання {7}\end{center}}
\vspace{-6mm}

\begin{center}
	\adjustimage{max size={0.8\linewidth}{0.8\paperheight}}
	{./15BinTree/trees_exp/48.png}
\end{center}
Указати послідовність міток вершин дерева, зображеного на рис., що утвориться в результаті обходу дерева в прямому порядку. Мітки записувати через пробіл.\par Алгоритм починає роботу з кореня (на рис. це верхня вершина).\par

\end{minipage}
\vline \hspace{8mm}
\begin{minipage}[t]{8.0cm}
               
{\begin{center}Завдання {8}\end{center}}
\vspace{-6mm}
Для графа на рисунку з вершини 2 запущено алгоритм пошуку в ширину, що будує кістякове дерево. Списки суміжності вершин переглядаються алгоритмом за зростанням міток вершин.\par
\begin{center}
	\adjustimage{max size={0.9\linewidth}{0.9\paperheight}}
	{./16Graphs/Traversals/133.png}
\end{center}
\par Які з наведених ребер входять до складу отриманого кістякового дерева?\par
1) (0, 1)\par
2) (1, 2)\par
3) (4, 6)\par
4) (4, 5)\par
5) (5, 6)\par
6) (0, 4)\par
{\vskip 15pt} \par

\end{minipage}

               
{\begin{center}Завдання {9}\end{center}}
\vspace{-6mm}
Рядки відомості через двокрапку містять порядковий номер (ціле), назву предмету (знаків пунктуації не містить), прізвище студента, ім'я студента, оцінку в балах за 100 бальною шкалою.

Білі символи навколо роздільників не допускаються.

Приклад такого рядка присвоєно змінній \kw{s}.



\begin{verbatim}
s = '13:algebra:surname:name:77'
\end{verbatim}


Нижче наведено код, за виконання якого значенням змінної \kw{out} має стати прізвище студента.

Для наведеного прикладу значенням змінної \kw{out} має стати рядок \code{surname}



\begin{verbatim}
res = re.fullmatch(r'XXX', s)
s = ''
out = YYY
\end{verbatim}


Код має працювати не тільки для конкретного рядка з прикладу, але й для кожного рядка з відомості.


Який рядок треба записати на місце \kw{XXX} ?
Який рядок треба записати на місце \kw{YYY} ?\par

%\end {large}
}
\newpage

{{23.05.2025 Контрольна робота, варіант  \No \textbf { 15 }} група К1{\parbox {0.5 cm} \dotfill}  ПІБ {\parbox {2.5 cm} \dotfill}\par}
{
%\begin {large}
\begin{minipage}[t]{9.7cm}
               
{\begin{center}Завдання {1}\end{center}}
\vspace{-6mm}
Укажіть ТИП значення виразу.\par\par
\begin{verbatim}
[7,2,3,4]
\end{verbatim}

\vspace{15mm}
У завданнях 2, 3 вказати, що буде виведено за виконання. Повідомлення інтерпретатора про помилку позначати ERROR

               
{\begin{center}Завдання {2}\end{center}}
\vspace{-6mm}

\begin{verbatim}
try:
    t = {72:9, 84:9, 73:7, 28:4, 84:1}
    t[72] = 6
    for x, y in t.items():
        print(x, y, sep=' ', end=' ')
except:
    print('ERR')
\end{verbatim}


\end{minipage}
\vline \hspace{8mm}
\begin{minipage}[t]{8.0cm}                
               
{\begin{center}Завдання {3}\end{center}}
\vspace{-6mm}

\begin{verbatim}
d = 0

class C:
    d = 1
    
    def __init__(self):
        self.d = 2
        d = 3
        C.d = 4

obj = C()
obj.d = 5
try:
    print(d, end=' ')
    print(obj.d, end=' ')
    print(C.d, end=' ')
    print(obj.d, end=' ')
    print(obj._C__d)
except:
    print('ERR')
\end{verbatim}

\end{minipage}

               
{\begin{center}Завдання {4}\end{center}}
\vspace{-6mm}
Нехай змінна \kw{s} позначає послідовність цілих, по якій можна ітерувати.

Написати вираз, що перевіряє, чи правда, що всі елементи послідовності більші за задане число num.\par

\vspace{10mm}
               
{\begin{center}Завдання {5}\end{center}}
\vspace{-6mm}
Змінні \kw{next} та \kw{prev} вузла списку позначають наступний та попередній за ним вузол відповідно. Починаючи з голови, у список записано послідовні цілі числа від~$-18$ до~$29$. Нехай змінна \kw{p} позначає вузол, що зберігає значення~$0$.
Яке число зберігається у вузлі \kw{p.next.next.next.prev.prev} ?\par

               
{\begin{center}Завдання {6}\end{center}}
\vspace{-6mm}
Значенням змінної \kw{a} є цілочисловий масив типу $numpy.ndarray$ розмірів $4{\times}5$. На скільки збільшиться сума елементів масиву \kw{a} після виконання
\begin{verbatim}
a += 1
\end{verbatim}


Якщо код некоректний, то в якості відповіді зазначити ERROR.\par

               
{\begin{center}Завдання {10}\end{center}}
\vspace{-6mm}

Написати функцію $f$, яка для файлу, заданого шляхом, розв’язує задачу:\\ Текстовий файл містить послідовність чисел $a_1$, $a_2$, $\ldots$, $a_t$, записану у форматі одне число на рядок. Знайти третє найменше значення, що записане у файлі. Кожне значення, що записане у файлі, враховується стільки разів, скільки входить у файл. Для послідовності 2, 4, 12, 5, 2, 0 таким значенням буде 2.

В усіх випадках розміри файлів такі, що їх вміст повністю завантажений у пам'ять бути не може. Якість коду також оцінюється.



\vspace{10mm}\par

\newpage
\begin{minipage}[t]{9.7cm}
               
{\begin{center}Завдання {7}\end{center}}
\vspace{-6mm}

\begin{center}
	\adjustimage{max size={0.8\linewidth}{0.8\paperheight}}
	{./15BinTree/trees_exp/60.png}
\end{center}
Указати послідовність міток вершин дерева, зображеного на рис., що утвориться в результаті обходу дерева в прямому порядку. Мітки записувати через пробіл.\par Алгоритм починає роботу з кореня (на рис. це верхня вершина).\par

\end{minipage}
\vline \hspace{8mm}
\begin{minipage}[t]{8.0cm}
               
{\begin{center}Завдання {8}\end{center}}
\vspace{-6mm}
Для графа на рисунку з вершини 2 запущено алгоритм пошуку в глибину, що будує кістякове дерево. Списки суміжності вершин переглядаються алгоритмом за зростанням міток вершин.\par
\begin{center}
	\adjustimage{max size={0.9\linewidth}{0.9\paperheight}}
	{./16Graphs/Traversals/173.png}
\end{center}
\par Які з наведених ребер входять до складу отриманого кістякового дерева?\par
1) (2, 3)\par
2) (4, 5)\par
3) (2, 6)\par
4) (0, 4)\par
5) (4, 6)\par
6) (1, 4)\par
{\vskip 15pt} \par

\end{minipage}

               
{\begin{center}Завдання {9}\end{center}}
\vspace{-6mm}
Файл \code{'1.txt'} містить журнал з рядками, в яких через кому записано поряд\-ко\-вий номер запису (ціле), ім'я користувача (складається з латиниці), час події (фор\-мат як у прикладі), тип події (ERROR, WARNING, INFO, STATUS), код події (ціле).

Білі символи навколо роздільників не допускаються.

Приклад рядка журналу



\begin{verbatim}
12,username,2022-05-05 13:26,ERROR,404
\end{verbatim}


Нижче наведено код, який шукає усі попередження (подія WARNING), що зустрі\-ча\-ються в журналі і відбулися під час обідньої перерви деканату (з 13:00 по 13:59 включно), та зберігає їх коди в змінній \kw{codes}.



\begin{verbatim}
codes = set()
regex = re.compile(r'XXX')
with open('1.txt', encoding='utf8') as f:
    for s in f:
        res = re.fullmatch(regex, s)
        if not res: continue
        s = ''
        codes.add(YYY)
\end{verbatim}





Який рядок треба записати на місце \kw{XXX} ?
Який рядок треба записати на місце \kw{YYY} ?\par

%\end {large}
}
\newpage

{{23.05.2025 Контрольна робота, варіант  \No \textbf { 16 }} група К1{\parbox {0.5 cm} \dotfill}  ПІБ {\parbox {2.5 cm} \dotfill}\par}
{
%\begin {large}
\begin{minipage}[t]{9.7cm}
               
{\begin{center}Завдання {1}\end{center}}
\vspace{-6mm}
Укажіть ТИП значення виразу.\par\par
\begin{verbatim}
{1:3, 5:8}
\end{verbatim}

\vspace{15mm}
У завданнях 2, 3 вказати, що буде виведено за виконання. Повідомлення інтерпретатора про помилку позначати ERROR

               
{\begin{center}Завдання {2}\end{center}}
\vspace{-6mm}

\begin{verbatim}
try:
    d = {31:6, 41:8, 10:9, 40:1, 31:9}
    d[83] = 5
    for x in d.items():
        print(x, sep=' ', end=' ')
except:
    print('ERR')
\end{verbatim}


\end{minipage}
\vline \hspace{8mm}
\begin{minipage}[t]{8.0cm}                
               
{\begin{center}Завдання {3}\end{center}}
\vspace{-6mm}

\begin{verbatim}
__a = 0

class C:
    __a = 1
    
    def __init__(self):
        self.__a = 2
        __a = 3
        C.__a = 4

obj = C()
obj.__a = 5
try:
    print(__a, end=' ')
    print(obj.__a, end=' ')
    print(C.__a, end=' ')
    print(obj.__a, end=' ')
    print(obj._C__a)
except:
    print('ERR')
\end{verbatim}

\end{minipage}

               
{\begin{center}Завдання {4}\end{center}}
\vspace{-6mm}
Нехай змінна \kw{s} позначає дуже довгу послідовність, по якій можна ітерувати.

Написати вираз-генератор, що задає підпослідовність її елементів, що не є числами типу \kw{float}.

Обмеження: усі елементи шуканої підпослідовності одночасно в пам'яті розміщені бути не можуть.\par

\vspace{10mm}
               
{\begin{center}Завдання {5}\end{center}}
\vspace{-6mm}
Змінні \kw{next} та \kw{prev} вузла списку позначають наступний та попередній за ним вузол відповідно. Починаючи з голови, у список записано послідовні цілі числа від~$-60$ до~$20$. Нехай змінна \kw{p} позначає вузол, що зберігає значення~$-3$.
Яке число зберігається у вузлі \kw{p.prev.prev.prev.next.next} ?\par

               
{\begin{center}Завдання {6}\end{center}}
\vspace{-6mm}
Значенням змінної \kw{a} є цілочисловий масив типу $numpy.ndarray$ розмірів $6{\times}4$. На скільки збільшиться сума елементів масиву \kw{a} після виконання
\begin{verbatim}
a.T[-3] += 1
\end{verbatim}


Якщо код некоректний, то в якості відповіді зазначити ERROR.\par

               
{\begin{center}Завдання {10}\end{center}}
\vspace{-6mm}

Написати функцію $f$, яка для файлу, заданого шляхом, розв’язує задачу:\\ Текстовий файл містить послідовність чисел $a_1$, $a_2$, $\ldots$, $a_t$, записану у форматі одне число на рядок. Знайти середнє арифметичне чисел $a_1$, $a_4$, $a_7$, $\ldots$ (починаючи з першого через два).

В усіх випадках розміри файлів такі, що їх вміст повністю завантажений у пам'ять бути не може. Якість коду також оцінюється.



\vspace{10mm}\par

\newpage
\begin{minipage}[t]{9.7cm}
               
{\begin{center}Завдання {7}\end{center}}
\vspace{-6mm}

\begin{center}
	\adjustimage{max size={0.8\linewidth}{0.8\paperheight}}
	{./15BinTree/trees_exp/59.png}
\end{center}
Указати послідовність міток вершин дерева, зображеного на рис., що утвориться в результаті обходу дерева в оберненому порядку. Мітки записувати через пробіл.\par Алгоритм починає роботу з кореня (на рис. це верхня вершина).\par

\end{minipage}
\vline \hspace{8mm}
\begin{minipage}[t]{8.0cm}
               
{\begin{center}Завдання {8}\end{center}}
\vspace{-6mm}
Для графа на рисунку з вершини 2 запущено алгоритм пошуку в глибину, що будує кістякове дерево. Списки суміжності вершин переглядаються алгоритмом за зростанням міток вершин.\par
\begin{center}
	\adjustimage{max size={0.9\linewidth}{0.9\paperheight}}
	{./16Graphs/Traversals/27.png}
\end{center}
\par Які з наведених ребер входять до складу отриманого кістякового дерева?\par
1) (2, 6)\par
2) (1, 2)\par
3) (1, 6)\par
4) (2, 3)\par
5) (3, 5)\par
6) (1, 3)\par
{\vskip 15pt} \par

\end{minipage}

               
{\begin{center}Завдання {9}\end{center}}
\vspace{-6mm}
Файл \code{'1.txt'} містить відомість з рядками, в яких через двокрапку записано поряд\-ковий номер (ціле), назву предмету (знаків пунктуації не містить), прізвище студента, ім'я студента, номер заліковки, оцінку в балах за 100 бальною шкалою.

Білі символи навколо роздільників не допускаються.

Приклад рядка файлу



\begin{verbatim}
13:algebra:surname:name:123456:77
\end{verbatim}


Нижче наведено код, який шукає усі предмети, який складав студент із заліковкою \code{777}, та зберігає їх типи в змінній \kw{codes}.



\begin{verbatim}
codes = set()
regex = re.compile(r'XXX')
with open('1.txt', encoding='utf8') as f:
    for s in f:
        res = re.fullmatch(regex, s)
        if not res: continue
        s = ''
        codes.add(YYY)
\end{verbatim}



Який рядок треба записати на місце \kw{XXX} ?
Який рядок треба записати на місце \kw{YYY} ?\par

%\end {large}
}
\newpage

{{23.05.2025 Контрольна робота, варіант  \No \textbf { 17 }} група К1{\parbox {0.5 cm} \dotfill}  ПІБ {\parbox {2.5 cm} \dotfill}\par}
{
%\begin {large}
\begin{minipage}[t]{9.7cm}
               
{\begin{center}Завдання {1}\end{center}}
\vspace{-6mm}
Укажіть ТИП значення виразу.\par\par
\begin{verbatim}
{21,2,4}
\end{verbatim}

\vspace{15mm}
У завданнях 2, 3 вказати, що буде виведено за виконання. Повідомлення інтерпретатора про помилку позначати ERROR

               
{\begin{center}Завдання {2}\end{center}}
\vspace{-6mm}

\begin{verbatim}
try:
    s = {38:7, 63:0, 45:0, 82:0, 45:0}
    s[18] = 0
    for x in s.keys():
        print(x, sep=' ', end=' ')
except:
    print('ERR')
\end{verbatim}


\end{minipage}
\vline \hspace{8mm}
\begin{minipage}[t]{8.0cm}                
               
{\begin{center}Завдання {3}\end{center}}
\vspace{-6mm}

\begin{verbatim}
__c = 0

class A:
    __c = 1
    
    def __init__(self):
        self.__c = 2
        __c = 3
        A.__c = 4

obj = A()
A.__c = 5
try:
    print(__c, end=' ')
    print(obj.__c, end=' ')
    print(A.__c, end=' ')
    print(obj.__c, end=' ')
    print(obj._A__c)
except:
    print('ERR')
\end{verbatim}

\end{minipage}

               
{\begin{center}Завдання {4}\end{center}}
\vspace{-6mm}
Нехай змінна \kw{s} позначає послідовність цілих, по якій можна ітерувати.

Написати вираз, що перевіряє, чи правда, що всі елементи послідовності є додатни\-ми.\par

\vspace{10mm}
               
{\begin{center}Завдання {5}\end{center}}
\vspace{-6mm}
Змінні \kw{next} та \kw{prev} вузла списку позначають наступний та попередній за ним вузол відповідно. Починаючи з голови, у список записано послідовні цілі числа від~$-57$ до~$49$. Нехай змінна \kw{p} позначає вузол, що зберігає значення~$-2$.
Яке число зберігається у вузлі \kw{p.next.next.next.prev.next} ?\par

               
{\begin{center}Завдання {6}\end{center}}
\vspace{-6mm}
Значенням змінної \kw{a} є цілочисловий масив типу $numpy.ndarray$ розмірів $5{\times}6$. На скільки збільшиться сума елементів масиву \kw{a} після виконання
\begin{verbatim}
a[1] += 3
\end{verbatim}


Якщо код некоректний, то в якості відповіді зазначити ERROR.\par

               
{\begin{center}Завдання {10}\end{center}}
\vspace{-6mm}

Написати функцію $f$, яка для файлу, заданого шляхом, розв’язує задачу:\\ Текстовий файл містить послідовність чисел $a_1$, $a_2$, $\ldots$, $a_t$, записану у форматі одне число на рядок. Знайти третє найменше значення, що записане у файлі. Кожне значення, що записане у файлі, враховується тільки один раз. \par
Для послідовності 1, 4, 12, 5, 1, 0 таким значенням буде 4.

В усіх випадках розміри файлів такі, що їх вміст повністю завантажений у пам'ять бути не може. Якість коду також оцінюється.



\vspace{10mm}\par

\newpage
\begin{minipage}[t]{9.7cm}
               
{\begin{center}Завдання {7}\end{center}}
\vspace{-6mm}

\begin{center}
	\adjustimage{max size={0.8\linewidth}{0.8\paperheight}}
	{./15BinTree/trees_exp/104.png}
\end{center}
Указати послідовність міток вершин дерева, зображеного на рис., що утвориться в результаті обходу дерева в прямому порядку. Мітки записувати через пробіл.\par Алгоритм починає роботу з кореня (на рис. це верхня вершина).\par

\end{minipage}
\vline \hspace{8mm}
\begin{minipage}[t]{8.0cm}
               
{\begin{center}Завдання {8}\end{center}}
\vspace{-6mm}
Для графа на рисунку з вершини 5 запущено алгоритм пошуку в глибину, що будує кістякове дерево. Списки суміжності вершин переглядаються алгоритмом за зростанням міток вершин.\par
\begin{center}
	\adjustimage{max size={0.9\linewidth}{0.9\paperheight}}
	{./16Graphs/Traversals/31.png}
\end{center}
\par Які з наведених ребер входять до складу отриманого кістякового дерева?\par
1) (2, 3)\par
2) (1, 5)\par
3) (0, 5)\par
4) (3, 4)\par
5) (2, 6)\par
6) (0, 6)\par
{\vskip 15pt} \par

\end{minipage}

               
{\begin{center}Завдання {9}\end{center}}
\vspace{-6mm}
Файл \code{'1.txt'} містить журнал з рядками, в яких через двокрапку записано поряд\-ко\-вий номер запису в журналі (ціле), тип події (ERROR, WARNING, INFO, STATUS), ім'я користувача (складається з латиниці), час події (фор\-мат як у прикладі), код події (ціле).

Білі символи навколо роздільників не допускаються.

Приклад рядка файлу



\begin{verbatim}
13:ERROR:username:2022-05-05 13-26:404
\end{verbatim}


Нижче наведено код, який шукає усі коди помилок (подія ERROR), що зустріча\-ють\-ся в журналі, та зберігає їх у змінній \kw{codes}.



\begin{verbatim}
codes = set()
regex = re.compile(r'XXX')
with open('1.txt', encoding='utf8') as f:
    for s in f:
        res = re.fullmatch(regex, s)
        if not res: continue
        s = ''
        codes.add(YYY)
\end{verbatim}



Який рядок треба записати на місце \kw{XXX} ?
Який рядок треба записати на місце \kw{YYY} ?\par

%\end {large}
}
\newpage

{{23.05.2025 Контрольна робота, варіант  \No \textbf { 18 }} група К1{\parbox {0.5 cm} \dotfill}  ПІБ {\parbox {2.5 cm} \dotfill}\par}
{
%\begin {large}
\begin{minipage}[t]{9.7cm}
               
{\begin{center}Завдання {1}\end{center}}
\vspace{-6mm}
Укажіть ТИП значення виразу.\par\par
\begin{verbatim}
{x for x in range(7) if x%3==1}
\end{verbatim}

\vspace{15mm}
У завданнях 2, 3 вказати, що буде виведено за виконання. Повідомлення інтерпретатора про помилку позначати ERROR

               
{\begin{center}Завдання {2}\end{center}}
\vspace{-6mm}

\begin{verbatim}
try:
    t = {50:3, 70:4, 87:3, 36:2, 87:6}
    t[87] = 2
    for x in t.items():
        print(*x, sep=' ', end=' ')
except:
    print('ERR')
\end{verbatim}


\end{minipage}
\vline \hspace{8mm}
\begin{minipage}[t]{8.0cm}                
               
{\begin{center}Завдання {3}\end{center}}
\vspace{-6mm}

\begin{verbatim}
b = 0
h = 1

class B:
    b = 2
    
    def __init__(self):
        self.h = 3
        b = 4
        B.b = 5

obj = B()
try:
    print(b, end=' ')
    print(h, end=' ')
    print(B.b, end=' ')
    print(obj.h)
except:
    print('ERR')
\end{verbatim}

\end{minipage}

               
{\begin{center}Завдання {4}\end{center}}
\vspace{-6mm}
Нехай змінна \kw{s} позначає дуже довгу послідовність, по якій можна ітерувати.

Написати вираз-генератор, що задає підпослідовність її елементів, що є числами типу \kw{float}.

Обмеження: усі елементи шуканої підпослідовності одночасно в пам'яті розміщені бути не можуть.\par

\vspace{10mm}
               
{\begin{center}Завдання {5}\end{center}}
\vspace{-6mm}
Змінні \kw{next} та \kw{prev} вузла списку позначають наступний та попередній за ним вузол відповідно. Починаючи з голови, у список записано послідовні цілі числа від~$-59$ до~$48$. Нехай змінна \kw{p} позначає вузол, що зберігає значення~$1$.
Яке число зберігається у вузлі \kw{p.prev.prev.prev.next.prev} ?\par

               
{\begin{center}Завдання {6}\end{center}}
\vspace{-6mm}
Значенням змінної \kw{a} є цілочисловий масив типу $numpy.ndarray$ розмірів $7{\times}3$. На скільки збільшиться сума елементів масиву \kw{a} після виконання
\begin{verbatim}
a.T[-1] += 2
\end{verbatim}


Якщо код некоректний, то в якості відповіді зазначити ERROR.\par

               
{\begin{center}Завдання {10}\end{center}}
\vspace{-6mm}

Написати функцію $f$, яка для файлу, заданого шляхом, розв’язує задачу:\\ Кожен з двох текстових файлів містить послідовність чисел, записану у форматі одне число на рядок.  Перевірити, чи задають вони одну ту саму послідовність.

В усіх випадках розміри файлів такі, що їх вміст повністю завантажений у пам'ять бути не може. Якість коду також оцінюється.



\vspace{10mm}\par

\newpage
\begin{minipage}[t]{9.7cm}
               
{\begin{center}Завдання {7}\end{center}}
\vspace{-6mm}

\begin{center}
	\adjustimage{max size={0.8\linewidth}{0.8\paperheight}}
	{./15BinTree/trees_exp/37.png}
\end{center}
Указати послідовність міток вершин дерева, зображеного на рис., що утвориться в результаті обходу дерева в оберненому порядку. Мітки записувати через пробіл.\par Алгоритм починає роботу з кореня (на рис. це верхня вершина).\par

\end{minipage}
\vline \hspace{8mm}
\begin{minipage}[t]{8.0cm}
               
{\begin{center}Завдання {8}\end{center}}
\vspace{-6mm}
Для графа на рисунку з вершини 0 запущено алгоритм пошуку в ширину, що будує кістякове дерево. Списки суміжності вершин переглядаються алгоритмом за зростанням міток вершин.\par
\begin{center}
	\adjustimage{max size={0.9\linewidth}{0.9\paperheight}}
	{./16Graphs/Traversals/235.png}
\end{center}
\par Які з наведених ребер входять до складу отриманого кістякового дерева?\par
1) (0, 2)\par
2) (3, 5)\par
3) (0, 3)\par
4) (2, 4)\par
5) (0, 5)\par
6) (1, 2)\par
{\vskip 15pt} \par

\end{minipage}

               
{\begin{center}Завдання {9}\end{center}}
\vspace{-6mm}
Файл \code{'1.txt'} містить відомість з рядками, в яких через кому записа\-но назву предмету (знаків пунктуації не містить), прізвище студента, ім'я студента, номер залі\-ков\-ки, оцінку в балах за 100 бальною шкалою.

Білі символи навколо роздільників не допускаються.

Приклад рядка файлу



\begin{verbatim}
algebra,surname,name,123456,77
\end{verbatim}


Нижче наведено код, який має замінити у файлі усі номери заліковок на рядок \code{***}.



\begin{verbatim}
regex = re.compile(r'XXX')
with open('1.txt', encoding='utf8') as f:
    lst = f.readlines()
for i in range(len(lst)):
    lst[i] = re.sub(regex, YYY, lst[i])
with open('1.txt', 'w', encoding='utf8') as f:
    f.writelines(lst)
\end{verbatim}





Який рядок треба записати на місце \kw{XXX} ?
Який рядок треба записати на місце \kw{YYY} ?\par

%\end {large}
}
\newpage

{{23.05.2025 Контрольна робота, варіант  \No \textbf { 19 }} група К1{\parbox {0.5 cm} \dotfill}  ПІБ {\parbox {2.5 cm} \dotfill}\par}
{
%\begin {large}
\begin{minipage}[t]{9.7cm}
               
{\begin{center}Завдання {1}\end{center}}
\vspace{-6mm}
Укажіть ТИП значення виразу.\par\par
\begin{verbatim}
(2, 7, -4)
\end{verbatim}

\vspace{15mm}
У завданнях 2, 3 вказати, що буде виведено за виконання. Повідомлення інтерпретатора про помилку позначати ERROR

               
{\begin{center}Завдання {2}\end{center}}
\vspace{-6mm}

\begin{verbatim}
try:
    t = {65:6, 81:1, 72:0, 72:4}
    t[81] = 1
    for x, y in t.items():
        print(x, y, sep=' ', end=' ')
except:
    print('ERR')
\end{verbatim}


\end{minipage}
\vline \hspace{8mm}
\begin{minipage}[t]{8.0cm}                
               
{\begin{center}Завдання {3}\end{center}}
\vspace{-6mm}

\begin{verbatim}
b = 0

class B:
    b = 1
    
    def __init__(self):
        self.b = 2
        b = 3
        B.b = 4

obj = B()
obj.b = 5
try:
    print(b, end=' ')
    print(obj.b, end=' ')
    print(B.b, end=' ')
    print(obj.b, end=' ')
    print(obj._B__b)
except:
    print('ERR')
\end{verbatim}

\end{minipage}

               
{\begin{center}Завдання {4}\end{center}}
\vspace{-6mm}
Нехай змінна \kw{s} позначає дуже довгу послідовність, по якій можна ітерувати.

Написати вираз-генератор, що задає підпослідовність її елементів, що не є числами типу \kw{int}.

Обмеження: усі елементи шуканої підпослідовності одночасно в пам'яті розміщені бути не можуть.\par

\vspace{10mm}
               
{\begin{center}Завдання {5}\end{center}}
\vspace{-6mm}
Змінні \kw{next} та \kw{prev} вузла списку позначають наступний та попередній за ним вузол відповідно. Починаючи з голови, у список записано послідовні цілі числа від~$-48$ до~$36$. Нехай змінна \kw{p} позначає вузол, що зберігає значення~$-4$.
Яке число зберігається у вузлі \kw{p.next.next.prev.prev.prev} ?\par

               
{\begin{center}Завдання {6}\end{center}}
\vspace{-6mm}
Значенням змінної \kw{a} є цілочисловий масив типу $numpy.ndarray$ розмірів $3{\times}5$. На скільки збільшиться сума елементів масиву \kw{a} після виконання
\begin{verbatim}
a.T[2, 3] += 1
\end{verbatim}


Якщо код некоректний, то в якості відповіді зазначити ERROR.\par

               
{\begin{center}Завдання {10}\end{center}}
\vspace{-6mm}

Написати функцію $f$, яка для файлу, заданого шляхом, розв’язує задачу:\\ Текстовий файл містить послідовність чисел $a_1$, $a_2$, $\ldots$, $a_t$, записану у форматі одне число на рядок. Знайти третє найменше значення, що записане у файлі. Кожне значення, що записане у файлі, враховується тільки один раз. \par
Для послідовності 1, 4, 12, 5, 1, 0 таким значенням буде 4.

В усіх випадках розміри файлів такі, що їх вміст повністю завантажений у пам'ять бути не може. Якість коду також оцінюється.



\vspace{10mm}\par

\newpage
\begin{minipage}[t]{9.7cm}
               
{\begin{center}Завдання {7}\end{center}}
\vspace{-6mm}

\begin{center}
	\adjustimage{max size={0.8\linewidth}{0.8\paperheight}}
	{./15BinTree/trees_exp/97.png}
\end{center}
Указати послідовність міток вершин дерева, зображеного на рис., що утвориться в результаті обходу дерева в оберненому порядку. Мітки записувати через пробіл.\par Алгоритм починає роботу з кореня (на рис. це верхня вершина).\par

\end{minipage}
\vline \hspace{8mm}
\begin{minipage}[t]{8.0cm}
               
{\begin{center}Завдання {8}\end{center}}
\vspace{-6mm}
Для графа на рисунку з вершини 6 запущено алгоритм пошуку в ширину, що будує кістякове дерево. Списки суміжності вершин переглядаються алгоритмом за зростанням міток вершин.\par
\begin{center}
	\adjustimage{max size={0.9\linewidth}{0.9\paperheight}}
	{./16Graphs/Traversals/57.png}
\end{center}
\par Які з наведених ребер входять до складу отриманого кістякового дерева?\par
1) (2, 4)\par
2) (0, 2)\par
3) (0, 5)\par
4) (4, 5)\par
5) (1, 3)\par
6) (1, 2)\par
{\vskip 15pt} \par

\end{minipage}

               
{\begin{center}Завдання {9}\end{center}}
\vspace{-6mm}
Файл \code{'1.txt'} містить журнал з рядками, в яких через двокрапку записано поряд\-ко\-вий номер запису в журналі (ціле), тип події (ERROR, WARNING, INFO, STATUS), ім'я користувача (складається з латиниці), час події (фор\-мат як у прикладі), код події (ціле).

Білі символи навколо роздільників не допускаються.

Приклад рядка файлу



\begin{verbatim}
13:ERROR:username:2022-05-05 13-26:404
\end{verbatim}


Нижче наведено код, який має замінити у файлі імена користувачів на рядок \code{unknown}.



\begin{verbatim}
regex = re.compile(r'XXX')
with open('1.txt', encoding='utf8') as f:
    lst = f.readlines()
for i in range(len(lst)):
    lst[i] = re.sub(regex, YYY, lst[i])
with open('1.txt', 'w', encoding='utf8') as f:
    f.writelines(lst)
\end{verbatim}





Який рядок треба записати на місце \kw{XXX} ?
Який рядок треба записати на місце \kw{YYY} ?\par

%\end {large}
}
\newpage

{{23.05.2025 Контрольна робота, варіант  \No \textbf { 20 }} група К1{\parbox {0.5 cm} \dotfill}  ПІБ {\parbox {2.5 cm} \dotfill}\par}
{
%\begin {large}
\begin{minipage}[t]{9.7cm}
               
{\begin{center}Завдання {1}\end{center}}
\vspace{-6mm}
Укажіть ТИП значення виразу.\par\par
\begin{verbatim}
{y:y*y for y in range(-7,7)}
\end{verbatim}

\vspace{15mm}
У завданнях 2, 3 вказати, що буде виведено за виконання. Повідомлення інтерпретатора про помилку позначати ERROR

               
{\begin{center}Завдання {2}\end{center}}
\vspace{-6mm}

\begin{verbatim}
try:
    d = {83:9, 66:0, 34:7, 29:0, 25:6}
    d[21] = 3
    for x in d.values():
        print(x, sep=' ', end=' ')
except:
    print('ERR')
\end{verbatim}


\end{minipage}
\vline \hspace{8mm}
\begin{minipage}[t]{8.0cm}                
               
{\begin{center}Завдання {3}\end{center}}
\vspace{-6mm}

\begin{verbatim}
_d = 0

class D:
    _d = 1
    
    def __init__(self):
        self._d = 2
        _d = 3
        D._d = 4

obj = D()
D._d = 5
try:
    print(_d, end=' ')
    print(obj._d, end=' ')
    print(D._d, end=' ')
    print(obj._d, end=' ')
    print(obj._D__d)
except:
    print('ERR')
\end{verbatim}

\end{minipage}

               
{\begin{center}Завдання {4}\end{center}}
\vspace{-6mm}
Нехай змінна \kw{s} позначає дуже довгу послідовність цілих, по якій можна ітерувати.

Написати вираз-генератор, що задає підпослідовність її елементів, що є додатни\-ми.

Обмеження: усі елементи шуканої підпослідовності одночасно в пам'яті розміщені бути не можуть.\par

\vspace{10mm}
               
{\begin{center}Завдання {5}\end{center}}
\vspace{-6mm}
Змінні \kw{next} та \kw{prev} вузла списку позначають наступний та попередній за ним вузол відповідно. Починаючи з голови, у список записано послідовні цілі числа від~$-28$ до~$43$. Нехай змінна \kw{p} позначає вузол, що зберігає значення~$4$.
Яке число зберігається у вузлі \kw{p.prev.prev.next.next.next} ?\par

               
{\begin{center}Завдання {6}\end{center}}
\vspace{-6mm}
Значенням змінної \kw{a} є цілочисловий масив типу $numpy.ndarray$ розмірів $4{\times}7$. На скільки збільшиться сума елементів масиву \kw{a} після виконання
\begin{verbatim}
a[-2] += 3
\end{verbatim}


Якщо код некоректний, то в якості відповіді зазначити ERROR.\par

               
{\begin{center}Завдання {10}\end{center}}
\vspace{-6mm}

Написати функцію $f$, яка для файлу, заданого шляхом, розв’язує задачу:\\ Текстовий файл містить послідовність чисел $a_1$, $a_2$, $\ldots$, $a_t$, записану у форматі одне число на рядок. Знайти середнє арифметичне чисел $a_1$, $a_4$, $a_7$, $\ldots$ (починаючи з першого через два).

В усіх випадках розміри файлів такі, що їх вміст повністю завантажений у пам'ять бути не може. Якість коду також оцінюється.



\vspace{10mm}\par

\newpage
\begin{minipage}[t]{9.7cm}
               
{\begin{center}Завдання {7}\end{center}}
\vspace{-6mm}

\begin{center}
	\adjustimage{max size={0.8\linewidth}{0.8\paperheight}}
	{./15BinTree/trees_exp/106.png}
\end{center}
Указати послідовність міток вершин дерева, зображеного на рис., що утвориться в результаті обходу дерева в прямому порядку. Мітки записувати через пробіл.\par Алгоритм починає роботу з кореня (на рис. це верхня вершина).\par

\end{minipage}
\vline \hspace{8mm}
\begin{minipage}[t]{8.0cm}
               
{\begin{center}Завдання {8}\end{center}}
\vspace{-6mm}
Для графа на рисунку з вершини 5 запущено алгоритм пошуку в глибину, що будує кістякове дерево. Списки суміжності вершин переглядаються алгоритмом за зростанням міток вершин.\par
\begin{center}
	\adjustimage{max size={0.9\linewidth}{0.9\paperheight}}
	{./16Graphs/Traversals/338.png}
\end{center}
\par Які з наведених ребер входять до складу отриманого кістякового дерева?\par
1) (3, 6)\par
2) (1, 2)\par
3) (1, 3)\par
4) (0, 4)\par
5) (2, 3)\par
6) (0, 6)\par
{\vskip 15pt} \par

\end{minipage}

               
{\begin{center}Завдання {9}\end{center}}
\vspace{-6mm}
Файл \code{'1.txt'} містить відомість з рядками, в яких через двокрапку записа\-но поряд\-ковий номер (ціле), назву предмету (знаків пунктуації не містить), прізвище студента, ім'я студента, номер заліковки, оцінку в балах за 100 бальною шкалою.

Білі символи навколо роздільників не допускаються.

Приклад рядка файлу



\begin{verbatim}
13:algebra:surname:name:123456:77
\end{verbatim}


Нижче наведено код, який має залишити у файлі в кожному рядку тільки номер заліковки, назву предмету та оцінку в балах за 100-бальною шкалою (у заданій послі\-дов\-ності).

Рядок з прикладу має бути замінений на такий



\begin{verbatim}
123456:algebra:77
\end{verbatim}




Код:



\begin{verbatim}
regex = re.compile(r'XXX')
with open('1.txt', encoding='utf8') as f:
    lst = f.readlines()
for i in range(len(lst)):
    lst[i] = re.sub(regex, YYY, lst[i])
with open('1.txt', 'w', encoding='utf8') as f:
    f.writelines(lst)
\end{verbatim}





Який рядок треба записати на місце \kw{XXX} ?
Який рядок треба записати на місце \kw{YYY} ?\par

%\end {large}
}
\newpage

{{23.05.2025 Контрольна робота, варіант  \No \textbf { 21 }} група К1{\parbox {0.5 cm} \dotfill}  ПІБ {\parbox {2.5 cm} \dotfill}\par}
{
%\begin {large}
\begin{minipage}[t]{9.7cm}
               
{\begin{center}Завдання {1}\end{center}}
\vspace{-6mm}
Укажіть ТИП значення виразу.\par\par
\begin{verbatim}
{x for x in range(7)}
\end{verbatim}

\vspace{15mm}
У завданнях 2, 3 вказати, що буде виведено за виконання. Повідомлення інтерпретатора про помилку позначати ERROR

               
{\begin{center}Завдання {2}\end{center}}
\vspace{-6mm}

\begin{verbatim}
try:
    s = {42:6, 87:8, 34:2, 42:5}
    s[34] = 2
    for x in s.items():
        print(*x, sep=' ', end=' ')
except:
    print('ERR')
\end{verbatim}


\end{minipage}
\vline \hspace{8mm}
\begin{minipage}[t]{8.0cm}                
               
{\begin{center}Завдання {3}\end{center}}
\vspace{-6mm}

\begin{verbatim}
b = 0
j = 1

class D:
    j = 2
    
    def __init__(self):
        self.j = 3
        j = 4
        D.j = 5

obj = D()
try:
    print(b, end=' ')
    print(j, end=' ')
    print(D.j, end=' ')
    print(obj.j)
except:
    print('ERR')
\end{verbatim}

\end{minipage}

               
{\begin{center}Завдання {4}\end{center}}
\vspace{-6mm}
Нехай змінна \kw{s} позначає послідовність цілих, по якій можна ітерувати.

Написати вираз, що перевіряє, чи правда, що всі елементи послідовності є цілими.\par

\vspace{10mm}
               
{\begin{center}Завдання {5}\end{center}}
\vspace{-6mm}
Змінні \kw{next} та \kw{prev} вузла списку позначають наступний та попередній за ним вузол відповідно. Починаючи з голови, у список записано послідовні цілі числа від~$-45$ до~$31$. Нехай змінна \kw{p} позначає вузол, що зберігає значення~$-1$.
Яке число зберігається у вузлі \kw{p.next.next.next.next.prev} ?\par

               
{\begin{center}Завдання {6}\end{center}}
\vspace{-6mm}
Значенням змінної \kw{a} є цілочисловий масив типу $numpy.ndarray$ розмірів $7{\times}6$. На скільки збільшиться сума елементів масиву \kw{a} після виконання
\begin{verbatim}
a.T[0] += 2
\end{verbatim}


Якщо код некоректний, то в якості відповіді зазначити ERROR.\par

               
{\begin{center}Завдання {10}\end{center}}
\vspace{-6mm}

Написати функцію $f$, яка для файлу, заданого шляхом, розв’язує задачу:\\ Текстовий файл містить послідовність чисел $a_1$, $a_2$, $\ldots$, $a_t$, записану у форматі одне число на рядок. Знайти третє найбільше значення, що записане у файлі. Кожне значення, що записане у файлі, враховується тільки один раз. Для послідовності 1, 4, 12, 5, 1, 0 таким значенням буде 4.

В усіх випадках розміри файлів такі, що їх вміст повністю завантажений у пам'ять бути не може. Якість коду також оцінюється.



\vspace{10mm}\par

\newpage
\begin{minipage}[t]{9.7cm}
               
{\begin{center}Завдання {7}\end{center}}
\vspace{-6mm}

\begin{center}
	\adjustimage{max size={0.8\linewidth}{0.8\paperheight}}
	{./15BinTree/trees_exp/50.png}
\end{center}
Указати послідовність міток вершин дерева, зображеного на рис., що утвориться в результаті обходу дерева в оберненому порядку. Мітки записувати через пробіл.\par Алгоритм починає роботу з кореня (на рис. це верхня вершина).\par

\end{minipage}
\vline \hspace{8mm}
\begin{minipage}[t]{8.0cm}
               
{\begin{center}Завдання {8}\end{center}}
\vspace{-6mm}
Для графа на рисунку з вершини 4 запущено алгоритм пошуку в глибину, що будує кістякове дерево. Списки суміжності вершин переглядаються алгоритмом за зростанням міток вершин.\par
\begin{center}
	\adjustimage{max size={0.9\linewidth}{0.9\paperheight}}
	{./16Graphs/Traversals/116.png}
\end{center}
\par Які з наведених ребер входять до складу отриманого кістякового дерева?\par
1) (2, 4)\par
2) (2, 3)\par
3) (1, 6)\par
4) (0, 3)\par
5) (0, 2)\par
6) (3, 6)\par
{\vskip 15pt} \par

\end{minipage}

               
{\begin{center}Завдання {9}\end{center}}
\vspace{-6mm}
Файл \code{'1.txt'} містить відомість з рядками, в яких через кому записано поряд\-ковий номер (ціле), назву предмету (знаків пунктуації не містить), прізвище студента, ім'я студента, номер заліковки, оцінку в балах за 100 бальною шкалою.

Білі символи навколо роздільників не допускаються.

Приклад рядка файлу



\begin{verbatim}
13,algebra,surname,name,123456,77
\end{verbatim}


Нижче наведено код, який шукає усі предмети, який складав студент із заліковкою \code{777}, та зберігає їх типи в змінній \kw{codes}.



\begin{verbatim}
codes = set()
regex = re.compile(r'XXX')
with open('1.txt', encoding='utf8') as f:
    for s in f:
        res = re.fullmatch(regex, s)
        if not res: continue
        s = ''
        codes.add(YYY)
\end{verbatim}



Який рядок треба записати на місце \kw{XXX} ?
Який рядок треба записати на місце \kw{YYY} ?\par

%\end {large}
}
\newpage

{{23.05.2025 Контрольна робота, варіант  \No \textbf { 22 }} група К1{\parbox {0.5 cm} \dotfill}  ПІБ {\parbox {2.5 cm} \dotfill}\par}
{
%\begin {large}
\begin{minipage}[t]{9.7cm}
               
{\begin{center}Завдання {1}\end{center}}
\vspace{-6mm}
Укажіть ТИП значення виразу.\par\par
\begin{verbatim}
[13,2,4]+[3,5]
\end{verbatim}

\vspace{15mm}
У завданнях 2, 3 вказати, що буде виведено за виконання. Повідомлення інтерпретатора про помилку позначати ERROR

               
{\begin{center}Завдання {2}\end{center}}
\vspace{-6mm}

\begin{verbatim}
try:
    s = {76:4, 56:0, 67:3, 56:0}
    s[56] = 7
    for x in s.items():
        print(x, sep=' ', end=' ')
except:
    print('ERR')
\end{verbatim}


\end{minipage}
\vline \hspace{8mm}
\begin{minipage}[t]{8.0cm}                
               
{\begin{center}Завдання {3}\end{center}}
\vspace{-6mm}

\begin{verbatim}
a = 0

class A:
    a = 1
    
    def __init__(self):
        self.a = 2
        a = 3
        A.a = 4

obj = A()
obj.a = 5
try:
    print(a, end=' ')
    print(obj.a, end=' ')
    print(A.a, end=' ')
    print(obj.a, end=' ')
    print(obj._A__a)
except:
    print('ERR')
\end{verbatim}

\end{minipage}

               
{\begin{center}Завдання {4}\end{center}}
\vspace{-6mm}
Нехай змінна \kw{s} позначає дуже довгу послідовність, по якій можна ітерувати.

Написати вираз-генератор, що задає підпослідовність її елементів, що є рядками типу \kw{str}.

Обмеження: усі елементи шуканої підпослідовності одночасно в пам'яті розміщені бути не можуть.\par

\vspace{10mm}
               
{\begin{center}Завдання {5}\end{center}}
\vspace{-6mm}
Змінні \kw{next} та \kw{prev} вузла списку позначають наступний та попередній за ним вузол відповідно. Починаючи з голови, у список записано послідовні цілі числа від~$-41$ до~$15$. Нехай змінна \kw{p} позначає вузол, що зберігає значення~$-7$.
Яке число зберігається у вузлі \kw{p.prev.prev.prev.prev.next} ?\par

               
{\begin{center}Завдання {6}\end{center}}
\vspace{-6mm}
Значенням змінної \kw{a} є цілочисловий масив типу $numpy.ndarray$ розмірів $4{\times}5$. На скільки збільшиться сума елементів масиву \kw{a} після виконання
\begin{verbatim}
a[2] += 1
\end{verbatim}


Якщо код некоректний, то в якості відповіді зазначити ERROR.\par

               
{\begin{center}Завдання {10}\end{center}}
\vspace{-6mm}

Написати функцію $f$, яка для файлу, заданого шляхом, розв’язує задачу:\\ Кожен з двох текстових файлів містить послідовність чисел, записану у форматі одне число на рядок.  Перевірити, чи задають вони одну ту саму послідовність.

В усіх випадках розміри файлів такі, що їх вміст повністю завантажений у пам'ять бути не може. Якість коду також оцінюється.



\vspace{10mm}\par

\newpage
\begin{minipage}[t]{9.7cm}
               
{\begin{center}Завдання {7}\end{center}}
\vspace{-6mm}

\begin{center}
	\adjustimage{max size={0.8\linewidth}{0.8\paperheight}}
	{./15BinTree/trees_exp/32.png}
\end{center}
Указати послідовність міток вершин дерева, зображеного на рис., що утвориться в результаті обходу дерева в прямому порядку. Мітки записувати через пробіл.\par Алгоритм починає роботу з кореня (на рис. це верхня вершина).\par

\end{minipage}
\vline \hspace{8mm}
\begin{minipage}[t]{8.0cm}
               
{\begin{center}Завдання {8}\end{center}}
\vspace{-6mm}
Для графа на рисунку з вершини 2 запущено алгоритм пошуку в ширину, що будує кістякове дерево. Списки суміжності вершин переглядаються алгоритмом за зростанням міток вершин.\par
\begin{center}
	\adjustimage{max size={0.9\linewidth}{0.9\paperheight}}
	{./16Graphs/Traversals/16.png}
\end{center}
\par Які з наведених ребер входять до складу отриманого кістякового дерева?\par
1) (1, 3)\par
2) (3, 5)\par
3) (1, 6)\par
4) (0, 1)\par
5) (2, 5)\par
6) (0, 6)\par
{\vskip 15pt} \par

\end{minipage}

               
{\begin{center}Завдання {9}\end{center}}
\vspace{-6mm}
Файл \code{'1.txt'} містить відомість з рядками, в яких через двокрапку записа\-но назву предмету (знаків пунктуації не містить), прізвище студента, ім'я студента, номер залі\-ков\-ки, оцінку в балах за 100 бальною шкалою.

Білі символи навколо роздільників не допускаються.

Приклад рядка файлу



\begin{verbatim}
algebra:surname:name:123456:77
\end{verbatim}


Нижче наведено код, який має замінити у файлі усі номери заліковок на рядок \code{***}.



\begin{verbatim}
regex = re.compile(r'XXX')
with open('1.txt', encoding='utf8') as f:
    lst = f.readlines()
for i in range(len(lst)):
    lst[i] = re.sub(regex, YYY, lst[i])
with open('1.txt', 'w', encoding='utf8') as f:
    f.writelines(lst)
\end{verbatim}





Який рядок треба записати на місце \kw{XXX} ?
Який рядок треба записати на місце \kw{YYY} ?\par

%\end {large}
}
\newpage

{{23.05.2025 Контрольна робота, варіант  \No \textbf { 23 }} група К1{\parbox {0.5 cm} \dotfill}  ПІБ {\parbox {2.5 cm} \dotfill}\par}
{
%\begin {large}
\begin{minipage}[t]{9.7cm}
               
{\begin{center}Завдання {1}\end{center}}
\vspace{-6mm}
Укажіть ТИП значення виразу.\par\par
\begin{verbatim}
{(1,4), 8}
\end{verbatim}

\vspace{15mm}
У завданнях 2, 3 вказати, що буде виведено за виконання. Повідомлення інтерпретатора про помилку позначати ERROR

               
{\begin{center}Завдання {2}\end{center}}
\vspace{-6mm}

\begin{verbatim}
try:
    d = {36:2, 85:1, 68:1, 85:6}
    d[36] = 2
    for x in d.keys():
        print(x, sep=' ', end=' ')
except:
    print('ERR')
\end{verbatim}


\end{minipage}
\vline \hspace{8mm}
\begin{minipage}[t]{8.0cm}                
               
{\begin{center}Завдання {3}\end{center}}
\vspace{-6mm}

\begin{verbatim}
__c = 0

class D:
    __c = 1
    
    def __init__(self):
        self.__c = 2
        __c = 3
        D.__c = 4

obj = D()
D.__c = 5
try:
    print(__c, end=' ')
    print(obj.__c, end=' ')
    print(D.__c, end=' ')
    print(obj.__c, end=' ')
    print(obj._D__c)
except:
    print('ERR')
\end{verbatim}

\end{minipage}

               
{\begin{center}Завдання {4}\end{center}}
\vspace{-6mm}
Нехай змінна \kw{s} позначає послідовність цілих, по якій можна ітерувати.

Написати вираз, що перевіряє, чи правда, що послідовність  містить хоча б одне додатнє число.\par

\vspace{10mm}
               
{\begin{center}Завдання {5}\end{center}}
\vspace{-6mm}
Змінні \kw{next} та \kw{prev} вузла списку позначають наступний та попередній за ним вузол відповідно. Починаючи з голови, у список записано послідовні цілі числа від~$-55$ до~$50$. Нехай змінна \kw{p} позначає вузол, що зберігає значення~$2$.
Яке число зберігається у вузлі \kw{p.prev.prev.prev.next.prev} ?\par

               
{\begin{center}Завдання {6}\end{center}}
\vspace{-6mm}
Значенням змінної \kw{a} є цілочисловий масив типу $numpy.ndarray$ розмірів $6{\times}4$. На скільки збільшиться сума елементів масиву \kw{a} після виконання
\begin{verbatim}
a.T[-2] += 3
\end{verbatim}


Якщо код некоректний, то в якості відповіді зазначити ERROR.\par

               
{\begin{center}Завдання {10}\end{center}}
\vspace{-6mm}

Написати функцію $f$, яка для файлу, заданого шляхом, розв’язує задачу:\\ Текстовий файл містить послідовність чисел $a_1$, $a_2$, $\ldots$, $a_t$, записану у форматі одне число на рядок. Знайти третє найменше значення, що записане у файлі. Кожне значення, що записане у файлі, враховується стільки разів, скільки входить у файл. Для послідовності 2, 4, 12, 5, 2, 0 таким значенням буде 2.

В усіх випадках розміри файлів такі, що їх вміст повністю завантажений у пам'ять бути не може. Якість коду також оцінюється.



\vspace{10mm}\par

\newpage
\begin{minipage}[t]{9.7cm}
               
{\begin{center}Завдання {7}\end{center}}
\vspace{-6mm}

\begin{center}
	\adjustimage{max size={0.8\linewidth}{0.8\paperheight}}
	{./15BinTree/trees_exp/47.png}
\end{center}
Указати послідовність міток вершин дерева, зображеного на рис., що утвориться в результаті обходу дерева в оберненому порядку. Мітки записувати через пробіл.\par Алгоритм починає роботу з кореня (на рис. це верхня вершина).\par

\end{minipage}
\vline \hspace{8mm}
\begin{minipage}[t]{8.0cm}
               
{\begin{center}Завдання {8}\end{center}}
\vspace{-6mm}
Для графа на рисунку з вершини 1 запущено алгоритм пошуку в глибину, що будує кістякове дерево. Списки суміжності вершин переглядаються алгоритмом за зростанням міток вершин.\par
\begin{center}
	\adjustimage{max size={0.9\linewidth}{0.9\paperheight}}
	{./16Graphs/Traversals/46.png}
\end{center}
\par Які з наведених ребер входять до складу отриманого кістякового дерева?\par
1) (1, 3)\par
2) (1, 5)\par
3) (4, 6)\par
4) (3, 4)\par
5) (0, 3)\par
6) (0, 4)\par
{\vskip 15pt} \par

\end{minipage}

               
{\begin{center}Завдання {9}\end{center}}
\vspace{-6mm}
Файл \code{'1.txt'} містить журнал з рядками, в яких через кому записано поряд\-ко\-вий номер запису в журналі (ціле), тип події (ERROR, WARNING, INFO, STATUS), ім'я користувача (складається з латиниці), час події (фор\-мат як у прикладі), код події (ціле).

Білі символи навколо роздільників не допускаються.

Приклад рядка файлу



\begin{verbatim}
13,ERROR,username,2022-05-05 13:26,404
\end{verbatim}


Нижче наведено код, який має замінити у файлі імена користувачів на рядок \code{unknown}.



\begin{verbatim}
regex = re.compile(r'XXX')
with open('1.txt', encoding='utf8') as f:
    lst = f.readlines()
for i in range(len(lst)):
    lst[i] = re.sub(regex, YYY, lst[i])
with open('1.txt', 'w', encoding='utf8') as f:
    f.writelines(lst)
\end{verbatim}





Який рядок треба записати на місце \kw{XXX} ?
Який рядок треба записати на місце \kw{YYY} ?\par

%\end {large}
}
\newpage

{{23.05.2025 Контрольна робота, варіант  \No \textbf { 24 }} група К1{\parbox {0.5 cm} \dotfill}  ПІБ {\parbox {2.5 cm} \dotfill}\par}
{
%\begin {large}
\begin{minipage}[t]{9.7cm}
               
{\begin{center}Завдання {1}\end{center}}
\vspace{-6mm}
Укажіть ТИП значення виразу.\par\par
\begin{verbatim}
{21,2,4}
\end{verbatim}

\vspace{15mm}
У завданнях 2, 3 вказати, що буде виведено за виконання. Повідомлення інтерпретатора про помилку позначати ERROR

               
{\begin{center}Завдання {2}\end{center}}
\vspace{-6mm}

\begin{verbatim}
try:
    t = {30:1, 71:2, 10:0, 10:6}
    t[37] = 4
    for x in t :
        print(x, sep=' ', end=' ')
except:
    print('ERR')
\end{verbatim}


\end{minipage}
\vline \hspace{8mm}
\begin{minipage}[t]{8.0cm}                
               
{\begin{center}Завдання {3}\end{center}}
\vspace{-6mm}

\begin{verbatim}
c = 0
k = 1

class A:
    c = 2
    
    def __init__(self):
        global c
        self.c = 3
        c = 4
        A.c = 5

obj = A()
try:
    print(c, end=' ')
    print(k, end=' ')
    print(A.c, end=' ')
    print(obj.c)
except:
    print('ERR')
\end{verbatim}

\end{minipage}

               
{\begin{center}Завдання {4}\end{center}}
\vspace{-6mm}
Нехай змінна \kw{s} позначає послідовність цілих, по якій можна ітерувати.

Написати вираз, що перевіряє, чи правда, що послідовність  містить хоча б одне число, що ділиться на 7.\par

\vspace{10mm}
               
{\begin{center}Завдання {5}\end{center}}
\vspace{-6mm}
Змінні \kw{next} та \kw{prev} вузла списку позначають наступний та попередній за ним вузол відповідно. Починаючи з голови, у список записано послідовні цілі числа від~$-21$ до~$18$. Нехай змінна \kw{p} позначає вузол, що зберігає значення~$-4$.
Яке число зберігається у вузлі \kw{p.next.next.next.prev.next} ?\par

               
{\begin{center}Завдання {6}\end{center}}
\vspace{-6mm}
Значенням змінної \kw{a} є цілочисловий масив типу $numpy.ndarray$ розмірів $3{\times}3$. На скільки збільшиться сума елементів масиву \kw{a} після виконання
\begin{verbatim}
a[1] += 2
\end{verbatim}


Якщо код некоректний, то в якості відповіді зазначити ERROR.\par

               
{\begin{center}Завдання {10}\end{center}}
\vspace{-6mm}

Написати функцію $f$, яка для файлу, заданого шляхом, розв’язує задачу:\\ Текстовий файл містить послідовність чисел $a_1$, $a_2$, $\ldots$, $a_t$, записану у форматі одне число на рядок. Знайти третє найбільше значення, що записане у файлі. Кожне значення, що записане у файлі, враховується стільки разів, скільки входить у файл. Для послідовності 2, 4, 12, 5, 2, 0 таким значенням буде 2.

В усіх випадках розміри файлів такі, що їх вміст повністю завантажений у пам'ять бути не може. Якість коду також оцінюється.



\vspace{10mm}\par

\newpage
\begin{minipage}[t]{9.7cm}
               
{\begin{center}Завдання {7}\end{center}}
\vspace{-6mm}

\begin{center}
	\adjustimage{max size={0.8\linewidth}{0.8\paperheight}}
	{./15BinTree/trees_exp/26.png}
\end{center}
Указати послідовність міток вершин дерева, зображеного на рис., що утвориться в результаті обходу дерева в прямому порядку. Мітки записувати через пробіл.\par Алгоритм починає роботу з кореня (на рис. це верхня вершина).\par

\end{minipage}
\vline \hspace{8mm}
\begin{minipage}[t]{8.0cm}
               
{\begin{center}Завдання {8}\end{center}}
\vspace{-6mm}
Для графа на рисунку з вершини 3 запущено алгоритм пошуку в ширину, що будує кістякове дерево. Списки суміжності вершин переглядаються алгоритмом за зростанням міток вершин.\par
\begin{center}
	\adjustimage{max size={0.9\linewidth}{0.9\paperheight}}
	{./16Graphs/Traversals/255.png}
\end{center}
\par Які з наведених ребер входять до складу отриманого кістякового дерева?\par
1) (5, 6)\par
2) (0, 4)\par
3) (4, 6)\par
4) (2, 3)\par
5) (2, 4)\par
6) (1, 3)\par
{\vskip 15pt} \par

\end{minipage}

               
{\begin{center}Завдання {9}\end{center}}
\vspace{-6mm}
Рядки відомості через кому містять порядковий номер (ціле), пріз\-ви\-ще студен\-та, ім'я студента, номер заліковки, оцінку в балах за 100 бальною шкалою.

Білі символи навколо роздільників не допускаються.

Приклад такого рядка присвоєно змінній \kw{s}.



\begin{verbatim}
s = '13,surname,name,123456,77'
\end{verbatim}


Нижче наведено код, за виконання якого для рядка \kw{s} значенням змінної \kw{out} має стати номер заліковки.

Для наведеного прикладу значенням змінної \kw{out} має стати рядок \code{123456} .



\begin{verbatim}
res = re.fullmatch(r'XXX', s)
s = ''
out = YYY
\end{verbatim}


Код має працювати не тільки для конкретного рядка з прикладу, але й для кожного рядка з відомості.


Який рядок треба записати на місце \kw{XXX} ?
Який рядок треба записати на місце \kw{YYY} ?\par

%\end {large}
}
\newpage

{{23.05.2025 Контрольна робота, варіант  \No \textbf { 25 }} група К1{\parbox {0.5 cm} \dotfill}  ПІБ {\parbox {2.5 cm} \dotfill}\par}
{
%\begin {large}
\begin{minipage}[t]{9.7cm}
               
{\begin{center}Завдання {1}\end{center}}
\vspace{-6mm}
Укажіть ТИП значення виразу.\par\par
\begin{verbatim}
{y:13 for y in range(-7,7)}
\end{verbatim}

\vspace{15mm}
У завданнях 2, 3 вказати, що буде виведено за виконання. Повідомлення інтерпретатора про помилку позначати ERROR

               
{\begin{center}Завдання {2}\end{center}}
\vspace{-6mm}

\begin{verbatim}
try:
    d = {46:3, 48:3, 74:4, 79:8, 48:6}
    d[79] = 3
    for x in d :
        print(x, sep=' ', end=' ')
except:
    print('ERR')
\end{verbatim}


\end{minipage}
\vline \hspace{8mm}
\begin{minipage}[t]{8.0cm}                
               
{\begin{center}Завдання {3}\end{center}}
\vspace{-6mm}

\begin{verbatim}
d = 0
k = 1

class C:
    d = 2
    
    def __init__(self):
        self.k = 3
        d = 4
        C.d = 5

obj = C()
try:
    print(d, end=' ')
    print(k, end=' ')
    print(C.d, end=' ')
    print(obj.k)
except:
    print('ERR')
\end{verbatim}

\end{minipage}

               
{\begin{center}Завдання {4}\end{center}}
\vspace{-6mm}
Нехай змінна \kw{s} позначає послідовність цілих, по якій можна ітерувати.

Написати вираз, що перевіряє, чи правда, що послідовність  містить числа, менші 13.\par

\vspace{10mm}
               
{\begin{center}Завдання {5}\end{center}}
\vspace{-6mm}
Змінні \kw{next} та \kw{prev} вузла списку позначають наступний та попередній за ним вузол відповідно. Починаючи з голови, у список записано послідовні цілі числа від~$-53$ до~$17$. Нехай змінна \kw{p} позначає вузол, що зберігає значення~$-6$.
Яке число зберігається у вузлі \kw{p.next.prev.prev.prev.prev} ?\par

               
{\begin{center}Завдання {6}\end{center}}
\vspace{-6mm}
Значенням змінної \kw{a} є цілочисловий масив типу $numpy.ndarray$ розмірів $5{\times}7$. На скільки збільшиться сума елементів масиву \kw{a} після виконання
\begin{verbatim}
a.T[-1] += 2
\end{verbatim}


Якщо код некоректний, то в якості відповіді зазначити ERROR.\par

               
{\begin{center}Завдання {10}\end{center}}
\vspace{-6mm}

Написати функцію $f$, яка для файлу, заданого шляхом, розв’язує задачу:\\ Текстовий файл містить послідовність чисел $a_1$, $a_2$, $\ldots$, $a_t$, записану у форматі одне число на рядок. Знайти третє найбільше значення, що записане у файлі. Кожне значення, що записане у файлі, враховується тільки один раз. Для послідовності 1, 4, 12, 5, 1, 0 таким значенням буде 4.

В усіх випадках розміри файлів такі, що їх вміст повністю завантажений у пам'ять бути не може. Якість коду також оцінюється.



\vspace{10mm}\par

\newpage
\begin{minipage}[t]{9.7cm}
               
{\begin{center}Завдання {7}\end{center}}
\vspace{-6mm}

\begin{center}
	\adjustimage{max size={0.8\linewidth}{0.8\paperheight}}
	{./15BinTree/trees_exp/164.png}
\end{center}
Указати послідовність міток вершин дерева, зображеного на рис., що утвориться в результаті обходу дерева в оберненому порядку. Мітки записувати через пробіл.\par Алгоритм починає роботу з кореня (на рис. це верхня вершина).\par

\end{minipage}
\vline \hspace{8mm}
\begin{minipage}[t]{8.0cm}
               
{\begin{center}Завдання {8}\end{center}}
\vspace{-6mm}
Для графа на рисунку з вершини 5 запущено алгоритм пошуку в глибину, що будує кістякове дерево. Списки суміжності вершин переглядаються алгоритмом за зростанням міток вершин.\par
\begin{center}
	\adjustimage{max size={0.9\linewidth}{0.9\paperheight}}
	{./16Graphs/Traversals/146.png}
\end{center}
\par Які з наведених ребер входять до складу отриманого кістякового дерева?\par
1) (1, 3)\par
2) (0, 5)\par
3) (1, 6)\par
4) (2, 5)\par
5) (2, 6)\par
6) (1, 4)\par
{\vskip 15pt} \par

\end{minipage}

               
{\begin{center}Завдання {9}\end{center}}
\vspace{-6mm}
Файл \code{'1.txt'} містить журнал з рядками, в яких через кому записано поряд\-ко\-вий номер запису в журналі (ціле), тип події (ERROR, WARNING, INFO, STATUS), ім'я користувача (складається з латиниці), час події (фор\-мат як у прикладі), опис події.

Білі символи навколо роздільників не допускаються.

Приклад рядка файлу



\begin{verbatim}
13,ERROR,username,2022-05-05 13:26,something happens
\end{verbatim}


Нижче наведено код, який має залишити у файлі в кожному рядку тільки тип події та її опис.

Рядок з прикладу має бути замінений на такий



\begin{verbatim}
ERROR,something happens
\end{verbatim}




Код:



\begin{verbatim}
regex = re.compile(r'XXX')
with open('1.txt', encoding='utf8') as f:
    lst = f.readlines()
for i in range(len(lst)):
    lst[i] = re.sub(regex, YYY, lst[i])
with open('1.txt', 'w', encoding='utf8') as f:
    f.writelines(lst)
\end{verbatim}







Який рядок треба записати на місце \kw{XXX} ?
Який рядок треба записати на місце \kw{YYY} ?\par

%\end {large}
}
\newpage

{{23.05.2025 Контрольна робота, варіант  \No \textbf { 26 }} група К1{\parbox {0.5 cm} \dotfill}  ПІБ {\parbox {2.5 cm} \dotfill}\par}
{
%\begin {large}
\begin{minipage}[t]{9.7cm}
               
{\begin{center}Завдання {1}\end{center}}
\vspace{-6mm}
Укажіть ТИП значення виразу.\par\par
\begin{verbatim}
{x:0 for x in range(7)}
\end{verbatim}

\vspace{15mm}
У завданнях 2, 3 вказати, що буде виведено за виконання. Повідомлення інтерпретатора про помилку позначати ERROR

               
{\begin{center}Завдання {2}\end{center}}
\vspace{-6mm}

\begin{verbatim}
try:
    s = {28:3, 34:5, 37:9, 37:2}
    s[66] = 1
    for x, y in s.items():
        print(x, y, sep=' ', end=' ')
except:
    print('ERR')
\end{verbatim}


\end{minipage}
\vline \hspace{8mm}
\begin{minipage}[t]{8.0cm}                
               
{\begin{center}Завдання {3}\end{center}}
\vspace{-6mm}

\begin{verbatim}
a = 0
h = 1

class B:
    h = 2
    
    def __init__(self):
        global h
        self.a = 3
        h = 4
        B.h = 5

obj = B()
try:
    print(a, end=' ')
    print(h, end=' ')
    print(B.h, end=' ')
    print(obj.a)
except:
    print('ERR')
\end{verbatim}

\end{minipage}

               
{\begin{center}Завдання {4}\end{center}}
\vspace{-6mm}
Нехай змінна \kw{s} позначає дуже довгу послідовність цілих, по якій можна ітерувати.

Написати вираз-генератор, що задає підпослідовність її елементів, що не діляться на жодне з чисел 2, 7, 5.

Обмеження: усі елементи шуканої підпослідовності одночасно в пам'яті розміщені бути не можуть.\par

\vspace{10mm}
               
{\begin{center}Завдання {5}\end{center}}
\vspace{-6mm}
Змінні \kw{next} та \kw{prev} вузла списку позначають наступний та попередній за ним вузол відповідно. Починаючи з голови, у список записано послідовні цілі числа від~$-24$ до~$32$. Нехай змінна \kw{p} позначає вузол, що зберігає значення~$-3$.
Яке число зберігається у вузлі \kw{p.prev.next.next.next.next} ?\par

               
{\begin{center}Завдання {6}\end{center}}
\vspace{-6mm}
Значенням змінної \kw{a} є цілочисловий масив типу $numpy.ndarray$ розмірів $4{\times}5$. На скільки збільшиться сума елементів масиву \kw{a} після виконання
\begin{verbatim}
a[-3] += 1
\end{verbatim}


Якщо код некоректний, то в якості відповіді зазначити ERROR.\par

               
{\begin{center}Завдання {10}\end{center}}
\vspace{-6mm}

Написати функцію $f$, яка для файлу, заданого шляхом, розв’язує задачу:\\ Текстовий файл містить послідовність чисел $a_1$, $a_2$, $\ldots$, $a_t$, записану у форматі одне число на рядок. Знайти третє найбільше значення, що записане у файлі. Кожне значення, що записане у файлі, враховується стільки разів, скільки входить у файл. Для послідовності 2, 4, 12, 5, 2, 0 таким значенням буде 2.

В усіх випадках розміри файлів такі, що їх вміст повністю завантажений у пам'ять бути не може. Якість коду також оцінюється.



\vspace{10mm}\par

\newpage
\begin{minipage}[t]{9.7cm}
               
{\begin{center}Завдання {7}\end{center}}
\vspace{-6mm}

\begin{center}
	\adjustimage{max size={0.8\linewidth}{0.8\paperheight}}
	{./15BinTree/trees_exp/137.png}
\end{center}
Указати послідовність міток вершин дерева, зображеного на рис., що утвориться в результаті обходу дерева в прямому порядку. Мітки записувати через пробіл.\par Алгоритм починає роботу з кореня (на рис. це верхня вершина).\par

\end{minipage}
\vline \hspace{8mm}
\begin{minipage}[t]{8.0cm}
               
{\begin{center}Завдання {8}\end{center}}
\vspace{-6mm}
Для графа на рисунку з вершини 3 запущено алгоритм пошуку в ширину, що будує кістякове дерево. Списки суміжності вершин переглядаються алгоритмом за зростанням міток вершин.\par
\begin{center}
	\adjustimage{max size={0.9\linewidth}{0.9\paperheight}}
	{./16Graphs/Traversals/45.png}
\end{center}
\par Які з наведених ребер входять до складу отриманого кістякового дерева?\par
1) (0, 4)\par
2) (1, 3)\par
3) (4, 6)\par
4) (2, 3)\par
5) (2, 4)\par
6) (0, 3)\par
{\vskip 15pt} \par

\end{minipage}

               
{\begin{center}Завдання {9}\end{center}}
\vspace{-6mm}
Файл \code{'1.txt'} містить журнал з рядками, в яких через двокрапку записано поряд\-ко\-вий номер запису (ціле), ім'я користувача (складається з латиниці), час події (фор\-мат як у прикладі), тип події (ERROR, WARNING, INFO, STATUS), код події (ціле).

Білі символи навколо роздільників не допускаються.

Приклад рядка журналу



\begin{verbatim}
12:username:2022-05-05 13-26:ERROR:404
\end{verbatim}


Нижче наведено код, який шукає усі попередження (подія WARNING), що зустрі\-ча\-ються в журналі і відбулися під час обідньої перерви деканату (з 13:00 по 13:59 включно), та зберігає їх коди в змінній \kw{codes}.



\begin{verbatim}
codes = set()
regex = re.compile(r'XXX')
with open('1.txt', encoding='utf8') as f:
    for s in f:
        res = re.fullmatch(regex, s)
        if not res: continue
        s = ''
        codes.add(YYY)
\end{verbatim}





Який рядок треба записати на місце \kw{XXX} ?
Який рядок треба записати на місце \kw{YYY} ?\par

%\end {large}
}
\newpage

{{23.05.2025 Контрольна робота, варіант  \No \textbf { 27 }} група К1{\parbox {0.5 cm} \dotfill}  ПІБ {\parbox {2.5 cm} \dotfill}\par}
{
%\begin {large}
\begin{minipage}[t]{9.7cm}
               
{\begin{center}Завдання {1}\end{center}}
\vspace{-6mm}
Укажіть ТИП значення виразу.\par\par
\begin{verbatim}
(11,2,3)
\end{verbatim}

\vspace{15mm}
У завданнях 2, 3 вказати, що буде виведено за виконання. Повідомлення інтерпретатора про помилку позначати ERROR

               
{\begin{center}Завдання {2}\end{center}}
\vspace{-6mm}

\begin{verbatim}
try:
    t = {77:0, 65:5, 30:1, 50:8}
    t[77] = 8
    for x in t.items():
        print(*x, sep=' ', end=' ')
except:
    print('ERR')
\end{verbatim}


\end{minipage}
\vline \hspace{8mm}
\begin{minipage}[t]{8.0cm}                
               
{\begin{center}Завдання {3}\end{center}}
\vspace{-6mm}

\begin{verbatim}
_b = 0

class B:
    _b = 1
    
    def __init__(self):
        self._b = 2
        _b = 3
        B._b = 4

obj = B()
B._b = 5
try:
    print(_b, end=' ')
    print(obj._b, end=' ')
    print(B._b, end=' ')
    print(obj._b, end=' ')
    print(obj._B__b)
except:
    print('ERR')
\end{verbatim}

\end{minipage}

               
{\begin{center}Завдання {4}\end{center}}
\vspace{-6mm}
Нехай змінна \kw{s} позначає послідовність цілих, по якій можна ітерувати.

Написати вираз, що перевіряє, чи правда, що послідовність  містить від’ємні числа.\par

\vspace{10mm}
               
{\begin{center}Завдання {5}\end{center}}
\vspace{-6mm}
Змінні \kw{next} та \kw{prev} вузла списку позначають наступний та попередній за ним вузол відповідно. Починаючи з голови, у список записано послідовні цілі числа від~$-36$ до~$26$. Нехай змінна \kw{p} позначає вузол, що зберігає значення~$-2$.
Яке число зберігається у вузлі \kw{p.prev.next.prev.next.prev} ?\par

               
{\begin{center}Завдання {6}\end{center}}
\vspace{-6mm}
Значенням змінної \kw{a} є цілочисловий масив типу $numpy.ndarray$ розмірів $5{\times}6$. На скільки збільшиться сума елементів масиву \kw{a} після виконання
\begin{verbatim}
a.T[-1] += 3
\end{verbatim}


Якщо код некоректний, то в якості відповіді зазначити ERROR.\par

               
{\begin{center}Завдання {10}\end{center}}
\vspace{-6mm}

Написати функцію $f$, яка для файлу, заданого шляхом, розв’язує задачу:\\ Текстовий файл містить послідовність чисел $a_1$, $a_2$, $\ldots$, $a_t$, записану у форматі одне число на рядок. Знайти третє найменше значення, що записане у файлі. Кожне значення, що записане у файлі, враховується тільки один раз. \par
Для послідовності 1, 4, 12, 5, 1, 0 таким значенням буде 4.

В усіх випадках розміри файлів такі, що їх вміст повністю завантажений у пам'ять бути не може. Якість коду також оцінюється.



\vspace{10mm}\par

\newpage
\begin{minipage}[t]{9.7cm}
               
{\begin{center}Завдання {7}\end{center}}
\vspace{-6mm}

\begin{center}
	\adjustimage{max size={0.8\linewidth}{0.8\paperheight}}
	{./15BinTree/trees_exp/80.png}
\end{center}
Указати послідовність міток вершин дерева, зображеного на рис., що утвориться в результаті обходу дерева в оберненому порядку. Мітки записувати через пробіл.\par Алгоритм починає роботу з кореня (на рис. це верхня вершина).\par

\end{minipage}
\vline \hspace{8mm}
\begin{minipage}[t]{8.0cm}
               
{\begin{center}Завдання {8}\end{center}}
\vspace{-6mm}
Для графа на рисунку з вершини 0 запущено алгоритм пошуку в ширину, що будує кістякове дерево. Списки суміжності вершин переглядаються алгоритмом за зростанням міток вершин.\par
\begin{center}
	\adjustimage{max size={0.9\linewidth}{0.9\paperheight}}
	{./16Graphs/Traversals/293.png}
\end{center}
\par Які з наведених ребер входять до складу отриманого кістякового дерева?\par
1) (0, 4)\par
2) (0, 2)\par
3) (1, 6)\par
4) (3, 6)\par
5) (2, 6)\par
6) (1, 3)\par
{\vskip 15pt} \par

\end{minipage}

               
{\begin{center}Завдання {9}\end{center}}
\vspace{-6mm}
Файл \code{'1.txt'} містить відомість з рядками, в яких через кому записано поряд\-ковий номер (ціле), назву предмету (знаків пунктуації не містить), прізвище студента, ім'я студента, номер заліковки, оцінку в балах за 100 бальною шкалою.

Білі символи навколо роздільників не допускаються.

Приклад рядка файлу



\begin{verbatim}
13,algebra,surname,name,123456,77
\end{verbatim}


Нижче наведено код, який шукає усі назви предметів, що зустрічаються у файлі, та зберігає їх у змінній \id{codes}.



\begin{verbatim}
codes = set()
regex = re.compile(r'XXX')
with open('1.txt', encoding='utf8') as f:
    for s in f:
        res = re.fullmatch(regex, s)
        if not res: continue
        s = ''
        codes.add(YYY)
\end{verbatim}



Який рядок треба записати на місце \kw{XXX} ?
Який рядок треба записати на місце \kw{YYY} ?\par

%\end {large}
}
\newpage

{{23.05.2025 Контрольна робота, варіант  \No \textbf { 28 }} група К1{\parbox {0.5 cm} \dotfill}  ПІБ {\parbox {2.5 cm} \dotfill}\par}
{
%\begin {large}
\begin{minipage}[t]{9.7cm}
               
{\begin{center}Завдання {1}\end{center}}
\vspace{-6mm}
Укажіть ТИП значення виразу.\par\par
\begin{verbatim}
{x*x:x for x in range(-7,7)}
\end{verbatim}

\vspace{15mm}
У завданнях 2, 3 вказати, що буде виведено за виконання. Повідомлення інтерпретатора про помилку позначати ERROR

               
{\begin{center}Завдання {2}\end{center}}
\vspace{-6mm}

\begin{verbatim}
try:
    s = {34:2, 21:6, 67:7, 68:2}
    s[24] = 5
    for x in s.items():
        print(x, sep=' ', end=' ')
except:
    print('ERR')
\end{verbatim}


\end{minipage}
\vline \hspace{8mm}
\begin{minipage}[t]{8.0cm}                
               
{\begin{center}Завдання {3}\end{center}}
\vspace{-6mm}

\begin{verbatim}
a = 0
n = 1

class A:
    a = 2
    
    def __init__(self):
        global a
        self.n = 3
        a = 4
        A.a = 5

obj = A()
try:
    print(a, end=' ')
    print(n, end=' ')
    print(A.a, end=' ')
    print(obj.n)
except:
    print('ERR')
\end{verbatim}

\end{minipage}

               
{\begin{center}Завдання {4}\end{center}}
\vspace{-6mm}
Нехай змінна \kw{s} позначає дуже довгу послідовність цілих, по якій можна ітерувати.

Написати вираз-генератор, що задає підпослідовність її елементів, що є точними квадратами.

Обмеження: усі елементи шуканої підпослідовності одночасно в пам'яті розміщені бути не можуть.\par

\vspace{10mm}
               
{\begin{center}Завдання {5}\end{center}}
\vspace{-6mm}
Змінні \kw{next} та \kw{prev} вузла списку позначають наступний та попередній за ним вузол відповідно. Починаючи з голови, у список записано послідовні цілі числа від~$-49$ до~$22$. Нехай змінна \kw{p} позначає вузол, що зберігає значення~$-8$.
Яке число зберігається у вузлі \kw{p.next.prev.next.prev.next} ?\par

               
{\begin{center}Завдання {6}\end{center}}
\vspace{-6mm}
Значенням змінної \kw{a} є цілочисловий масив типу $numpy.ndarray$ розмірів $3{\times}4$. На скільки збільшиться сума елементів масиву \kw{a} після виконання
\begin{verbatim}
a[-2] += 3
\end{verbatim}


Якщо код некоректний, то в якості відповіді зазначити ERROR.\par

               
{\begin{center}Завдання {10}\end{center}}
\vspace{-6mm}

Написати функцію $f$, яка для файлу, заданого шляхом, розв’язує задачу:\\ Кожен з двох текстових файлів містить послідовність чисел, записану у форматі одне число на рядок.  Перевірити, чи задають вони одну ту саму послідовність.

В усіх випадках розміри файлів такі, що їх вміст повністю завантажений у пам'ять бути не може. Якість коду також оцінюється.



\vspace{10mm}\par

\newpage
\begin{minipage}[t]{9.7cm}
               
{\begin{center}Завдання {7}\end{center}}
\vspace{-6mm}

\begin{center}
	\adjustimage{max size={0.8\linewidth}{0.8\paperheight}}
	{./15BinTree/trees_exp/162.png}
\end{center}
Указати послідовність міток вершин дерева, зображеного на рис., що утвориться в результаті обходу дерева в прямому порядку. Мітки записувати через пробіл.\par Алгоритм починає роботу з кореня (на рис. це верхня вершина).\par

\end{minipage}
\vline \hspace{8mm}
\begin{minipage}[t]{8.0cm}
               
{\begin{center}Завдання {8}\end{center}}
\vspace{-6mm}
Для графа на рисунку з вершини 2 запущено алгоритм пошуку в глибину, що будує кістякове дерево. Списки суміжності вершин переглядаються алгоритмом за зростанням міток вершин.\par
\begin{center}
	\adjustimage{max size={0.9\linewidth}{0.9\paperheight}}
	{./16Graphs/Traversals/152.png}
\end{center}
\par Які з наведених ребер входять до складу отриманого кістякового дерева?\par
1) (3, 4)\par
2) (2, 5)\par
3) (1, 2)\par
4) (3, 6)\par
5) (1, 4)\par
6) (1, 3)\par
{\vskip 15pt} \par

\end{minipage}

               
{\begin{center}Завдання {9}\end{center}}
\vspace{-6mm}
Файл \code{'1.txt'} містить відомість з рядками, в яких через двокрапку записано поряд\-ковий номер (ціле), назву предмету (знаків пунктуації не містить), прізвище студента, ім'я студента, номер заліковки, оцінку в балах за 100 бальною шкалою.

Білі символи навколо роздільників не допускаються.

Приклад рядка файлу



\begin{verbatim}
13:algebra:surname:name:123456:77
\end{verbatim}


Нижче наведено код, який має залишити у файлі в кожному рядку тільки номер заліковки та предмет.

Рядок з прикладу має бути замінений на такий



\begin{verbatim}
123456:algebra
\end{verbatim}




Код:



\begin{verbatim}
regex = re.compile(r'XXX')
with open('1.txt', encoding='utf8') as f:
    lst = f.readlines()
for i in range(len(lst)):
    lst[i] = re.sub(regex, YYY, lst[i])
with open('1.txt', 'w', encoding='utf8') as f:
    f.writelines(lst)
\end{verbatim}







Який рядок треба записати на місце \kw{XXX} ?
Який рядок треба записати на місце \kw{YYY} ?\par

%\end {large}
}
\newpage

{{23.05.2025 Контрольна робота, варіант  \No \textbf { 29 }} група К1{\parbox {0.5 cm} \dotfill}  ПІБ {\parbox {2.5 cm} \dotfill}\par}
{
%\begin {large}
\begin{minipage}[t]{9.7cm}
               
{\begin{center}Завдання {1}\end{center}}
\vspace{-6mm}
Укажіть ТИП значення виразу.\par\par
\begin{verbatim}
[7,2,3,4]
\end{verbatim}

\vspace{15mm}
У завданнях 2, 3 вказати, що буде виведено за виконання. Повідомлення інтерпретатора про помилку позначати ERROR

               
{\begin{center}Завдання {2}\end{center}}
\vspace{-6mm}

\begin{verbatim}
try:
    t = {83:9, 53:8, 85:7, 53:9}
    t[83] = 5
    for x in t.values():
        print(x, sep=' ', end=' ')
except:
    print('ERR')
\end{verbatim}


\end{minipage}
\vline \hspace{8mm}
\begin{minipage}[t]{8.0cm}                
               
{\begin{center}Завдання {3}\end{center}}
\vspace{-6mm}

\begin{verbatim}
b = 0
j = 1

class D:
    j = 2
    
    def __init__(self):
        self.b = 3
        j = 4
        D.j = 5

obj = D()
try:
    print(b, end=' ')
    print(j, end=' ')
    print(D.j, end=' ')
    print(obj.b)
except:
    print('ERR')
\end{verbatim}

\end{minipage}

               
{\begin{center}Завдання {4}\end{center}}
\vspace{-6mm}
Нехай змінна \kw{s} позначає послідовність цілих, по якій можна ітерувати.

Написати вираз, що перевіряє, чи правда, що послідовність  містить числа, менші 13.\par

\vspace{10mm}
               
{\begin{center}Завдання {5}\end{center}}
\vspace{-6mm}
Змінні \kw{next} та \kw{prev} вузла списку позначають наступний та попередній за ним вузол відповідно. Починаючи з голови, у список записано послідовні цілі числа від~$-26$ до~$39$. Нехай змінна \kw{p} позначає вузол, що зберігає значення~$0$.
Яке число зберігається у вузлі \kw{p.next.prev.next.prev.next} ?\par

               
{\begin{center}Завдання {6}\end{center}}
\vspace{-6mm}
Значенням змінної \kw{a} є цілочисловий масив типу $numpy.ndarray$ розмірів $7{\times}3$. На скільки збільшиться сума елементів масиву \kw{a} після виконання
\begin{verbatim}
a += 2
\end{verbatim}


Якщо код некоректний, то в якості відповіді зазначити ERROR.\par

               
{\begin{center}Завдання {10}\end{center}}
\vspace{-6mm}

Написати функцію $f$, яка для файлу, заданого шляхом, розв’язує задачу:\\ Текстовий файл містить послідовність чисел $a_1$, $a_2$, $\ldots$, $a_t$, записану у форматі одне число на рядок. Знайти середнє арифметичне чисел $a_1$, $a_4$, $a_7$, $\ldots$ (починаючи з першого через два).

В усіх випадках розміри файлів такі, що їх вміст повністю завантажений у пам'ять бути не може. Якість коду також оцінюється.



\vspace{10mm}\par

\newpage
\begin{minipage}[t]{9.7cm}
               
{\begin{center}Завдання {7}\end{center}}
\vspace{-6mm}

\begin{center}
	\adjustimage{max size={0.8\linewidth}{0.8\paperheight}}
	{./15BinTree/trees_exp/103.png}
\end{center}
Указати послідовність міток вершин дерева, зображеного на рис., що утвориться в результаті обходу дерева в оберненому порядку. Мітки записувати через пробіл.\par Алгоритм починає роботу з кореня (на рис. це верхня вершина).\par

\end{minipage}
\vline \hspace{8mm}
\begin{minipage}[t]{8.0cm}
               
{\begin{center}Завдання {8}\end{center}}
\vspace{-6mm}
Для графа на рисунку з вершини 3 запущено алгоритм пошуку в глибину, що будує кістякове дерево. Списки суміжності вершин переглядаються алгоритмом за зростанням міток вершин.\par
\begin{center}
	\adjustimage{max size={0.9\linewidth}{0.9\paperheight}}
	{./16Graphs/Traversals/55.png}
\end{center}
\par Які з наведених ребер входять до складу отриманого кістякового дерева?\par
1) (0, 2)\par
2) (3, 6)\par
3) (2, 4)\par
4) (4, 6)\par
5) (2, 5)\par
6) (3, 5)\par
{\vskip 15pt} \par

\end{minipage}

               
{\begin{center}Завдання {9}\end{center}}
\vspace{-6mm}
Рядки журналу через двокрапку містять ім'я користувача (складається з латиниці), час події (фор\-мат як у прикладі), тип події (ERROR, WARNING, INFO, STATUS), код події (ціле).

Білі символи навколо роздільників не допускаються.

Приклад такого рядка присвоєно змінній \kw{s}.



\begin{verbatim}
s = 'username:2022-05-05 13-26:ERROR:404'
\end{verbatim}


Нижче наведено код, за виконання якого значенням змінної \kw{out} має стати код події.

Для наведеного прикладу значенням змінної \kw{out} має стати рядок \code{404}



\begin{verbatim}
res = re.fullmatch(r'XXX', s)
s = ''
out = YYY
\end{verbatim}


Код має працювати не тільки для конкретного рядка з прикладу, але й для кожного рядка з журналу.


Який рядок треба записати на місце \kw{XXX} ?
Який рядок треба записати на місце \kw{YYY} ?\par

%\end {large}
}
\newpage

{{23.05.2025 Контрольна робота, варіант  \No \textbf { 30 }} група К1{\parbox {0.5 cm} \dotfill}  ПІБ {\parbox {2.5 cm} \dotfill}\par}
{
%\begin {large}
\begin{minipage}[t]{9.7cm}
               
{\begin{center}Завдання {1}\end{center}}
\vspace{-6mm}
Укажіть ТИП значення виразу.\par\par
\begin{verbatim}
{x:x for x in range(7)}
\end{verbatim}

\vspace{15mm}
У завданнях 2, 3 вказати, що буде виведено за виконання. Повідомлення інтерпретатора про помилку позначати ERROR

               
{\begin{center}Завдання {2}\end{center}}
\vspace{-6mm}

\begin{verbatim}
try:
    d = {75:1, 56:1, 63:1, 47:2, 47:9}
    d[64] = 1
    for x in d.keys():
        print(x, sep=' ', end=' ')
except:
    print('ERR')
\end{verbatim}


\end{minipage}
\vline \hspace{8mm}
\begin{minipage}[t]{8.0cm}                
               
{\begin{center}Завдання {3}\end{center}}
\vspace{-6mm}

\begin{verbatim}
_d = 0

class C:
    _d = 1
    
    def __init__(self):
        self._d = 2
        _d = 3
        C._d = 4

obj = C()
obj._d = 5
try:
    print(_d, end=' ')
    print(obj._d, end=' ')
    print(C._d, end=' ')
    print(obj._d, end=' ')
    print(obj._C__d)
except:
    print('ERR')
\end{verbatim}

\end{minipage}

               
{\begin{center}Завдання {4}\end{center}}
\vspace{-6mm}
Нехай змінна \kw{s} позначає дуже довгу послідовність, по якій можна ітерувати.

Написати вираз-генератор, що задає підпослідовність її елементів, що не є числами типу \kw{int}.

Обмеження: усі елементи шуканої підпослідовності одночасно в пам'яті розміщені бути не можуть.\par

\vspace{10mm}
               
{\begin{center}Завдання {5}\end{center}}
\vspace{-6mm}
Змінні \kw{next} та \kw{prev} вузла списку позначають наступний та попередній за ним вузол відповідно. Починаючи з голови, у список записано послідовні цілі числа від~$-32$ до~$35$. Нехай змінна \kw{p} позначає вузол, що зберігає значення~$6$.
Яке число зберігається у вузлі \kw{p.prev.next.prev.next.prev} ?\par

               
{\begin{center}Завдання {6}\end{center}}
\vspace{-6mm}
Значенням змінної \kw{a} є цілочисловий масив типу $numpy.ndarray$ розмірів $6{\times}7$. На скільки збільшиться сума елементів масиву \kw{a} після виконання
\begin{verbatim}
a.T[-3] += 1
\end{verbatim}


Якщо код некоректний, то в якості відповіді зазначити ERROR.\par

               
{\begin{center}Завдання {10}\end{center}}
\vspace{-6mm}

Написати функцію $f$, яка для файлу, заданого шляхом, розв’язує задачу:\\ Текстовий файл містить послідовність чисел $a_1$, $a_2$, $\ldots$, $a_t$, записану у форматі одне число на рядок. Знайти третє найменше значення, що записане у файлі. Кожне значення, що записане у файлі, враховується стільки разів, скільки входить у файл. Для послідовності 2, 4, 12, 5, 2, 0 таким значенням буде 2.

В усіх випадках розміри файлів такі, що їх вміст повністю завантажений у пам'ять бути не може. Якість коду також оцінюється.



\vspace{10mm}\par

\newpage
\begin{minipage}[t]{9.7cm}
               
{\begin{center}Завдання {7}\end{center}}
\vspace{-6mm}

\begin{center}
	\adjustimage{max size={0.8\linewidth}{0.8\paperheight}}
	{./15BinTree/trees_exp/178.png}
\end{center}
Указати послідовність міток вершин дерева, зображеного на рис., що утвориться в результаті обходу дерева в прямому порядку. Мітки записувати через пробіл.\par Алгоритм починає роботу з кореня (на рис. це верхня вершина).\par

\end{minipage}
\vline \hspace{8mm}
\begin{minipage}[t]{8.0cm}
               
{\begin{center}Завдання {8}\end{center}}
\vspace{-6mm}
Для графа на рисунку з вершини 3 запущено алгоритм пошуку в глибину, що будує кістякове дерево. Списки суміжності вершин переглядаються алгоритмом за зростанням міток вершин.\par
\begin{center}
	\adjustimage{max size={0.9\linewidth}{0.9\paperheight}}
	{./16Graphs/Traversals/309.png}
\end{center}
\par Які з наведених ребер входять до складу отриманого кістякового дерева?\par
1) (2, 6)\par
2) (2, 5)\par
3) (2, 3)\par
4) (0, 4)\par
5) (2, 4)\par
6) (1, 4)\par
{\vskip 15pt} \par

\end{minipage}

               
{\begin{center}Завдання {9}\end{center}}
\vspace{-6mm}
Файл \code{'1.txt'} містить журнал з рядками, в яких через кому записано поряд\-ко\-вий номер запису в журналі (ціле), тип події (ERROR, WARNING, INFO, STATUS), ім'я користувача (складається з латиниці), час події (фор\-мат як у прикладі), код події (ціле).

Білі символи навколо роздільників не допускаються.

Приклад рядка файлу



\begin{verbatim}
13,ERROR,username,2022-05-05 13:26,404
\end{verbatim}


Нижче наведено код, який шукає усі коди помилок (подія ERROR), що зустріча\-ють\-ся в журналі, та зберігає їх у змінній \kw{codes}.



\begin{verbatim}
codes = set()
regex = re.compile(r'XXX')
with open('1.txt', encoding='utf8') as f:
    for s in f:
        res = re.fullmatch(regex, s)
        if not res: continue
        s = ''
        codes.add(YYY)
\end{verbatim}



Який рядок треба записати на місце \kw{XXX} ?
Який рядок треба записати на місце \kw{YYY} ?\par

%\end {large}
}
\newpage

{{23.05.2025 Контрольна робота, варіант  \No \textbf { 31 }} група К1{\parbox {0.5 cm} \dotfill}  ПІБ {\parbox {2.5 cm} \dotfill}\par}
{
%\begin {large}
\begin{minipage}[t]{9.7cm}
               
{\begin{center}Завдання {1}\end{center}}
\vspace{-6mm}
Укажіть ТИП значення виразу.\par\par
\begin{verbatim}
{1:3, 5:8}
\end{verbatim}

\vspace{15mm}
У завданнях 2, 3 вказати, що буде виведено за виконання. Повідомлення інтерпретатора про помилку позначати ERROR

               
{\begin{center}Завдання {2}\end{center}}
\vspace{-6mm}

\begin{verbatim}
try:
    t = {55:6, 19:9, 22:3, 53:1, 22:2}
    t[11] = 0
    for x in t.keys():
        print(x, sep=' ', end=' ')
except:
    print('ERR')
\end{verbatim}


\end{minipage}
\vline \hspace{8mm}
\begin{minipage}[t]{8.0cm}                
               
{\begin{center}Завдання {3}\end{center}}
\vspace{-6mm}

\begin{verbatim}
c = 0
m = 1

class B:
    c = 2
    
    def __init__(self):
        self.m = 3
        c = 4
        B.c = 5

obj = B()
try:
    print(c, end=' ')
    print(m, end=' ')
    print(B.c, end=' ')
    print(obj.m)
except:
    print('ERR')
\end{verbatim}

\end{minipage}

               
{\begin{center}Завдання {4}\end{center}}
\vspace{-6mm}
Записати ВИРАЗ, значенням якого є множина, що складається з кубів цілих чисел від 20 до 2015 (включно), що діляться на 7.\par

\vspace{10mm}
               
{\begin{center}Завдання {5}\end{center}}
\vspace{-6mm}
Змінні \kw{next} та \kw{prev} вузла списку позначають наступний та попередній за ним вузол відповідно. Починаючи з голови, у список записано послідовні цілі числа від~$-51$ до~$30$. Нехай змінна \kw{p} позначає вузол, що зберігає значення~$1$.
Яке число зберігається у вузлі \kw{p.prev.prev.next.next.next} ?\par

               
{\begin{center}Завдання {6}\end{center}}
\vspace{-6mm}
Значенням змінної \kw{a} є цілочисловий масив типу $numpy.ndarray$ розмірів $7{\times}5$. На скільки збільшиться сума елементів масиву \kw{a} після виконання
\begin{verbatim}
a[2] += 2
\end{verbatim}


Якщо код некоректний, то в якості відповіді зазначити ERROR.\par

               
{\begin{center}Завдання {10}\end{center}}
\vspace{-6mm}

Написати функцію $f$, яка для файлу, заданого шляхом, розв’язує задачу:\\ Текстовий файл містить послідовність чисел $a_1$, $a_2$, $\ldots$, $a_t$, записану у форматі одне число на рядок. Знайти третє найбільше значення, що записане у файлі. Кожне значення, що записане у файлі, враховується тільки один раз. Для послідовності 1, 4, 12, 5, 1, 0 таким значенням буде 4.

В усіх випадках розміри файлів такі, що їх вміст повністю завантажений у пам'ять бути не може. Якість коду також оцінюється.



\vspace{10mm}\par

\newpage
\begin{minipage}[t]{9.7cm}
               
{\begin{center}Завдання {7}\end{center}}
\vspace{-6mm}

\begin{center}
	\adjustimage{max size={0.8\linewidth}{0.8\paperheight}}
	{./15BinTree/trees_exp/55.png}
\end{center}
Указати послідовність міток вершин дерева, зображеного на рис., що утвориться в результаті обходу дерева в оберненому порядку. Мітки записувати через пробіл.\par Алгоритм починає роботу з кореня (на рис. це верхня вершина).\par

\end{minipage}
\vline \hspace{8mm}
\begin{minipage}[t]{8.0cm}
               
{\begin{center}Завдання {8}\end{center}}
\vspace{-6mm}
Для графа на рисунку з вершини 1 запущено алгоритм пошуку в ширину, що будує кістякове дерево. Списки суміжності вершин переглядаються алгоритмом за зростанням міток вершин.\par
\begin{center}
	\adjustimage{max size={0.9\linewidth}{0.9\paperheight}}
	{./16Graphs/Traversals/74.png}
\end{center}
\par Які з наведених ребер входять до складу отриманого кістякового дерева?\par
1) (5, 6)\par
2) (0, 3)\par
3) (3, 6)\par
4) (1, 3)\par
5) (3, 4)\par
6) (1, 5)\par
{\vskip 15pt} \par

\end{minipage}

               
{\begin{center}Завдання {9}\end{center}}
\vspace{-6mm}
Файл \code{'1.txt'} містить відомість з рядками, в яких через кому записано поряд\-ковий номер (ціле), назву предмету (знаків пунктуації не містить), прізвище студента, ім'я студента, номер заліковки, оцінку в балах за 100 бальною шкалою.

Білі символи навколо роздільників не допускаються.

Приклад рядка файлу



\begin{verbatim}
13,algebra,surname,name,123456,77
\end{verbatim}


Нижче наведено код, який має залишити у файлі в кожному рядку тільки номер заліковки та предмет.

Рядок з прикладу має бути замінений на такий



\begin{verbatim}
123456,algebra
\end{verbatim}




Код



\begin{verbatim}
regex = re.compile(r'XXX')
with open('1.txt', encoding='utf8') as f:
    lst = f.readlines()
for i in range(len(lst)):
    lst[i] = re.sub(regex, YYY, lst[i])
with open('1.txt', 'w', encoding='utf8') as f:
    f.writelines(lst)
\end{verbatim}







Який рядок треба записати на місце \kw{XXX} ?
Який рядок треба записати на місце \kw{YYY} ?\par

%\end {large}
}
\newpage

{{23.05.2025 Контрольна робота, варіант  \No \textbf { 32 }} група К1{\parbox {0.5 cm} \dotfill}  ПІБ {\parbox {2.5 cm} \dotfill}\par}
{
%\begin {large}
\begin{minipage}[t]{9.7cm}
               
{\begin{center}Завдання {1}\end{center}}
\vspace{-6mm}
Укажіть ТИП значення виразу.\par\par
\begin{verbatim}
[{1,2}, 5]
\end{verbatim}

\vspace{15mm}
У завданнях 2, 3 вказати, що буде виведено за виконання. Повідомлення інтерпретатора про помилку позначати ERROR

               
{\begin{center}Завдання {2}\end{center}}
\vspace{-6mm}

\begin{verbatim}
try:
    s = {89:2, 53:8, 80:8, 53:7}
    s[30] = 2
    for x in s :
        print(x, sep=' ', end=' ')
except:
    print('ERR')
\end{verbatim}


\end{minipage}
\vline \hspace{8mm}
\begin{minipage}[t]{8.0cm}                
               
{\begin{center}Завдання {3}\end{center}}
\vspace{-6mm}

\begin{verbatim}
a = 0

class A:
    a = 1
    
    def __init__(self):
        self.a = 2
        a = 3
        A.a = 4

obj = A()
A.a = 5
try:
    print(a, end=' ')
    print(obj.a, end=' ')
    print(A.a, end=' ')
    print(obj.a, end=' ')
    print(obj._A__a)
except:
    print('ERR')
\end{verbatim}

\end{minipage}

               
{\begin{center}Завдання {4}\end{center}}
\vspace{-6mm}
Записати ВИРАЗ, значенням якого є множина, що складається з цілих чисел від 18 до 2001 (включно), що не діляться на жодне з чисел 2, 3.\par

\vspace{10mm}
               
{\begin{center}Завдання {5}\end{center}}
\vspace{-6mm}
Змінні \kw{next} та \kw{prev} вузла списку позначають наступний та попередній за ним вузол відповідно. Починаючи з голови, у список записано послідовні цілі числа від~$-63$ до~$45$. Нехай змінна \kw{p} позначає вузол, що зберігає значення~$-9$.
Яке число зберігається у вузлі \kw{p.next.next.prev.prev.prev} ?\par

               
{\begin{center}Завдання {6}\end{center}}
\vspace{-6mm}
Значенням змінної \kw{a} є цілочисловий масив типу $numpy.ndarray$ розмірів $5{\times}6$. На скільки збільшиться сума елементів масиву \kw{a} після виконання
\begin{verbatim}
a.T[0] += 3
\end{verbatim}


Якщо код некоректний, то в якості відповіді зазначити ERROR.\par

               
{\begin{center}Завдання {10}\end{center}}
\vspace{-6mm}

Написати функцію $f$, яка для файлу, заданого шляхом, розв’язує задачу:\\ Текстовий файл містить послідовність чисел $a_1$, $a_2$, $\ldots$, $a_t$, записану у форматі одне число на рядок. Знайти середнє арифметичне чисел $a_1$, $a_4$, $a_7$, $\ldots$ (починаючи з першого через два).

В усіх випадках розміри файлів такі, що їх вміст повністю завантажений у пам'ять бути не може. Якість коду також оцінюється.



\vspace{10mm}\par

\newpage
\begin{minipage}[t]{9.7cm}
               
{\begin{center}Завдання {7}\end{center}}
\vspace{-6mm}

\begin{center}
	\adjustimage{max size={0.8\linewidth}{0.8\paperheight}}
	{./15BinTree/trees_exp/124.png}
\end{center}
Указати послідовність міток вершин дерева, зображеного на рис., що утвориться в результаті обходу дерева в прямому порядку. Мітки записувати через пробіл.\par Алгоритм починає роботу з кореня (на рис. це верхня вершина).\par

\end{minipage}
\vline \hspace{8mm}
\begin{minipage}[t]{8.0cm}
               
{\begin{center}Завдання {8}\end{center}}
\vspace{-6mm}
Для графа на рисунку з вершини 6 запущено алгоритм пошуку в ширину, що будує кістякове дерево. Списки суміжності вершин переглядаються алгоритмом за зростанням міток вершин.\par
\begin{center}
	\adjustimage{max size={0.9\linewidth}{0.9\paperheight}}
	{./16Graphs/Traversals/306.png}
\end{center}
\par Які з наведених ребер входять до складу отриманого кістякового дерева?\par
1) (4, 6)\par
2) (3, 5)\par
3) (1, 6)\par
4) (0, 3)\par
5) (0, 1)\par
6) (0, 5)\par
{\vskip 15pt} \par

\end{minipage}

               
{\begin{center}Завдання {9}\end{center}}
\vspace{-6mm}
Рядки журналу через двокрапку містять тип події (ERROR, WARNING, INFO, STATUS), ім'я користувача (складається з латиниці), час події (фор\-мат як у прикладі), код події (ціле).

Білі символи навколо роздільників не допускаються.

Приклад такого рядка присвоєно змінній \kw{s}.



\begin{verbatim}
s = 'ERROR:username:2022-05-05 13-26:404'
\end{verbatim}


Нижче наведено код, за виконання якого значенням змінної \kw{out} має стати ім'я користувача.

Для наведеного прикладу значенням змінної \kw{out} має стати рядок \code{username}



\begin{verbatim}
res = re.fullmatch(r'XXX', s)
s = ''
out = YYY
\end{verbatim}


Код має працювати не тільки для конкретного рядка з прикладу, але й для кожного рядка з журналу.


Який рядок треба записати на місце \kw{XXX} ?
Який рядок треба записати на місце \kw{YYY} ?\par

%\end {large}
}
\newpage

{{23.05.2025 Контрольна робота, варіант  \No \textbf { 33 }} група К1{\parbox {0.5 cm} \dotfill}  ПІБ {\parbox {2.5 cm} \dotfill}\par}
{
%\begin {large}
\begin{minipage}[t]{9.7cm}
               
{\begin{center}Завдання {1}\end{center}}
\vspace{-6mm}
Укажіть ТИП значення виразу.\par\par
\begin{verbatim}
{y:None for y in range(-8,3)}
\end{verbatim}

\vspace{15mm}
У завданнях 2, 3 вказати, що буде виведено за виконання. Повідомлення інтерпретатора про помилку позначати ERROR

               
{\begin{center}Завдання {2}\end{center}}
\vspace{-6mm}

\begin{verbatim}
try:
    d = {87:5, 64:6, 24:5, 35:7, 35:2}
    d[16] = 2
    for x in d.values():
        print(x, sep=' ', end=' ')
except:
    print('ERR')
\end{verbatim}


\end{minipage}
\vline \hspace{8mm}
\begin{minipage}[t]{8.0cm}                
               
{\begin{center}Завдання {3}\end{center}}
\vspace{-6mm}

\begin{verbatim}
d = 0
h = 1

class C:
    h = 2
    
    def __init__(self):
        global h
        self.d = 3
        h = 4
        C.h = 5

obj = C()
try:
    print(d, end=' ')
    print(h, end=' ')
    print(C.h, end=' ')
    print(obj.d)
except:
    print('ERR')
\end{verbatim}

\end{minipage}

               
{\begin{center}Завдання {4}\end{center}}
\vspace{-6mm}
Нехай змінна \kw{s} позначає дуже довгу послідовність, по якій можна ітерувати.

Написати вираз-генератор, що задає підпослідовність її елементів, що не є числами типу \kw{float}.

Обмеження: усі елементи шуканої підпослідовності одночасно в пам'яті розміщені бути не можуть.\par

\vspace{10mm}
               
{\begin{center}Завдання {5}\end{center}}
\vspace{-6mm}
Змінні \kw{next} та \kw{prev} вузла списку позначають наступний та попередній за ним вузол відповідно. Починаючи з голови, у список записано послідовні цілі числа від~$-61$ до~$41$. Нехай змінна \kw{p} позначає вузол, що зберігає значення~$-5$.
Яке число зберігається у вузлі \kw{p.next.next.prev.next.next} ?\par

               
{\begin{center}Завдання {6}\end{center}}
\vspace{-6mm}
Значенням змінної \kw{a} є цілочисловий масив типу $numpy.ndarray$ розмірів $6{\times}4$. На скільки збільшиться сума елементів масиву \kw{a} після виконання
\begin{verbatim}
a.T[1] += 1
\end{verbatim}


Якщо код некоректний, то в якості відповіді зазначити ERROR.\par

               
{\begin{center}Завдання {10}\end{center}}
\vspace{-6mm}

Написати функцію $f$, яка для файлу, заданого шляхом, розв’язує задачу:\\ Текстовий файл містить послідовність чисел $a_1$, $a_2$, $\ldots$, $a_t$, записану у форматі одне число на рядок. Знайти третє найбільше значення, що записане у файлі. Кожне значення, що записане у файлі, враховується стільки разів, скільки входить у файл. Для послідовності 2, 4, 12, 5, 2, 0 таким значенням буде 2.

В усіх випадках розміри файлів такі, що їх вміст повністю завантажений у пам'ять бути не може. Якість коду також оцінюється.



\vspace{10mm}\par

\newpage
\begin{minipage}[t]{9.7cm}
               
{\begin{center}Завдання {7}\end{center}}
\vspace{-6mm}

\begin{center}
	\adjustimage{max size={0.8\linewidth}{0.8\paperheight}}
	{./15BinTree/trees_exp/130.png}
\end{center}
Указати послідовність міток вершин дерева, зображеного на рис., що утвориться в результаті обходу дерева в прямому порядку. Мітки записувати через пробіл.\par Алгоритм починає роботу з кореня (на рис. це верхня вершина).\par

\end{minipage}
\vline \hspace{8mm}
\begin{minipage}[t]{8.0cm}
               
{\begin{center}Завдання {8}\end{center}}
\vspace{-6mm}
Для графа на рисунку з вершини 6 запущено алгоритм пошуку в глибину, що будує кістякове дерево. Списки суміжності вершин переглядаються алгоритмом за зростанням міток вершин.\par
\begin{center}
	\adjustimage{max size={0.9\linewidth}{0.9\paperheight}}
	{./16Graphs/Traversals/23.png}
\end{center}
\par Які з наведених ребер входять до складу отриманого кістякового дерева?\par
1) (3, 6)\par
2) (5, 6)\par
3) (2, 5)\par
4) (3, 4)\par
5) (0, 4)\par
6) (3, 5)\par
{\vskip 15pt} \par

\end{minipage}

               
{\begin{center}Завдання {9}\end{center}}
\vspace{-6mm}
Рядки відомості через кому містять порядковий номер (ціле), назву предмету (зна\-ків пунктуації не містить), прізвище студента, ім'я студента, оцінку в балах за 100 бальною шкалою.

Білі символи навколо роздільників не допускаються.

Приклад такого рядка присвоєно змінній \kw{s}.



\begin{verbatim}
s = '13,algebra,surname,name,77'
\end{verbatim}


Нижче наведено код, за виконання якого значенням змінної \kw{out} має стати прізвище студента.

Для наведеного прикладу значенням змінної \kw{out} має стати рядок \code{surname}



\begin{verbatim}
res = re.fullmatch(r'XXX', s)
s = ''
out = YYY
\end{verbatim}


Код має працювати не тільки для конкретного рядка з прикладу, але й для кожного рядка з відомості.


Який рядок треба записати на місце \kw{XXX} ?
Який рядок треба записати на місце \kw{YYY} ?\par

%\end {large}
}
\newpage

{{23.05.2025 Контрольна робота, варіант  \No \textbf { 34 }} група К1{\parbox {0.5 cm} \dotfill}  ПІБ {\parbox {2.5 cm} \dotfill}\par}
{
%\begin {large}
\begin{minipage}[t]{9.7cm}
               
{\begin{center}Завдання {1}\end{center}}
\vspace{-6mm}
Укажіть ТИП значення виразу.\par\par
\begin{verbatim}
[x for x in range(7)]
\end{verbatim}

\vspace{15mm}
У завданнях 2, 3 вказати, що буде виведено за виконання. Повідомлення інтерпретатора про помилку позначати ERROR

               
{\begin{center}Завдання {2}\end{center}}
\vspace{-6mm}

\begin{verbatim}
try:
    d = {63:3, 57:0, 65:0, 46:3}
    d[18] = 4
    for x in d.items():
        print(x, sep=' ', end=' ')
except:
    print('ERR')
\end{verbatim}


\end{minipage}
\vline \hspace{8mm}
\begin{minipage}[t]{8.0cm}                
               
{\begin{center}Завдання {3}\end{center}}
\vspace{-6mm}

\begin{verbatim}
c = 0
j = 1

class C:
    c = 2
    
    def __init__(self):
        global c
        self.j = 3
        c = 4
        C.c = 5

obj = C()
try:
    print(c, end=' ')
    print(j, end=' ')
    print(C.c, end=' ')
    print(obj.j)
except:
    print('ERR')
\end{verbatim}

\end{minipage}

               
{\begin{center}Завдання {4}\end{center}}
\vspace{-6mm}
Нехай змінна \kw{s} позначає дуже довгу послідовність цілих, по якій можна ітерувати.

Написати вираз-генератор, що задає підпослідовність її елементів, що менші за задане число num.

Обмеження: усі елементи шуканої підпослідовності одночасно в пам'яті розміщені бути не можуть.\par

\vspace{10mm}
               
{\begin{center}Завдання {5}\end{center}}
\vspace{-6mm}
Змінні \kw{next} та \kw{prev} вузла списку позначають наступний та попередній за ним вузол відповідно. Починаючи з голови, у список записано послідовні цілі числа від~$-66$ до~$23$. Нехай змінна \kw{p} позначає вузол, що зберігає значення~$8$.
Яке число зберігається у вузлі \kw{p.prev.prev.next.prev.prev} ?\par

               
{\begin{center}Завдання {6}\end{center}}
\vspace{-6mm}
Значенням змінної \kw{a} є цілочисловий масив типу $numpy.ndarray$ розмірів $3{\times}7$. На скільки збільшиться сума елементів масиву \kw{a} після виконання
\begin{verbatim}
a[-2] += 1
\end{verbatim}


Якщо код некоректний, то в якості відповіді зазначити ERROR.\par

               
{\begin{center}Завдання {10}\end{center}}
\vspace{-6mm}

Написати функцію $f$, яка для файлу, заданого шляхом, розв’язує задачу:\\ Кожен з двох текстових файлів містить послідовність чисел, записану у форматі одне число на рядок.  Перевірити, чи задають вони одну ту саму послідовність.

В усіх випадках розміри файлів такі, що їх вміст повністю завантажений у пам'ять бути не може. Якість коду також оцінюється.



\vspace{10mm}\par

\newpage
\begin{minipage}[t]{9.7cm}
               
{\begin{center}Завдання {7}\end{center}}
\vspace{-6mm}

\begin{center}
	\adjustimage{max size={0.8\linewidth}{0.8\paperheight}}
	{./15BinTree/trees_exp/76.png}
\end{center}
Указати послідовність міток вершин дерева, зображеного на рис., що утвориться в результаті обходу дерева в оберненому порядку. Мітки записувати через пробіл.\par Алгоритм починає роботу з кореня (на рис. це верхня вершина).\par

\end{minipage}
\vline \hspace{8mm}
\begin{minipage}[t]{8.0cm}
               
{\begin{center}Завдання {8}\end{center}}
\vspace{-6mm}
Для графа на рисунку з вершини 4 запущено алгоритм пошуку в ширину, що будує кістякове дерево. Списки суміжності вершин переглядаються алгоритмом за зростанням міток вершин.\par
\begin{center}
	\adjustimage{max size={0.9\linewidth}{0.9\paperheight}}
	{./16Graphs/Traversals/151.png}
\end{center}
\par Які з наведених ребер входять до складу отриманого кістякового дерева?\par
1) (5, 6)\par
2) (2, 5)\par
3) (4, 5)\par
4) (3, 4)\par
5) (0, 4)\par
6) (2, 3)\par
{\vskip 15pt} \par

\end{minipage}

               
{\begin{center}Завдання {9}\end{center}}
\vspace{-6mm}
Рядки відомості через двокрапку містять порядковий номер (ціле), пріз\-ви\-ще студен\-та, ім'я студента, номер заліковки, оцінку в балах за 100 бальною шкалою.

Білі символи навколо роздільників не допускаються.

Приклад такого рядка присвоєно змінній \kw{s}.



\begin{verbatim}
s = '13:surname:name:123456:77'
\end{verbatim}


Нижче наведено код, за виконання якого для рядка \kw{s} значенням змінної \kw{out} має стати номер заліковки.

Для наведеного прикладу значенням змінної \kw{out} має стати рядок \code{123456} .



\begin{verbatim}
res = re.fullmatch(r'XXX', s)
s = ''
out = YYY
\end{verbatim}


Код має працювати не тільки для конкретного рядка з прикладу, але й для кожного рядка з відомості.


Який рядок треба записати на місце \kw{XXX} ?
Який рядок треба записати на місце \kw{YYY} ?\par

%\end {large}
}
\newpage

{{23.05.2025 Контрольна робота, варіант  \No \textbf { 35 }} група К1{\parbox {0.5 cm} \dotfill}  ПІБ {\parbox {2.5 cm} \dotfill}\par}
{
%\begin {large}
\begin{minipage}[t]{9.7cm}
               
{\begin{center}Завдання {1}\end{center}}
\vspace{-6mm}
Укажіть ТИП значення виразу.\par\par
\begin{verbatim}
(11,2,3)
\end{verbatim}

\vspace{15mm}
У завданнях 2, 3 вказати, що буде виведено за виконання. Повідомлення інтерпретатора про помилку позначати ERROR

               
{\begin{center}Завдання {2}\end{center}}
\vspace{-6mm}

\begin{verbatim}
try:
    t = {10:9, 24:6, 37:8, 10:2}
    t[77] = 1
    for x in t.items():
        print(*x, sep=' ', end=' ')
except:
    print('ERR')
\end{verbatim}


\end{minipage}
\vline \hspace{8mm}
\begin{minipage}[t]{8.0cm}                
               
{\begin{center}Завдання {3}\end{center}}
\vspace{-6mm}

\begin{verbatim}
__d = 0

class B:
    __d = 1
    
    def __init__(self):
        self.__d = 2
        __d = 3
        B.__d = 4

obj = B()
obj.__d = 5
try:
    print(__d, end=' ')
    print(obj.__d, end=' ')
    print(B.__d, end=' ')
    print(obj.__d, end=' ')
    print(obj._B__d)
except:
    print('ERR')
\end{verbatim}

\end{minipage}

               
{\begin{center}Завдання {4}\end{center}}
\vspace{-6mm}
Нехай змінна \kw{s} позначає послідовність цілих, по якій можна ітерувати.

Написати вираз, що перевіряє, чи правда, що послідовність  містить числа, більші за 5.\par

\vspace{10mm}
               
{\begin{center}Завдання {5}\end{center}}
\vspace{-6mm}
Змінні \kw{next} та \kw{prev} вузла списку позначають наступний та попередній за ним вузол відповідно. Починаючи з голови, у список записано послідовні цілі числа від~$-67$ до~$47$. Нехай змінна \kw{p} позначає вузол, що зберігає значення~$9$.
Яке число зберігається у вузлі \kw{p.prev.next.prev.prev.next} ?\par

               
{\begin{center}Завдання {6}\end{center}}
\vspace{-6mm}
Значенням змінної \kw{a} є цілочисловий масив типу $numpy.ndarray$ розмірів $4{\times}3$. На скільки збільшиться сума елементів масиву \kw{a} після виконання
\begin{verbatim}
a[-1, 2] += 2
\end{verbatim}


Якщо код некоректний, то в якості відповіді зазначити ERROR.\par

               
{\begin{center}Завдання {10}\end{center}}
\vspace{-6mm}

Написати функцію $f$, яка для файлу, заданого шляхом, розв’язує задачу:\\ Текстовий файл містить послідовність чисел $a_1$, $a_2$, $\ldots$, $a_t$, записану у форматі одне число на рядок. Знайти третє найменше значення, що записане у файлі. Кожне значення, що записане у файлі, враховується стільки разів, скільки входить у файл. Для послідовності 2, 4, 12, 5, 2, 0 таким значенням буде 2.

В усіх випадках розміри файлів такі, що їх вміст повністю завантажений у пам'ять бути не може. Якість коду також оцінюється.



\vspace{10mm}\par

\newpage
\begin{minipage}[t]{9.7cm}
               
{\begin{center}Завдання {7}\end{center}}
\vspace{-6mm}

\begin{center}
	\adjustimage{max size={0.8\linewidth}{0.8\paperheight}}
	{./15BinTree/trees_exp/145.png}
\end{center}
Указати послідовність міток вершин дерева, зображеного на рис., що утвориться в результаті обходу дерева в прямому порядку. Мітки записувати через пробіл.\par Алгоритм починає роботу з кореня (на рис. це верхня вершина).\par

\end{minipage}
\vline \hspace{8mm}
\begin{minipage}[t]{8.0cm}
               
{\begin{center}Завдання {8}\end{center}}
\vspace{-6mm}
Для графа на рисунку з вершини 0 запущено алгоритм пошуку в ширину, що будує кістякове дерево. Списки суміжності вершин переглядаються алгоритмом за зростанням міток вершин.\par
\begin{center}
	\adjustimage{max size={0.9\linewidth}{0.9\paperheight}}
	{./16Graphs/Traversals/199.png}
\end{center}
\par Які з наведених ребер входять до складу отриманого кістякового дерева?\par
1) (2, 3)\par
2) (4, 5)\par
3) (2, 4)\par
4) (3, 5)\par
5) (3, 4)\par
6) (5, 6)\par
{\vskip 15pt} \par

\end{minipage}

               
{\begin{center}Завдання {9}\end{center}}
\vspace{-6mm}
Файл \code{'1.txt'} містить журнал з рядками, в яких через двокрапку записано поряд\-ко\-вий номер запису (ціле), ім'я користувача (складається з латиниці), час події (фор\-мат як у прикладі), тип події (ERROR, WARNING, INFO, STATUS), код події (ціле).

Білі символи навколо роздільників не допускаються.

Приклад рядка журналу



\begin{verbatim}
12:username:2022-05-05 13-26:ERROR:404
\end{verbatim}


Нижче наведено код, який шукає усі попередження (подія WARNING), що зустрі\-ча\-ються в журналі і відбулися під час обідньої перерви деканату (з 13:00 по 13:59 включно), та зберігає їх коди в змінній \kw{codes}.



\begin{verbatim}
codes = set()
regex = re.compile(r'XXX')
with open('1.txt', encoding='utf8') as f:
    for s in f:
        res = re.fullmatch(regex, s)
        if not res: continue
        s = ''
        codes.add(YYY)
\end{verbatim}





Який рядок треба записати на місце \kw{XXX} ?
Який рядок треба записати на місце \kw{YYY} ?\par

%\end {large}
}
\newpage

{{23.05.2025 Контрольна робота, варіант  \No \textbf { 36 }} група К1{\parbox {0.5 cm} \dotfill}  ПІБ {\parbox {2.5 cm} \dotfill}\par}
{
%\begin {large}
\begin{minipage}[t]{9.7cm}
               
{\begin{center}Завдання {1}\end{center}}
\vspace{-6mm}
Укажіть ТИП значення виразу.\par\par
\begin{verbatim}
{y:13 for y in range(-7,7)}
\end{verbatim}

\vspace{15mm}
У завданнях 2, 3 вказати, що буде виведено за виконання. Повідомлення інтерпретатора про помилку позначати ERROR

               
{\begin{center}Завдання {2}\end{center}}
\vspace{-6mm}

\begin{verbatim}
try:
    s = {34:2, 13:9, 22:0, 34:8}
    s[69] = 2
    for x, y in s.items():
        print(x, y, sep=' ', end=' ')
except:
    print('ERR')
\end{verbatim}


\end{minipage}
\vline \hspace{8mm}
\begin{minipage}[t]{8.0cm}                
               
{\begin{center}Завдання {3}\end{center}}
\vspace{-6mm}

\begin{verbatim}
b = 0
k = 1

class D:
    k = 2
    
    def __init__(self):
        self.b = 3
        k = 4
        D.k = 5

obj = D()
try:
    print(b, end=' ')
    print(k, end=' ')
    print(D.k, end=' ')
    print(obj.b)
except:
    print('ERR')
\end{verbatim}

\end{minipage}

               
{\begin{center}Завдання {4}\end{center}}
\vspace{-6mm}

 Задано числову послідовність \kw{s} довжини 6000. Записати ВИРАЗ, значенням якого є список, що складається з другої половини послідовності, записаної в оберненому порядку.\par

\vspace{10mm}
               
{\begin{center}Завдання {5}\end{center}}
\vspace{-6mm}
Змінні \kw{next} та \kw{prev} вузла списку позначають наступний та попередній за ним вузол відповідно. Починаючи з голови, у список записано послідовні цілі числа від~$-47$ до~$34$. Нехай змінна \kw{p} позначає вузол, що зберігає значення~$5$.
Яке число зберігається у вузлі \kw{p.next.prev.next.next.prev} ?\par

               
{\begin{center}Завдання {6}\end{center}}
\vspace{-6mm}
Значенням змінної \kw{a} є цілочисловий масив типу $numpy.ndarray$ розмірів $5{\times}6$. На скільки збільшиться сума елементів масиву \kw{a} після виконання
\begin{verbatim}
a.T[1] += 3
\end{verbatim}


Якщо код некоректний, то в якості відповіді зазначити ERROR.\par

               
{\begin{center}Завдання {10}\end{center}}
\vspace{-6mm}

Написати функцію $f$, яка для файлу, заданого шляхом, розв’язує задачу:\\ Текстовий файл містить послідовність чисел $a_1$, $a_2$, $\ldots$, $a_t$, записану у форматі одне число на рядок. Знайти третє найменше значення, що записане у файлі. Кожне значення, що записане у файлі, враховується тільки один раз. \par
Для послідовності 1, 4, 12, 5, 1, 0 таким значенням буде 4.

В усіх випадках розміри файлів такі, що їх вміст повністю завантажений у пам'ять бути не може. Якість коду також оцінюється.



\vspace{10mm}\par

\newpage
\begin{minipage}[t]{9.7cm}
               
{\begin{center}Завдання {7}\end{center}}
\vspace{-6mm}

\begin{center}
	\adjustimage{max size={0.8\linewidth}{0.8\paperheight}}
	{./15BinTree/trees_exp/100.png}
\end{center}
Указати послідовність міток вершин дерева, зображеного на рис., що утвориться в результаті обходу дерева в оберненому порядку. Мітки записувати через пробіл.\par Алгоритм починає роботу з кореня (на рис. це верхня вершина).\par

\end{minipage}
\vline \hspace{8mm}
\begin{minipage}[t]{8.0cm}
               
{\begin{center}Завдання {8}\end{center}}
\vspace{-6mm}
Для графа на рисунку з вершини 1 запущено алгоритм пошуку в ширину, що будує кістякове дерево. Списки суміжності вершин переглядаються алгоритмом за зростанням міток вершин.\par
\begin{center}
	\adjustimage{max size={0.9\linewidth}{0.9\paperheight}}
	{./16Graphs/Traversals/347.png}
\end{center}
\par Які з наведених ребер входять до складу отриманого кістякового дерева?\par
1) (1, 6)\par
2) (4, 6)\par
3) (1, 3)\par
4) (0, 5)\par
5) (3, 5)\par
6) (0, 6)\par
{\vskip 15pt} \par

\end{minipage}

               
{\begin{center}Завдання {9}\end{center}}
\vspace{-6mm}
Файл \code{'1.txt'} містить відомість з рядками, в яких через двокрапку записа\-но поряд\-ковий номер (ціле), назву предмету (знаків пунктуації не містить), прізвище студента, ім'я студента, номер заліковки, оцінку в балах за 100 бальною шкалою.

Білі символи навколо роздільників не допускаються.

Приклад рядка файлу



\begin{verbatim}
13:algebra:surname:name:123456:77
\end{verbatim}


Нижче наведено код, який має залишити у файлі в кожному рядку тільки номер заліковки, назву предмету та оцінку в балах за 100-бальною шкалою (у заданій послі\-дов\-ності).

Рядок з прикладу має бути замінений на такий



\begin{verbatim}
123456:algebra:77
\end{verbatim}




Код:



\begin{verbatim}
regex = re.compile(r'XXX')
with open('1.txt', encoding='utf8') as f:
    lst = f.readlines()
for i in range(len(lst)):
    lst[i] = re.sub(regex, YYY, lst[i])
with open('1.txt', 'w', encoding='utf8') as f:
    f.writelines(lst)
\end{verbatim}





Який рядок треба записати на місце \kw{XXX} ?
Який рядок треба записати на місце \kw{YYY} ?\par

%\end {large}
}
\newpage

{{23.05.2025 Контрольна робота, варіант  \No \textbf { 37 }} група К1{\parbox {0.5 cm} \dotfill}  ПІБ {\parbox {2.5 cm} \dotfill}\par}
{
%\begin {large}
\begin{minipage}[t]{9.7cm}
               
{\begin{center}Завдання {1}\end{center}}
\vspace{-6mm}
Укажіть ТИП значення виразу.\par\par
\begin{verbatim}
[x for x in range(7)]
\end{verbatim}

\vspace{15mm}
У завданнях 2, 3 вказати, що буде виведено за виконання. Повідомлення інтерпретатора про помилку позначати ERROR

               
{\begin{center}Завдання {2}\end{center}}
\vspace{-6mm}

\begin{verbatim}
try:
    d = {21:9, 57:5, 89:5, 89:9}
    d[32] = 8
    for x in d.items():
        print(x, sep=' ', end=' ')
except:
    print('ERR')
\end{verbatim}


\end{minipage}
\vline \hspace{8mm}
\begin{minipage}[t]{8.0cm}                
               
{\begin{center}Завдання {3}\end{center}}
\vspace{-6mm}

\begin{verbatim}
a = 0
n = 1

class B:
    a = 2
    
    def __init__(self):
        global a
        self.a = 3
        a = 4
        B.a = 5

obj = B()
try:
    print(a, end=' ')
    print(n, end=' ')
    print(B.a, end=' ')
    print(obj.a)
except:
    print('ERR')
\end{verbatim}

\end{minipage}

               
{\begin{center}Завдання {4}\end{center}}
\vspace{-6mm}
Нехай змінна \kw{s} позначає послідовність цілих, по якій можна ітерувати.

Написати вираз, що перевіряє, чи правда, що послідовність  містить від’ємні числа.\par

\vspace{10mm}
               
{\begin{center}Завдання {5}\end{center}}
\vspace{-6mm}
Змінні \kw{next} та \kw{prev} вузла списку позначають наступний та попередній за ним вузол відповідно. Починаючи з голови, у список записано послідовні цілі числа від~$-22$ до~$28$. Нехай змінна \kw{p} позначає вузол, що зберігає значення~$3$.
Яке число зберігається у вузлі \kw{p.prev.next.prev.next.next} ?\par

               
{\begin{center}Завдання {6}\end{center}}
\vspace{-6mm}
Значенням змінної \kw{a} є цілочисловий масив типу $numpy.ndarray$ розмірів $4{\times}4$. На скільки збільшиться сума елементів масиву \kw{a} після виконання
\begin{verbatim}
a[0] += 2
\end{verbatim}


Якщо код некоректний, то в якості відповіді зазначити ERROR.\par

               
{\begin{center}Завдання {10}\end{center}}
\vspace{-6mm}

Написати функцію $f$, яка для файлу, заданого шляхом, розв’язує задачу:\\ Текстовий файл містить послідовність чисел $a_1$, $a_2$, $\ldots$, $a_t$, записану у форматі одне число на рядок. Знайти третє найбільше значення, що записане у файлі. Кожне значення, що записане у файлі, враховується стільки разів, скільки входить у файл. Для послідовності 2, 4, 12, 5, 2, 0 таким значенням буде 2.

В усіх випадках розміри файлів такі, що їх вміст повністю завантажений у пам'ять бути не може. Якість коду також оцінюється.



\vspace{10mm}\par

\newpage
\begin{minipage}[t]{9.7cm}
               
{\begin{center}Завдання {7}\end{center}}
\vspace{-6mm}

\begin{center}
	\adjustimage{max size={0.8\linewidth}{0.8\paperheight}}
	{./15BinTree/trees_exp/122.png}
\end{center}
Указати послідовність міток вершин дерева, зображеного на рис., що утвориться в результаті обходу дерева в прямому порядку. Мітки записувати через пробіл.\par Алгоритм починає роботу з кореня (на рис. це верхня вершина).\par

\end{minipage}
\vline \hspace{8mm}
\begin{minipage}[t]{8.0cm}
               
{\begin{center}Завдання {8}\end{center}}
\vspace{-6mm}
Для графа на рисунку з вершини 2 запущено алгоритм пошуку в глибину, що будує кістякове дерево. Списки суміжності вершин переглядаються алгоритмом за зростанням міток вершин.\par
\begin{center}
	\adjustimage{max size={0.9\linewidth}{0.9\paperheight}}
	{./16Graphs/Traversals/25.png}
\end{center}
\par Які з наведених ребер входять до складу отриманого кістякового дерева?\par
1) (1, 3)\par
2) (5, 6)\par
3) (0, 4)\par
4) (4, 6)\par
5) (3, 4)\par
6) (2, 3)\par
{\vskip 15pt} \par

\end{minipage}

               
{\begin{center}Завдання {9}\end{center}}
\vspace{-6mm}
Файл \code{'1.txt'} містить журнал з рядками, в яких через кому записано поряд\-ко\-вий номер запису (ціле), ім'я користувача (складається з латиниці), час події (фор\-мат як у прикладі), тип події (ERROR, WARNING, INFO, STATUS), код події (ціле).

Білі символи навколо роздільників не допускаються.

Приклад рядка журналу



\begin{verbatim}
12,username,2022-05-05 13:26,ERROR,404
\end{verbatim}


Нижче наведено код, який шукає усі події, що відбулись\\ для користувача \code{anonymous}, та зберігає їх типи в змінній \kw{codes}.



\begin{verbatim}
codes = set()
regex = re.compile(r'XXX')
with open('1.txt', encoding='utf8') as f:
    for s in f:
        res = re.fullmatch(regex, s)
        if not res: continue
        s = ''
        codes.add(YYY)
\end{verbatim}



Який рядок треба записати на місце \kw{XXX} ?
Який рядок треба записати на місце \kw{YYY} ?\par

%\end {large}
}
\newpage

{{23.05.2025 Контрольна робота, варіант  \No \textbf { 38 }} група К1{\parbox {0.5 cm} \dotfill}  ПІБ {\parbox {2.5 cm} \dotfill}\par}
{
%\begin {large}
\begin{minipage}[t]{9.7cm}
               
{\begin{center}Завдання {1}\end{center}}
\vspace{-6mm}
Укажіть ТИП значення виразу.\par\par
\begin{verbatim}
{x:0 for x in range(7)}
\end{verbatim}

\vspace{15mm}
У завданнях 2, 3 вказати, що буде виведено за виконання. Повідомлення інтерпретатора про помилку позначати ERROR

               
{\begin{center}Завдання {2}\end{center}}
\vspace{-6mm}

\begin{verbatim}
try:
    t = {28:2, 34:9, 44:2, 34:7}
    t[34] = 3
    for x in t.keys():
        print(x, sep=' ', end=' ')
except:
    print('ERR')
\end{verbatim}


\end{minipage}
\vline \hspace{8mm}
\begin{minipage}[t]{8.0cm}                
               
{\begin{center}Завдання {3}\end{center}}
\vspace{-6mm}

\begin{verbatim}
d = 0
m = 1

class A:
    m = 2
    
    def __init__(self):
        self.m = 3
        m = 4
        A.m = 5

obj = A()
try:
    print(d, end=' ')
    print(m, end=' ')
    print(A.m, end=' ')
    print(obj.m)
except:
    print('ERR')
\end{verbatim}

\end{minipage}

               
{\begin{center}Завдання {4}\end{center}}
\vspace{-6mm}
Нехай змінна \kw{s} позначає дуже довгу послідовність цілих, по якій можна ітерувати.

Написати вираз-генератор, що задає підпослідовність її елементів, що є від'єм\-ни\-ми.

Обмеження: усі елементи шуканої підпослідовності одночасно в пам'яті розміщені бути не можуть.\par

\vspace{10mm}
               
{\begin{center}Завдання {5}\end{center}}
\vspace{-6mm}
Змінні \kw{next} та \kw{prev} вузла списку позначають наступний та попередній за ним вузол відповідно. Починаючи з голови, у список записано послідовні цілі числа від~$-68$ до~$49$. Нехай змінна \kw{p} позначає вузол, що зберігає значення~$7$.
Яке число зберігається у вузлі \kw{p.next.prev.next.prev.prev} ?\par

               
{\begin{center}Завдання {6}\end{center}}
\vspace{-6mm}
Значенням змінної \kw{a} є цілочисловий масив типу $numpy.ndarray$ розмірів $3{\times}3$. На скільки збільшиться сума елементів масиву \kw{a} після виконання
\begin{verbatim}
a.T[-3] += 3
\end{verbatim}


Якщо код некоректний, то в якості відповіді зазначити ERROR.\par

               
{\begin{center}Завдання {10}\end{center}}
\vspace{-6mm}

Написати функцію $f$, яка для файлу, заданого шляхом, розв’язує задачу:\\ Текстовий файл містить послідовність чисел $a_1$, $a_2$, $\ldots$, $a_t$, записану у форматі одне число на рядок. Знайти середнє арифметичне чисел $a_1$, $a_4$, $a_7$, $\ldots$ (починаючи з першого через два).

В усіх випадках розміри файлів такі, що їх вміст повністю завантажений у пам'ять бути не може. Якість коду також оцінюється.



\vspace{10mm}\par

\newpage
\begin{minipage}[t]{9.7cm}
               
{\begin{center}Завдання {7}\end{center}}
\vspace{-6mm}

\begin{center}
	\adjustimage{max size={0.8\linewidth}{0.8\paperheight}}
	{./15BinTree/trees_exp/53.png}
\end{center}
Указати послідовність міток вершин дерева, зображеного на рис., що утвориться в результаті обходу дерева в оберненому порядку. Мітки записувати через пробіл.\par Алгоритм починає роботу з кореня (на рис. це верхня вершина).\par

\end{minipage}
\vline \hspace{8mm}
\begin{minipage}[t]{8.0cm}
               
{\begin{center}Завдання {8}\end{center}}
\vspace{-6mm}
Для графа на рисунку з вершини 4 запущено алгоритм пошуку в ширину, що будує кістякове дерево. Списки суміжності вершин переглядаються алгоритмом за зростанням міток вершин.\par
\begin{center}
	\adjustimage{max size={0.9\linewidth}{0.9\paperheight}}
	{./16Graphs/Traversals/149.png}
\end{center}
\par Які з наведених ребер входять до складу отриманого кістякового дерева?\par
1) (0, 5)\par
2) (5, 6)\par
3) (2, 4)\par
4) (0, 1)\par
5) (0, 4)\par
6) (0, 2)\par
{\vskip 15pt} \par

\end{minipage}

               
{\begin{center}Завдання {9}\end{center}}
\vspace{-6mm}
Файл \code{'1.txt'} містить журнал з рядками, в яких через двокрапку записано поряд\-ко\-вий номер запису (ціле), ім'я користувача (складається з латиниці), час події (фор\-мат як у прикладі), тип події (ERROR, WARNING, INFO, STATUS), код події (ціле).

Білі символи навколо роздільників не допускаються.

Приклад рядка журналу



\begin{verbatim}
12:username:2022-05-05 13-26:ERROR:404
\end{verbatim}


Нижче наведено код, який шукає усі події, що відбулись\\ для користувача \code{anonymous}, та зберігає їх типи в змінній \kw{codes}.



\begin{verbatim}
codes = set()
regex = re.compile(r'XXX')
with open('1.txt', encoding='utf8') as f:
    for s in f:
        res = re.fullmatch(regex, s)
        if not res: continue
        s = ''
        codes.add(YYY)
\end{verbatim}



Який рядок треба записати на місце \kw{XXX} ?
Який рядок треба записати на місце \kw{YYY} ?\par

%\end {large}
}
\newpage

{{23.05.2025 Контрольна робота, варіант  \No \textbf { 39 }} група К1{\parbox {0.5 cm} \dotfill}  ПІБ {\parbox {2.5 cm} \dotfill}\par}
{
%\begin {large}
\begin{minipage}[t]{9.7cm}
               
{\begin{center}Завдання {1}\end{center}}
\vspace{-6mm}
Укажіть ТИП значення виразу.\par\par
\begin{verbatim}
[7,2,3,4]
\end{verbatim}

\vspace{15mm}
У завданнях 2, 3 вказати, що буде виведено за виконання. Повідомлення інтерпретатора про помилку позначати ERROR

               
{\begin{center}Завдання {2}\end{center}}
\vspace{-6mm}

\begin{verbatim}
try:
    s = {12:8, 71:9, 22:6, 12:3}
    s[35] = 2
    for x in s.values():
        print(x, sep=' ', end=' ')
except:
    print('ERR')
\end{verbatim}


\end{minipage}
\vline \hspace{8mm}
\begin{minipage}[t]{8.0cm}                
               
{\begin{center}Завдання {3}\end{center}}
\vspace{-6mm}

\begin{verbatim}
a = 0
j = 1

class C:
    j = 2
    
    def __init__(self):
        global j
        self.j = 3
        j = 4
        C.j = 5

obj = C()
try:
    print(a, end=' ')
    print(j, end=' ')
    print(C.j, end=' ')
    print(obj.j)
except:
    print('ERR')
\end{verbatim}

\end{minipage}

               
{\begin{center}Завдання {4}\end{center}}
\vspace{-6mm}

 Задано числову послідовність \kw{s} довжини 6000. Записати ВИРАЗ, значенням якого є список, що складається з другої третини послідовності, записаної в оберненому порядку.\par

\vspace{10mm}
               
{\begin{center}Завдання {5}\end{center}}
\vspace{-6mm}
Змінні \kw{next} та \kw{prev} вузла списку позначають наступний та попередній за ним вузол відповідно. Починаючи з голови, у список записано послідовні цілі числа від~$-34$ до~$27$. Нехай змінна \kw{p} позначає вузол, що зберігає значення~$-7$.
Яке число зберігається у вузлі \kw{p.next.next.prev.prev.next} ?\par

               
{\begin{center}Завдання {6}\end{center}}
\vspace{-6mm}
Значенням змінної \kw{a} є цілочисловий масив типу $numpy.ndarray$ розмірів $7{\times}7$. На скільки збільшиться сума елементів масиву \kw{a} після виконання
\begin{verbatim}
a.T[2] += 1
\end{verbatim}


Якщо код некоректний, то в якості відповіді зазначити ERROR.\par

               
{\begin{center}Завдання {10}\end{center}}
\vspace{-6mm}

Написати функцію $f$, яка для файлу, заданого шляхом, розв’язує задачу:\\ Текстовий файл містить послідовність чисел $a_1$, $a_2$, $\ldots$, $a_t$, записану у форматі одне число на рядок. Знайти третє найменше значення, що записане у файлі. Кожне значення, що записане у файлі, враховується тільки один раз. \par
Для послідовності 1, 4, 12, 5, 1, 0 таким значенням буде 4.

В усіх випадках розміри файлів такі, що їх вміст повністю завантажений у пам'ять бути не може. Якість коду також оцінюється.



\vspace{10mm}\par

\newpage
\begin{minipage}[t]{9.7cm}
               
{\begin{center}Завдання {7}\end{center}}
\vspace{-6mm}

\begin{center}
	\adjustimage{max size={0.8\linewidth}{0.8\paperheight}}
	{./15BinTree/trees_exp/15.png}
\end{center}
Указати послідовність міток вершин дерева, зображеного на рис., що утвориться в результаті обходу дерева в прямому порядку. Мітки записувати через пробіл.\par Алгоритм починає роботу з кореня (на рис. це верхня вершина).\par

\end{minipage}
\vline \hspace{8mm}
\begin{minipage}[t]{8.0cm}
               
{\begin{center}Завдання {8}\end{center}}
\vspace{-6mm}
Для графа на рисунку з вершини 0 запущено алгоритм пошуку в глибину, що будує кістякове дерево. Списки суміжності вершин переглядаються алгоритмом за зростанням міток вершин.\par
\begin{center}
	\adjustimage{max size={0.9\linewidth}{0.9\paperheight}}
	{./16Graphs/Traversals/77.png}
\end{center}
\par Які з наведених ребер входять до складу отриманого кістякового дерева?\par
1) (2, 3)\par
2) (0, 2)\par
3) (0, 6)\par
4) (0, 1)\par
5) (0, 5)\par
6) (1, 6)\par
{\vskip 15pt} \par

\end{minipage}

               
{\begin{center}Завдання {9}\end{center}}
\vspace{-6mm}
Файл \code{'1.txt'} містить відомість з рядками, в яких через двокрапку записа\-но назву предмету (знаків пунктуації не містить), прізвище студента, ім'я студента, номер залі\-ков\-ки, оцінку в балах за 100 бальною шкалою.

Білі символи навколо роздільників не допускаються.

Приклад рядка файлу



\begin{verbatim}
algebra:surname:name:123456:77
\end{verbatim}


Нижче наведено код, який має замінити у файлі усі номери заліковок на рядок \code{***}.



\begin{verbatim}
regex = re.compile(r'XXX')
with open('1.txt', encoding='utf8') as f:
    lst = f.readlines()
for i in range(len(lst)):
    lst[i] = re.sub(regex, YYY, lst[i])
with open('1.txt', 'w', encoding='utf8') as f:
    f.writelines(lst)
\end{verbatim}





Який рядок треба записати на місце \kw{XXX} ?
Який рядок треба записати на місце \kw{YYY} ?\par

%\end {large}
}
\newpage

{{23.05.2025 Контрольна робота, варіант  \No \textbf { 40 }} група К1{\parbox {0.5 cm} \dotfill}  ПІБ {\parbox {2.5 cm} \dotfill}\par}
{
%\begin {large}
\begin{minipage}[t]{9.7cm}
               
{\begin{center}Завдання {1}\end{center}}
\vspace{-6mm}
Укажіть ТИП значення виразу.\par\par
\begin{verbatim}
{x for x in range(7)}
\end{verbatim}

\vspace{15mm}
У завданнях 2, 3 вказати, що буде виведено за виконання. Повідомлення інтерпретатора про помилку позначати ERROR

               
{\begin{center}Завдання {2}\end{center}}
\vspace{-6mm}

\begin{verbatim}
try:
    t = {10:0, 48:0, 64:4, 35:3, 10:4}
    t[76] = 9
    for x, y in t.items():
        print(x, y, sep=' ', end=' ')
except:
    print('ERR')
\end{verbatim}


\end{minipage}
\vline \hspace{8mm}
\begin{minipage}[t]{8.0cm}                
               
{\begin{center}Завдання {3}\end{center}}
\vspace{-6mm}

\begin{verbatim}
_c = 0

class D:
    _c = 1
    
    def __init__(self):
        self._c = 2
        _c = 3
        D._c = 4

obj = D()
obj._c = 5
try:
    print(_c, end=' ')
    print(obj._c, end=' ')
    print(D._c, end=' ')
    print(obj._c, end=' ')
    print(obj._D__c)
except:
    print('ERR')
\end{verbatim}

\end{minipage}

               
{\begin{center}Завдання {4}\end{center}}
\vspace{-6mm}
Нехай змінна \kw{s} позначає дуже довгу послідовність, по якій можна ітерувати.

Написати вираз-генератор, що задає підпослідовність її елементів, що є числами типу \kw{int}.

Обмеження: усі елементи шуканої підпослідовності одночасно в пам'яті розміщені бути не можуть.\par

\vspace{10mm}
               
{\begin{center}Завдання {5}\end{center}}
\vspace{-6mm}
Змінні \kw{next} та \kw{prev} вузла списку позначають наступний та попередній за ним вузол відповідно. Починаючи з голови, у список записано послідовні цілі числа від~$-43$ до~$45$. Нехай змінна \kw{p} позначає вузол, що зберігає значення~$-6$.
Яке число зберігається у вузлі \kw{p.prev.prev.next.next.prev} ?\par

               
{\begin{center}Завдання {6}\end{center}}
\vspace{-6mm}
Значенням змінної \kw{a} є цілочисловий масив типу $numpy.ndarray$ розмірів $6{\times}5$. На скільки збільшиться сума елементів масиву \kw{a} після виконання
\begin{verbatim}
a[0] += 1
\end{verbatim}


Якщо код некоректний, то в якості відповіді зазначити ERROR.\par

               
{\begin{center}Завдання {10}\end{center}}
\vspace{-6mm}

Написати функцію $f$, яка для файлу, заданого шляхом, розв’язує задачу:\\ Текстовий файл містить послідовність чисел $a_1$, $a_2$, $\ldots$, $a_t$, записану у форматі одне число на рядок. Знайти третє найбільше значення, що записане у файлі. Кожне значення, що записане у файлі, враховується тільки один раз. Для послідовності 1, 4, 12, 5, 1, 0 таким значенням буде 4.

В усіх випадках розміри файлів такі, що їх вміст повністю завантажений у пам'ять бути не може. Якість коду також оцінюється.



\vspace{10mm}\par

\newpage
\begin{minipage}[t]{9.7cm}
               
{\begin{center}Завдання {7}\end{center}}
\vspace{-6mm}

\begin{center}
	\adjustimage{max size={0.8\linewidth}{0.8\paperheight}}
	{./15BinTree/trees_exp/54.png}
\end{center}
Указати послідовність міток вершин дерева, зображеного на рис., що утвориться в результаті обходу дерева в оберненому порядку. Мітки записувати через пробіл.\par Алгоритм починає роботу з кореня (на рис. це верхня вершина).\par

\end{minipage}
\vline \hspace{8mm}
\begin{minipage}[t]{8.0cm}
               
{\begin{center}Завдання {8}\end{center}}
\vspace{-6mm}
Для графа на рисунку з вершини 1 запущено алгоритм пошуку в глибину, що будує кістякове дерево. Списки суміжності вершин переглядаються алгоритмом за зростанням міток вершин.\par
\begin{center}
	\adjustimage{max size={0.9\linewidth}{0.9\paperheight}}
	{./16Graphs/Traversals/111.png}
\end{center}
\par Які з наведених ребер входять до складу отриманого кістякового дерева?\par
1) (0, 2)\par
2) (0, 1)\par
3) (5, 6)\par
4) (1, 6)\par
5) (2, 6)\par
6) (2, 4)\par
{\vskip 15pt} \par

\end{minipage}

               
{\begin{center}Завдання {9}\end{center}}
\vspace{-6mm}
Файл \code{'1.txt'} містить журнал з рядками, в яких через кому записано поряд\-ко\-вий номер запису в журналі (ціле), тип події (ERROR, WARNING, INFO, STATUS), ім'я користувача (складається з латиниці), час події (фор\-мат як у прикладі), код події (ціле).

Білі символи навколо роздільників не допускаються.

Приклад рядка файлу



\begin{verbatim}
13,ERROR,username,2022-05-05 13:26,404
\end{verbatim}


Нижче наведено код, який має замінити у файлі імена користувачів на рядок \code{unknown}.



\begin{verbatim}
regex = re.compile(r'XXX')
with open('1.txt', encoding='utf8') as f:
    lst = f.readlines()
for i in range(len(lst)):
    lst[i] = re.sub(regex, YYY, lst[i])
with open('1.txt', 'w', encoding='utf8') as f:
    f.writelines(lst)
\end{verbatim}





Який рядок треба записати на місце \kw{XXX} ?
Який рядок треба записати на місце \kw{YYY} ?\par

%\end {large}
}
\newpage

{{23.05.2025 Контрольна робота, варіант  \No \textbf { 41 }} група К1{\parbox {0.5 cm} \dotfill}  ПІБ {\parbox {2.5 cm} \dotfill}\par}
{
%\begin {large}
\begin{minipage}[t]{9.7cm}
               
{\begin{center}Завдання {1}\end{center}}
\vspace{-6mm}
Укажіть ТИП значення виразу.\par\par
\begin{verbatim}
[13,2,4]+[3,5]
\end{verbatim}

\vspace{15mm}
У завданнях 2, 3 вказати, що буде виведено за виконання. Повідомлення інтерпретатора про помилку позначати ERROR

               
{\begin{center}Завдання {2}\end{center}}
\vspace{-6mm}

\begin{verbatim}
try:
    d = {47:1, 20:1, 87:1, 87:3}
    d[50] = 8
    for x in d.items():
        print(*x, sep=' ', end=' ')
except:
    print('ERR')
\end{verbatim}


\end{minipage}
\vline \hspace{8mm}
\begin{minipage}[t]{8.0cm}                
               
{\begin{center}Завдання {3}\end{center}}
\vspace{-6mm}

\begin{verbatim}
b = 0
n = 1

class D:
    b = 2
    
    def __init__(self):
        self.b = 3
        b = 4
        D.b = 5

obj = D()
try:
    print(b, end=' ')
    print(n, end=' ')
    print(D.b, end=' ')
    print(obj.b)
except:
    print('ERR')
\end{verbatim}

\end{minipage}

               
{\begin{center}Завдання {4}\end{center}}
\vspace{-6mm}
Нехай змінна \kw{s} позначає послідовність цілих, по якій можна ітерувати.

Написати вираз, що перевіряє, чи правда, що всі елементи послідовності менші за задане число num.\par

\vspace{10mm}
               
{\begin{center}Завдання {5}\end{center}}
\vspace{-6mm}
Змінні \kw{next} та \kw{prev} вузла списку позначають наступний та попередній за ним вузол відповідно. Починаючи з голови, у список записано послідовні цілі числа від~$-62$ до~$31$. Нехай змінна \kw{p} позначає вузол, що зберігає значення~$8$.
Яке число зберігається у вузлі \kw{p.next.prev.prev.next.prev} ?\par

               
{\begin{center}Завдання {6}\end{center}}
\vspace{-6mm}
Значенням змінної \kw{a} є цілочисловий масив типу $numpy.ndarray$ розмірів $7{\times}5$. На скільки збільшиться сума елементів масиву \kw{a} після виконання
\begin{verbatim}
a.T[-2, 0] += 2
\end{verbatim}


Якщо код некоректний, то в якості відповіді зазначити ERROR.\par

               
{\begin{center}Завдання {10}\end{center}}
\vspace{-6mm}

Написати функцію $f$, яка для файлу, заданого шляхом, розв’язує задачу:\\ Текстовий файл містить послідовність чисел $a_1$, $a_2$, $\ldots$, $a_t$, записану у форматі одне число на рядок. Знайти третє найменше значення, що записане у файлі. Кожне значення, що записане у файлі, враховується стільки разів, скільки входить у файл. Для послідовності 2, 4, 12, 5, 2, 0 таким значенням буде 2.

В усіх випадках розміри файлів такі, що їх вміст повністю завантажений у пам'ять бути не може. Якість коду також оцінюється.



\vspace{10mm}\par

\newpage
\begin{minipage}[t]{9.7cm}
               
{\begin{center}Завдання {7}\end{center}}
\vspace{-6mm}

\begin{center}
	\adjustimage{max size={0.8\linewidth}{0.8\paperheight}}
	{./15BinTree/trees_exp/127.png}
\end{center}
Указати послідовність міток вершин дерева, зображеного на рис., що утвориться в результаті обходу дерева в прямому порядку. Мітки записувати через пробіл.\par Алгоритм починає роботу з кореня (на рис. це верхня вершина).\par

\end{minipage}
\vline \hspace{8mm}
\begin{minipage}[t]{8.0cm}
               
{\begin{center}Завдання {8}\end{center}}
\vspace{-6mm}
Для графа на рисунку з вершини 4 запущено алгоритм пошуку в ширину, що будує кістякове дерево. Списки суміжності вершин переглядаються алгоритмом за зростанням міток вершин.\par
\begin{center}
	\adjustimage{max size={0.9\linewidth}{0.9\paperheight}}
	{./16Graphs/Traversals/311.png}
\end{center}
\par Які з наведених ребер входять до складу отриманого кістякового дерева?\par
1) (3, 5)\par
2) (4, 6)\par
3) (5, 6)\par
4) (2, 3)\par
5) (1, 4)\par
6) (3, 4)\par
{\vskip 15pt} \par

\end{minipage}

               
{\begin{center}Завдання {9}\end{center}}
\vspace{-6mm}
Файл \code{'1.txt'} містить відомість з рядками, в яких через двокрапку записано поряд\-ковий номер (ціле), назву предмету (знаків пунктуації не містить), прізвище студента, ім'я студента, номер заліковки, оцінку в балах за 100 бальною шкалою.

Білі символи навколо роздільників не допускаються.

Приклад рядка файлу



\begin{verbatim}
13:algebra:surname:name:123456:77
\end{verbatim}


Нижче наведено код, який шукає усі назви предметів, що зустрічаються у файлі, та зберігає їх у змінній \id{codes}.



\begin{verbatim}
codes = set()
regex = re.compile(r'XXX')
with open('1.txt', encoding='utf8') as f:
    for s in f:
        res = re.fullmatch(regex, s)
        if not res: continue
        s = ''
        codes.add(YYY)
\end{verbatim}



Який рядок треба записати на місце \kw{XXX} ?
Який рядок треба записати на місце \kw{YYY} ?\par

%\end {large}
}
\newpage

{{23.05.2025 Контрольна робота, варіант  \No \textbf { 42 }} група К1{\parbox {0.5 cm} \dotfill}  ПІБ {\parbox {2.5 cm} \dotfill}\par}
{
%\begin {large}
\begin{minipage}[t]{9.7cm}
               
{\begin{center}Завдання {1}\end{center}}
\vspace{-6mm}
Укажіть ТИП значення виразу.\par\par
\begin{verbatim}
[{1,2}, 5]
\end{verbatim}

\vspace{15mm}
У завданнях 2, 3 вказати, що буде виведено за виконання. Повідомлення інтерпретатора про помилку позначати ERROR

               
{\begin{center}Завдання {2}\end{center}}
\vspace{-6mm}

\begin{verbatim}
try:
    s = {84:7, 25:8, 51:0, 47:5}
    s[84] = 8
    for x in s :
        print(x, sep=' ', end=' ')
except:
    print('ERR')
\end{verbatim}


\end{minipage}
\vline \hspace{8mm}
\begin{minipage}[t]{8.0cm}                
               
{\begin{center}Завдання {3}\end{center}}
\vspace{-6mm}

\begin{verbatim}
b = 0

class C:
    b = 1
    
    def __init__(self):
        self.b = 2
        b = 3
        C.b = 4

obj = C()
C.b = 5
try:
    print(b, end=' ')
    print(obj.b, end=' ')
    print(C.b, end=' ')
    print(obj.b, end=' ')
    print(obj._C__b)
except:
    print('ERR')
\end{verbatim}

\end{minipage}

               
{\begin{center}Завдання {4}\end{center}}
\vspace{-6mm}
Нехай змінна \kw{s} позначає дуже довгу послідовність цілих, по якій можна ітерувати.

Написати вираз-генератор, що задає підпослідовність її елементів, що не діляться на жодне з чисел 2, 7, 5.

Обмеження: усі елементи шуканої підпослідовності одночасно в пам'яті розміщені бути не можуть.\par

\vspace{10mm}
               
{\begin{center}Завдання {5}\end{center}}
\vspace{-6mm}
Змінні \kw{next} та \kw{prev} вузла списку позначають наступний та попередній за ним вузол відповідно. Починаючи з голови, у список записано послідовні цілі числа від~$-56$ до~$44$. Нехай змінна \kw{p} позначає вузол, що зберігає значення~$2$.
Яке число зберігається у вузлі \kw{p.prev.next.next.prev.next} ?\par

               
{\begin{center}Завдання {6}\end{center}}
\vspace{-6mm}
Значенням змінної \kw{a} є цілочисловий масив типу $numpy.ndarray$ розмірів $4{\times}6$. На скільки збільшиться сума елементів масиву \kw{a} після виконання
\begin{verbatim}
a[-1] += 3
\end{verbatim}


Якщо код некоректний, то в якості відповіді зазначити ERROR.\par

               
{\begin{center}Завдання {10}\end{center}}
\vspace{-6mm}

Написати функцію $f$, яка для файлу, заданого шляхом, розв’язує задачу:\\ Кожен з двох текстових файлів містить послідовність чисел, записану у форматі одне число на рядок.  Перевірити, чи задають вони одну ту саму послідовність.

В усіх випадках розміри файлів такі, що їх вміст повністю завантажений у пам'ять бути не може. Якість коду також оцінюється.



\vspace{10mm}\par

\newpage
\begin{minipage}[t]{9.7cm}
               
{\begin{center}Завдання {7}\end{center}}
\vspace{-6mm}

\begin{center}
	\adjustimage{max size={0.8\linewidth}{0.8\paperheight}}
	{./15BinTree/trees_exp/98.png}
\end{center}
Указати послідовність міток вершин дерева, зображеного на рис., що утвориться в результаті обходу дерева в оберненому порядку. Мітки записувати через пробіл.\par Алгоритм починає роботу з кореня (на рис. це верхня вершина).\par

\end{minipage}
\vline \hspace{8mm}
\begin{minipage}[t]{8.0cm}
               
{\begin{center}Завдання {8}\end{center}}
\vspace{-6mm}
Для графа на рисунку з вершини 3 запущено алгоритм пошуку в глибину, що будує кістякове дерево. Списки суміжності вершин переглядаються алгоритмом за зростанням міток вершин.\par
\begin{center}
	\adjustimage{max size={0.9\linewidth}{0.9\paperheight}}
	{./16Graphs/Traversals/298.png}
\end{center}
\par Які з наведених ребер входять до складу отриманого кістякового дерева?\par
1) (3, 5)\par
2) (2, 6)\par
3) (0, 5)\par
4) (0, 2)\par
5) (1, 2)\par
6) (3, 4)\par
{\vskip 15pt} \par

\end{minipage}

               
{\begin{center}Завдання {9}\end{center}}
\vspace{-6mm}
Файл \code{'1.txt'} містить журнал з рядками, в яких через кому записано поряд\-ко\-вий номер запису (ціле), ім'я користувача (складається з латиниці), час події (фор\-мат як у прикладі), тип події (ERROR, WARNING, INFO, STATUS), код події (ціле).

Білі символи навколо роздільників не допускаються.

Приклад рядка журналу



\begin{verbatim}
12,username,2022-05-05 13:26,ERROR,404
\end{verbatim}


Нижче наведено код, який шукає усі попередження (подія WARNING), що зустрі\-ча\-ються в журналі і відбулися під час обідньої перерви деканату (з 13:00 по 13:59 включно), та зберігає їх коди в змінній \kw{codes}.



\begin{verbatim}
codes = set()
regex = re.compile(r'XXX')
with open('1.txt', encoding='utf8') as f:
    for s in f:
        res = re.fullmatch(regex, s)
        if not res: continue
        s = ''
        codes.add(YYY)
\end{verbatim}





Який рядок треба записати на місце \kw{XXX} ?
Який рядок треба записати на місце \kw{YYY} ?\par

%\end {large}
}
\newpage

{{23.05.2025 Контрольна робота, варіант  \No \textbf { 43 }} група К1{\parbox {0.5 cm} \dotfill}  ПІБ {\parbox {2.5 cm} \dotfill}\par}
{
%\begin {large}
\begin{minipage}[t]{9.7cm}
               
{\begin{center}Завдання {1}\end{center}}
\vspace{-6mm}
Укажіть ТИП значення виразу.\par\par
\begin{verbatim}
{x*x:x for x in range(-7,7)}
\end{verbatim}

\vspace{15mm}
У завданнях 2, 3 вказати, що буде виведено за виконання. Повідомлення інтерпретатора про помилку позначати ERROR

               
{\begin{center}Завдання {2}\end{center}}
\vspace{-6mm}

\begin{verbatim}
try:
    t = {45:7, 13:9, 27:9, 41:3, 13:0}
    t[22] = 3
    for x in t.values():
        print(x, sep=' ', end=' ')
except:
    print('ERR')
\end{verbatim}


\end{minipage}
\vline \hspace{8mm}
\begin{minipage}[t]{8.0cm}                
               
{\begin{center}Завдання {3}\end{center}}
\vspace{-6mm}

\begin{verbatim}
__c = 0

class C:
    __c = 1
    
    def __init__(self):
        self.__c = 2
        __c = 3
        C.__c = 4

obj = C()
C.__c = 5
try:
    print(__c, end=' ')
    print(obj.__c, end=' ')
    print(C.__c, end=' ')
    print(obj.__c, end=' ')
    print(obj._C__c)
except:
    print('ERR')
\end{verbatim}

\end{minipage}

               
{\begin{center}Завдання {4}\end{center}}
\vspace{-6mm}
Задано словник \kw{d} з цілими ключами та значеннями. Записати ВИРАЗ, значенням якого є множина, що складається  з тих значень, чиї ключі діляться на 4.\par

\vspace{10mm}
               
{\begin{center}Завдання {5}\end{center}}
\vspace{-6mm}
Змінні \kw{next} та \kw{prev} вузла списку позначають наступний та попередній за ним вузол відповідно. Починаючи з голови, у список записано послідовні цілі числа від~$-64$ до~$34$. Нехай змінна \kw{p} позначає вузол, що зберігає значення~$3$.
Яке число зберігається у вузлі \kw{p.prev.prev.prev.prev.prev} ?\par

               
{\begin{center}Завдання {6}\end{center}}
\vspace{-6mm}
Значенням змінної \kw{a} є цілочисловий масив типу $numpy.ndarray$ розмірів $5{\times}3$. На скільки збільшиться сума елементів масиву \kw{a} після виконання
\begin{verbatim}
a.T[-3] += 3
\end{verbatim}


Якщо код некоректний, то в якості відповіді зазначити ERROR.\par

               
{\begin{center}Завдання {10}\end{center}}
\vspace{-6mm}

Написати функцію $f$, яка для файлу, заданого шляхом, розв’язує задачу:\\ Кожен з двох текстових файлів містить послідовність чисел, записану у форматі одне число на рядок.  Перевірити, чи задають вони одну ту саму послідовність.

В усіх випадках розміри файлів такі, що їх вміст повністю завантажений у пам'ять бути не може. Якість коду також оцінюється.



\vspace{10mm}\par

\newpage
\begin{minipage}[t]{9.7cm}
               
{\begin{center}Завдання {7}\end{center}}
\vspace{-6mm}

\begin{center}
	\adjustimage{max size={0.8\linewidth}{0.8\paperheight}}
	{./15BinTree/trees_exp/21.png}
\end{center}
Указати послідовність міток вершин дерева, зображеного на рис., що утвориться в результаті обходу дерева в оберненому порядку. Мітки записувати через пробіл.\par Алгоритм починає роботу з кореня (на рис. це верхня вершина).\par

\end{minipage}
\vline \hspace{8mm}
\begin{minipage}[t]{8.0cm}
               
{\begin{center}Завдання {8}\end{center}}
\vspace{-6mm}
Для графа на рисунку з вершини 1 запущено алгоритм пошуку в ширину, що будує кістякове дерево. Списки суміжності вершин переглядаються алгоритмом за зростанням міток вершин.\par
\begin{center}
	\adjustimage{max size={0.9\linewidth}{0.9\paperheight}}
	{./16Graphs/Traversals/112.png}
\end{center}
\par Які з наведених ребер входять до складу отриманого кістякового дерева?\par
1) (1, 2)\par
2) (0, 6)\par
3) (0, 1)\par
4) (2, 6)\par
5) (3, 6)\par
6) (5, 6)\par
{\vskip 15pt} \par

\end{minipage}

               
{\begin{center}Завдання {9}\end{center}}
\vspace{-6mm}
Файл \code{'1.txt'} містить журнал з рядками, в яких через двокрапку записано поряд\-ко\-вий номер запису в журналі (ціле), код події (ціле), тип події (ERROR, \\WARNING, INFO, STATUS), ім'я користувача (складається з латиниці), час події (фор\-мат як у прикладі), опис події.

Білі символи навколо роздільників не допускаються.

Приклад рядка файлу



\begin{verbatim}
13:404:ERROR:username:2022-05-05 13-26:something happens
\end{verbatim}


Нижче наведено код, який має залишити у файлі в кожному рядку тільки код події та ім'я користувача. Рядок з прикладу має бути замінений на такий



\begin{verbatim}
404:username
\end{verbatim}




Код:



\begin{verbatim}
regex = re.compile(r'XXX')
with open('1.txt', encoding='utf8') as f:
    lst = f.readlines()
for i in range(len(lst)):
    lst[i] = re.sub(regex, YYY, lst[i])
with open('1.txt', 'w', encoding='utf8') as f:
    f.writelines(lst)
\end{verbatim}





Який рядок треба записати на місце \kw{XXX} ?
Який рядок треба записати на місце \kw{YYY} ?\par

%\end {large}
}
\newpage

{{23.05.2025 Контрольна робота, варіант  \No \textbf { 44 }} група К1{\parbox {0.5 cm} \dotfill}  ПІБ {\parbox {2.5 cm} \dotfill}\par}
{
%\begin {large}
\begin{minipage}[t]{9.7cm}
               
{\begin{center}Завдання {1}\end{center}}
\vspace{-6mm}
Укажіть ТИП значення виразу.\par\par
\begin{verbatim}
{y:None for y in range(-8,3)}
\end{verbatim}

\vspace{15mm}
У завданнях 2, 3 вказати, що буде виведено за виконання. Повідомлення інтерпретатора про помилку позначати ERROR

               
{\begin{center}Завдання {2}\end{center}}
\vspace{-6mm}

\begin{verbatim}
try:
    d = {37:4, 39:7, 75:1, 72:8, 72:6}
    d[77] = 4
    for x in d.keys():
        print(x, sep=' ', end=' ')
except:
    print('ERR')
\end{verbatim}


\end{minipage}
\vline \hspace{8mm}
\begin{minipage}[t]{8.0cm}                
               
{\begin{center}Завдання {3}\end{center}}
\vspace{-6mm}

\begin{verbatim}
_a = 0

class B:
    _a = 1
    
    def __init__(self):
        self._a = 2
        _a = 3
        B._a = 4

obj = B()
obj._a = 5
try:
    print(_a, end=' ')
    print(obj._a, end=' ')
    print(B._a, end=' ')
    print(obj._a, end=' ')
    print(obj._B__a)
except:
    print('ERR')
\end{verbatim}

\end{minipage}

               
{\begin{center}Завдання {4}\end{center}}
\vspace{-6mm}
Нехай змінна \kw{s} позначає послідовність цілих, по якій можна ітерувати.

Написати вираз, що перевіряє, чи правда, що всі елементи послідовності є від’ємни\-ми.\par

\vspace{10mm}
               
{\begin{center}Завдання {5}\end{center}}
\vspace{-6mm}
Змінні \kw{next} та \kw{prev} вузла списку позначають наступний та попередній за ним вузол відповідно. Починаючи з голови, у список записано послідовні цілі числа від~$-58$ до~$33$. Нехай змінна \kw{p} позначає вузол, що зберігає значення~$5$.
Яке число зберігається у вузлі \kw{p.next.next.next.next.next} ?\par

               
{\begin{center}Завдання {6}\end{center}}
\vspace{-6mm}
Значенням змінної \kw{a} є цілочисловий масив типу $numpy.ndarray$ розмірів $6{\times}4$. На скільки збільшиться сума елементів масиву \kw{a} після виконання
\begin{verbatim}
a[1] += 1
\end{verbatim}


Якщо код некоректний, то в якості відповіді зазначити ERROR.\par

               
{\begin{center}Завдання {10}\end{center}}
\vspace{-6mm}

Написати функцію $f$, яка для файлу, заданого шляхом, розв’язує задачу:\\ Текстовий файл містить послідовність чисел $a_1$, $a_2$, $\ldots$, $a_t$, записану у форматі одне число на рядок. Знайти третє найменше значення, що записане у файлі. Кожне значення, що записане у файлі, враховується тільки один раз. \par
Для послідовності 1, 4, 12, 5, 1, 0 таким значенням буде 4.

В усіх випадках розміри файлів такі, що їх вміст повністю завантажений у пам'ять бути не може. Якість коду також оцінюється.



\vspace{10mm}\par

\newpage
\begin{minipage}[t]{9.7cm}
               
{\begin{center}Завдання {7}\end{center}}
\vspace{-6mm}

\begin{center}
	\adjustimage{max size={0.8\linewidth}{0.8\paperheight}}
	{./15BinTree/trees_exp/167.png}
\end{center}
Указати послідовність міток вершин дерева, зображеного на рис., що утвориться в результаті обходу дерева в прямому порядку. Мітки записувати через пробіл.\par Алгоритм починає роботу з кореня (на рис. це верхня вершина).\par

\end{minipage}
\vline \hspace{8mm}
\begin{minipage}[t]{8.0cm}
               
{\begin{center}Завдання {8}\end{center}}
\vspace{-6mm}
Для графа на рисунку з вершини 2 запущено алгоритм пошуку в глибину, що будує кістякове дерево. Списки суміжності вершин переглядаються алгоритмом за зростанням міток вершин.\par
\begin{center}
	\adjustimage{max size={0.9\linewidth}{0.9\paperheight}}
	{./16Graphs/Traversals/248.png}
\end{center}
\par Які з наведених ребер входять до складу отриманого кістякового дерева?\par
1) (0, 4)\par
2) (3, 6)\par
3) (4, 5)\par
4) (0, 6)\par
5) (0, 1)\par
6) (2, 6)\par
{\vskip 15pt} \par

\end{minipage}

               
{\begin{center}Завдання {9}\end{center}}
\vspace{-6mm}
Файл \code{'1.txt'} містить журнал з рядками, в яких через двокрапку записано поряд\-ко\-вий номер запису в журналі (ціле), код події (ціле), тип події (ERROR,\\ WARNING, INFO, STATUS), ім'я користувача (складається з латиниці), час події (фор\-мат як у прикладі), опис події.

Білі символи навколо роздільників не допускаються.

Приклад рядка файлу



\begin{verbatim}
13:404:ERROR:username:2022-05-05 13:26:something happens
\end{verbatim}


Нижче наведено код, який має в кожному рядку файли переставити місцями код та тип події.

Рядок з прикладу має бути замінений на такий



\begin{verbatim}
13:ERROR:404:username:2022-05-05 13:26:something happens
\end{verbatim}




Код:



\begin{verbatim}
regex = re.compile(r'XXX')
with open('1.txt', encoding='utf8') as f:
    lst = f.readlines()
for i in range(len(lst)):
    lst[i] = re.sub(regex, YYY, lst[i])
with open('1.txt', 'w', encoding='utf8') as f:
    f.writelines(lst)
\end{verbatim}





Який рядок треба записати на місце \kw{XXX} ?
Який рядок треба записати на місце \kw{YYY} ?\par

%\end {large}
}
\newpage

{{23.05.2025 Контрольна робота, варіант  \No \textbf { 45 }} група К1{\parbox {0.5 cm} \dotfill}  ПІБ {\parbox {2.5 cm} \dotfill}\par}
{
%\begin {large}
\begin{minipage}[t]{9.7cm}
               
{\begin{center}Завдання {1}\end{center}}
\vspace{-6mm}
Укажіть ТИП значення виразу.\par\par
\begin{verbatim}
{x for x in range(7) if x%3==1}
\end{verbatim}

\vspace{15mm}
У завданнях 2, 3 вказати, що буде виведено за виконання. Повідомлення інтерпретатора про помилку позначати ERROR

               
{\begin{center}Завдання {2}\end{center}}
\vspace{-6mm}

\begin{verbatim}
try:
    s = {11:2, 88:2, 40:5, 11:5}
    s[11] = 4
    for x in s :
        print(x, sep=' ', end=' ')
except:
    print('ERR')
\end{verbatim}


\end{minipage}
\vline \hspace{8mm}
\begin{minipage}[t]{8.0cm}                
               
{\begin{center}Завдання {3}\end{center}}
\vspace{-6mm}

\begin{verbatim}
d = 0
k = 1

class A:
    d = 2
    
    def __init__(self):
        global d
        self.d = 3
        d = 4
        A.d = 5

obj = A()
try:
    print(d, end=' ')
    print(k, end=' ')
    print(A.d, end=' ')
    print(obj.d)
except:
    print('ERR')
\end{verbatim}

\end{minipage}

               
{\begin{center}Завдання {4}\end{center}}
\vspace{-6mm}
Нехай змінна \kw{s} позначає послідовність цілих, по якій можна ітерувати.

Написати вираз, що перевіряє, чи правда, що всі елементи послідовності діляться на 11.\par

\vspace{10mm}
               
{\begin{center}Завдання {5}\end{center}}
\vspace{-6mm}
Змінні \kw{next} та \kw{prev} вузла списку позначають наступний та попередній за ним вузол відповідно. Починаючи з голови, у список записано послідовні цілі числа від~$-42$ до~$39$. Нехай змінна \kw{p} позначає вузол, що зберігає значення~$6$.
Яке число зберігається у вузлі \kw{p.next.next.prev.next.next} ?\par

               
{\begin{center}Завдання {6}\end{center}}
\vspace{-6mm}
Значенням змінної \kw{a} є цілочисловий масив типу $numpy.ndarray$ розмірів $3{\times}7$. На скільки збільшиться сума елементів масиву \kw{a} після виконання
\begin{verbatim}
a.T[2] += 2
\end{verbatim}


Якщо код некоректний, то в якості відповіді зазначити ERROR.\par

               
{\begin{center}Завдання {10}\end{center}}
\vspace{-6mm}

Написати функцію $f$, яка для файлу, заданого шляхом, розв’язує задачу:\\ Текстовий файл містить послідовність чисел $a_1$, $a_2$, $\ldots$, $a_t$, записану у форматі одне число на рядок. Знайти третє найбільше значення, що записане у файлі. Кожне значення, що записане у файлі, враховується стільки разів, скільки входить у файл. Для послідовності 2, 4, 12, 5, 2, 0 таким значенням буде 2.

В усіх випадках розміри файлів такі, що їх вміст повністю завантажений у пам'ять бути не може. Якість коду також оцінюється.



\vspace{10mm}\par

\newpage
\begin{minipage}[t]{9.7cm}
               
{\begin{center}Завдання {7}\end{center}}
\vspace{-6mm}

\begin{center}
	\adjustimage{max size={0.8\linewidth}{0.8\paperheight}}
	{./15BinTree/trees_exp/78.png}
\end{center}
Указати послідовність міток вершин дерева, зображеного на рис., що утвориться в результаті обходу дерева в прямому порядку. Мітки записувати через пробіл.\par Алгоритм починає роботу з кореня (на рис. це верхня вершина).\par

\end{minipage}
\vline \hspace{8mm}
\begin{minipage}[t]{8.0cm}
               
{\begin{center}Завдання {8}\end{center}}
\vspace{-6mm}
Для графа на рисунку з вершини 0 запущено алгоритм пошуку в глибину, що будує кістякове дерево. Списки суміжності вершин переглядаються алгоритмом за зростанням міток вершин.\par
\begin{center}
	\adjustimage{max size={0.9\linewidth}{0.9\paperheight}}
	{./16Graphs/Traversals/137.png}
\end{center}
\par Які з наведених ребер входять до складу отриманого кістякового дерева?\par
1) (5, 6)\par
2) (2, 4)\par
3) (1, 6)\par
4) (1, 2)\par
5) (3, 5)\par
6) (1, 3)\par
{\vskip 15pt} \par

\end{minipage}

               
{\begin{center}Завдання {9}\end{center}}
\vspace{-6mm}
Рядки журналу через кому містять тип події (ERROR, WARNING, INFO, STATUS), ім'я користувача (складається з латиниці), час події (фор\-мат як у прикладі), код події (ціле).

Білі символи навколо роздільників не допускаються.

Приклад такого рядка присвоєно змінній \kw{s}.



\begin{verbatim}
s = 'ERROR,username,2022-05-05 13:26,404'
\end{verbatim}


Нижче наведено код, за виконання якого значенням змінної \kw{out} має стати ім'я користувача.

Для наведеного прикладу значенням змінної \kw{out} має стати рядок \code{username}



\begin{verbatim}
res = re.fullmatch(r'XXX', s)
s = ''
out = YYY
\end{verbatim}


Код має працювати не тільки для конкретного рядка з прикладу, але й для кожного рядка з журналу.


Який рядок треба записати на місце \kw{XXX} ?
Який рядок треба записати на місце \kw{YYY} ?\par

%\end {large}
}
\newpage

{{23.05.2025 Контрольна робота, варіант  \No \textbf { 46 }} група К1{\parbox {0.5 cm} \dotfill}  ПІБ {\parbox {2.5 cm} \dotfill}\par}
{
%\begin {large}
\begin{minipage}[t]{9.7cm}
               
{\begin{center}Завдання {1}\end{center}}
\vspace{-6mm}
Укажіть ТИП значення виразу.\par\par
\begin{verbatim}
{21,2,4}
\end{verbatim}

\vspace{15mm}
У завданнях 2, 3 вказати, що буде виведено за виконання. Повідомлення інтерпретатора про помилку позначати ERROR

               
{\begin{center}Завдання {2}\end{center}}
\vspace{-6mm}

\begin{verbatim}
try:
    t = {39:9, 55:3, 30:4, 70:0, 55:9}
    t[58] = 3
    for x, y in t.items():
        print(x, y, sep=' ', end=' ')
except:
    print('ERR')
\end{verbatim}


\end{minipage}
\vline \hspace{8mm}
\begin{minipage}[t]{8.0cm}                
               
{\begin{center}Завдання {3}\end{center}}
\vspace{-6mm}

\begin{verbatim}
c = 0
h = 1

class B:
    h = 2
    
    def __init__(self):
        self.h = 3
        h = 4
        B.h = 5

obj = B()
try:
    print(c, end=' ')
    print(h, end=' ')
    print(B.h, end=' ')
    print(obj.h)
except:
    print('ERR')
\end{verbatim}

\end{minipage}

               
{\begin{center}Завдання {4}\end{center}}
\vspace{-6mm}
Нехай змінна \kw{s} позначає дуже довгу послідовність, по якій можна ітерувати.

Написати вираз-генератор, що задає підпослідовність її елементів, що є рядками типу \kw{str}.

Обмеження: усі елементи шуканої підпослідовності одночасно в пам'яті розміщені бути не можуть.\par

\vspace{10mm}
               
{\begin{center}Завдання {5}\end{center}}
\vspace{-6mm}
Змінні \kw{next} та \kw{prev} вузла списку позначають наступний та попередній за ним вузол відповідно. Починаючи з голови, у список записано послідовні цілі числа від~$-52$ до~$48$. Нехай змінна \kw{p} позначає вузол, що зберігає значення~$-5$.
Яке число зберігається у вузлі \kw{p.prev.prev.next.prev.prev} ?\par

               
{\begin{center}Завдання {6}\end{center}}
\vspace{-6mm}
Значенням змінної \kw{a} є цілочисловий масив типу $numpy.ndarray$ розмірів $3{\times}5$. На скільки збільшиться сума елементів масиву \kw{a} після виконання
\begin{verbatim}
a += 3
\end{verbatim}


Якщо код некоректний, то в якості відповіді зазначити ERROR.\par

               
{\begin{center}Завдання {10}\end{center}}
\vspace{-6mm}

Написати функцію $f$, яка для файлу, заданого шляхом, розв’язує задачу:\\ Текстовий файл містить послідовність чисел $a_1$, $a_2$, $\ldots$, $a_t$, записану у форматі одне число на рядок. Знайти третє найбільше значення, що записане у файлі. Кожне значення, що записане у файлі, враховується тільки один раз. Для послідовності 1, 4, 12, 5, 1, 0 таким значенням буде 4.

В усіх випадках розміри файлів такі, що їх вміст повністю завантажений у пам'ять бути не може. Якість коду також оцінюється.



\vspace{10mm}\par

\newpage
\begin{minipage}[t]{9.7cm}
               
{\begin{center}Завдання {7}\end{center}}
\vspace{-6mm}

\begin{center}
	\adjustimage{max size={0.8\linewidth}{0.8\paperheight}}
	{./15BinTree/trees_exp/17.png}
\end{center}
Указати послідовність міток вершин дерева, зображеного на рис., що утвориться в результаті обходу дерева в оберненому порядку. Мітки записувати через пробіл.\par Алгоритм починає роботу з кореня (на рис. це верхня вершина).\par

\end{minipage}
\vline \hspace{8mm}
\begin{minipage}[t]{8.0cm}
               
{\begin{center}Завдання {8}\end{center}}
\vspace{-6mm}
Для графа на рисунку з вершини 0 запущено алгоритм пошуку в ширину, що будує кістякове дерево. Списки суміжності вершин переглядаються алгоритмом за зростанням міток вершин.\par
\begin{center}
	\adjustimage{max size={0.9\linewidth}{0.9\paperheight}}
	{./16Graphs/Traversals/75.png}
\end{center}
\par Які з наведених ребер входять до складу отриманого кістякового дерева?\par
1) (2, 6)\par
2) (2, 3)\par
3) (0, 4)\par
4) (1, 6)\par
5) (2, 4)\par
6) (3, 4)\par
{\vskip 15pt} \par

\end{minipage}

               
{\begin{center}Завдання {9}\end{center}}
\vspace{-6mm}
Файл \code{'1.txt'} містить відомість з рядками, в яких через двокрапку записано поряд\-ковий номер (ціле), назву предмету (знаків пунктуації не містить), прізвище студента, ім'я студента, номер заліковки, оцінку в балах за 100 бальною шкалою.

Білі символи навколо роздільників не допускаються.

Приклад рядка файлу



\begin{verbatim}
13:algebra:surname:name:123456:77
\end{verbatim}


Нижче наведено код, який шукає усі предмети, який складав студент із заліковкою \code{777}, та зберігає їх типи в змінній \kw{codes}.



\begin{verbatim}
codes = set()
regex = re.compile(r'XXX')
with open('1.txt', encoding='utf8') as f:
    for s in f:
        res = re.fullmatch(regex, s)
        if not res: continue
        s = ''
        codes.add(YYY)
\end{verbatim}



Який рядок треба записати на місце \kw{XXX} ?
Який рядок треба записати на місце \kw{YYY} ?\par

%\end {large}
}
\newpage

{{23.05.2025 Контрольна робота, варіант  \No \textbf { 47 }} група К1{\parbox {0.5 cm} \dotfill}  ПІБ {\parbox {2.5 cm} \dotfill}\par}
{
%\begin {large}
\begin{minipage}[t]{9.7cm}
               
{\begin{center}Завдання {1}\end{center}}
\vspace{-6mm}
Укажіть ТИП значення виразу.\par\par
\begin{verbatim}
{1:3, 5:8}
\end{verbatim}

\vspace{15mm}
У завданнях 2, 3 вказати, що буде виведено за виконання. Повідомлення інтерпретатора про помилку позначати ERROR

               
{\begin{center}Завдання {2}\end{center}}
\vspace{-6mm}

\begin{verbatim}
try:
    s = {52:0, 55:9, 21:9, 36:5, 52:0}
    s[67] = 5
    for x in s.items():
        print(x, sep=' ', end=' ')
except:
    print('ERR')
\end{verbatim}


\end{minipage}
\vline \hspace{8mm}
\begin{minipage}[t]{8.0cm}                
               
{\begin{center}Завдання {3}\end{center}}
\vspace{-6mm}

\begin{verbatim}
__b = 0

class A:
    __b = 1
    
    def __init__(self):
        self.__b = 2
        __b = 3
        A.__b = 4

obj = A()
obj.__b = 5
try:
    print(__b, end=' ')
    print(obj.__b, end=' ')
    print(A.__b, end=' ')
    print(obj.__b, end=' ')
    print(obj._A__b)
except:
    print('ERR')
\end{verbatim}

\end{minipage}

               
{\begin{center}Завдання {4}\end{center}}
\vspace{-6mm}

Записати ВИРАЗ, значенням якого є список, що складається з квадратних коренів цілих чисел від 17 до 2005 (включно), що діляться на 7.\par

\vspace{10mm}
               
{\begin{center}Завдання {5}\end{center}}
\vspace{-6mm}
Змінні \kw{next} та \kw{prev} вузла списку позначають наступний та попередній за ним вузол відповідно. Починаючи з голови, у список записано послідовні цілі числа від~$-54$ до~$26$. Нехай змінна \kw{p} позначає вузол, що зберігає значення~$7$.
Яке число зберігається у вузлі \kw{p.prev.prev.prev.next.prev} ?\par

               
{\begin{center}Завдання {6}\end{center}}
\vspace{-6mm}
Значенням змінної \kw{a} є цілочисловий масив типу $numpy.ndarray$ розмірів $7{\times}7$. На скільки збільшиться сума елементів масиву \kw{a} після виконання
\begin{verbatim}
a.T[2] += 2
\end{verbatim}


Якщо код некоректний, то в якості відповіді зазначити ERROR.\par

               
{\begin{center}Завдання {10}\end{center}}
\vspace{-6mm}

Написати функцію $f$, яка для файлу, заданого шляхом, розв’язує задачу:\\ Текстовий файл містить послідовність чисел $a_1$, $a_2$, $\ldots$, $a_t$, записану у форматі одне число на рядок. Знайти середнє арифметичне чисел $a_1$, $a_4$, $a_7$, $\ldots$ (починаючи з першого через два).

В усіх випадках розміри файлів такі, що їх вміст повністю завантажений у пам'ять бути не може. Якість коду також оцінюється.



\vspace{10mm}\par

\newpage
\begin{minipage}[t]{9.7cm}
               
{\begin{center}Завдання {7}\end{center}}
\vspace{-6mm}

\begin{center}
	\adjustimage{max size={0.8\linewidth}{0.8\paperheight}}
	{./15BinTree/trees_exp/46.png}
\end{center}
Указати послідовність міток вершин дерева, зображеного на рис., що утвориться в результаті обходу дерева в оберненому порядку. Мітки записувати через пробіл.\par Алгоритм починає роботу з кореня (на рис. це верхня вершина).\par

\end{minipage}
\vline \hspace{8mm}
\begin{minipage}[t]{8.0cm}
               
{\begin{center}Завдання {8}\end{center}}
\vspace{-6mm}
Для графа на рисунку з вершини 3 запущено алгоритм пошуку в ширину, що будує кістякове дерево. Списки суміжності вершин переглядаються алгоритмом за зростанням міток вершин.\par
\begin{center}
	\adjustimage{max size={0.9\linewidth}{0.9\paperheight}}
	{./16Graphs/Traversals/212.png}
\end{center}
\par Які з наведених ребер входять до складу отриманого кістякового дерева?\par
1) (0, 5)\par
2) (1, 5)\par
3) (0, 6)\par
4) (1, 4)\par
5) (3, 6)\par
6) (0, 2)\par
{\vskip 15pt} \par

\end{minipage}

               
{\begin{center}Завдання {9}\end{center}}
\vspace{-6mm}
Файл \code{'1.txt'} містить журнал з рядками, в яких через двокрапку записано поряд\-ко\-вий номер запису в журналі (ціле), тип події (ERROR, WARNING, INFO, STATUS), ім'я користувача (складається з латиниці), час події (фор\-мат як у прикладі), код події (ціле).

Білі символи навколо роздільників не допускаються.

Приклад рядка файлу



\begin{verbatim}
13:ERROR:username:2022-05-05 13-26:404
\end{verbatim}


Нижче наведено код, який шукає усі коди помилок (подія ERROR), що зустріча\-ють\-ся в журналі, та зберігає їх у змінній \kw{codes}.



\begin{verbatim}
codes = set()
regex = re.compile(r'XXX')
with open('1.txt', encoding='utf8') as f:
    for s in f:
        res = re.fullmatch(regex, s)
        if not res: continue
        s = ''
        codes.add(YYY)
\end{verbatim}



Який рядок треба записати на місце \kw{XXX} ?
Який рядок треба записати на місце \kw{YYY} ?\par

%\end {large}
}
\newpage

{{23.05.2025 Контрольна робота, варіант  \No \textbf { 48 }} група К1{\parbox {0.5 cm} \dotfill}  ПІБ {\parbox {2.5 cm} \dotfill}\par}
{
%\begin {large}
\begin{minipage}[t]{9.7cm}
               
{\begin{center}Завдання {1}\end{center}}
\vspace{-6mm}
Укажіть ТИП значення виразу.\par\par
\begin{verbatim}
(2, 7, -4)
\end{verbatim}

\vspace{15mm}
У завданнях 2, 3 вказати, що буде виведено за виконання. Повідомлення інтерпретатора про помилку позначати ERROR

               
{\begin{center}Завдання {2}\end{center}}
\vspace{-6mm}

\begin{verbatim}
try:
    d = {32:5, 70:2, 67:7, 32:3}
    d[83] = 7
    for x in d.items():
        print(*x, sep=' ', end=' ')
except:
    print('ERR')
\end{verbatim}


\end{minipage}
\vline \hspace{8mm}
\begin{minipage}[t]{8.0cm}                
               
{\begin{center}Завдання {3}\end{center}}
\vspace{-6mm}

\begin{verbatim}
d = 0

class D:
    d = 1
    
    def __init__(self):
        self.d = 2
        d = 3
        D.d = 4

obj = D()
D.d = 5
try:
    print(d, end=' ')
    print(obj.d, end=' ')
    print(D.d, end=' ')
    print(obj.d, end=' ')
    print(obj._D__d)
except:
    print('ERR')
\end{verbatim}

\end{minipage}

               
{\begin{center}Завдання {4}\end{center}}
\vspace{-6mm}
Нехай змінна \kw{s} позначає послідовність цілих, по якій можна ітерувати.

Написати вираз, що перевіряє, чи правда, що послідовність  містить хоча б одне число, що ділиться на 7.\par

\vspace{10mm}
               
{\begin{center}Завдання {5}\end{center}}
\vspace{-6mm}
Змінні \kw{next} та \kw{prev} вузла списку позначають наступний та попередній за ним вузол відповідно. Починаючи з голови, у список записано послідовні цілі числа від~$-70$ до~$17$. Нехай змінна \kw{p} позначає вузол, що зберігає значення~$-9$.
Яке число зберігається у вузлі \kw{p.next.next.next.prev.next} ?\par

               
{\begin{center}Завдання {6}\end{center}}
\vspace{-6mm}
Значенням змінної \kw{a} є цілочисловий масив типу $numpy.ndarray$ розмірів $4{\times}6$. На скільки збільшиться сума елементів масиву \kw{a} після виконання
\begin{verbatim}
a[-3] += 1
\end{verbatim}


Якщо код некоректний, то в якості відповіді зазначити ERROR.\par

               
{\begin{center}Завдання {10}\end{center}}
\vspace{-6mm}

Написати функцію $f$, яка для файлу, заданого шляхом, розв’язує задачу:\\ Текстовий файл містить послідовність чисел $a_1$, $a_2$, $\ldots$, $a_t$, записану у форматі одне число на рядок. Знайти третє найменше значення, що записане у файлі. Кожне значення, що записане у файлі, враховується стільки разів, скільки входить у файл. Для послідовності 2, 4, 12, 5, 2, 0 таким значенням буде 2.

В усіх випадках розміри файлів такі, що їх вміст повністю завантажений у пам'ять бути не може. Якість коду також оцінюється.



\vspace{10mm}\par

\newpage
\begin{minipage}[t]{9.7cm}
               
{\begin{center}Завдання {7}\end{center}}
\vspace{-6mm}

\begin{center}
	\adjustimage{max size={0.8\linewidth}{0.8\paperheight}}
	{./15BinTree/trees_exp/113.png}
\end{center}
Указати послідовність міток вершин дерева, зображеного на рис., що утвориться в результаті обходу дерева в прямому порядку. Мітки записувати через пробіл.\par Алгоритм починає роботу з кореня (на рис. це верхня вершина).\par

\end{minipage}
\vline \hspace{8mm}
\begin{minipage}[t]{8.0cm}
               
{\begin{center}Завдання {8}\end{center}}
\vspace{-6mm}
Для графа на рисунку з вершини 3 запущено алгоритм пошуку в ширину, що будує кістякове дерево. Списки суміжності вершин переглядаються алгоритмом за зростанням міток вершин.\par
\begin{center}
	\adjustimage{max size={0.9\linewidth}{0.9\paperheight}}
	{./16Graphs/Traversals/282.png}
\end{center}
\par Які з наведених ребер входять до складу отриманого кістякового дерева?\par
1) (1, 5)\par
2) (2, 6)\par
3) (3, 6)\par
4) (0, 2)\par
5) (1, 3)\par
6) (0, 1)\par
{\vskip 15pt} \par

\end{minipage}

               
{\begin{center}Завдання {9}\end{center}}
\vspace{-6mm}
Файл \code{'1.txt'} містить журнал з рядками, в яких через кому записано поряд\-ко\-вий номер запису в журналі (ціле), тип події (ERROR, WARNING, INFO, STATUS), ім'я користувача (складається з латиниці), час події (фор\-мат як у прикладі), опис події.

Білі символи навколо роздільників не допускаються.

Приклад рядка файлу



\begin{verbatim}
13,ERROR,username,2022-05-05 13:26,something happens
\end{verbatim}


Нижче наведено код, який має залишити у файлі в кожному рядку тільки тип події та її опис.

Рядок з прикладу має бути замінений на такий



\begin{verbatim}
ERROR,something happens
\end{verbatim}




Код:



\begin{verbatim}
regex = re.compile(r'XXX')
with open('1.txt', encoding='utf8') as f:
    lst = f.readlines()
for i in range(len(lst)):
    lst[i] = re.sub(regex, YYY, lst[i])
with open('1.txt', 'w', encoding='utf8') as f:
    f.writelines(lst)
\end{verbatim}







Який рядок треба записати на місце \kw{XXX} ?
Який рядок треба записати на місце \kw{YYY} ?\par

%\end {large}
}
\newpage

{{23.05.2025 Контрольна робота, варіант  \No \textbf { 49 }} група К1{\parbox {0.5 cm} \dotfill}  ПІБ {\parbox {2.5 cm} \dotfill}\par}
{
%\begin {large}
\begin{minipage}[t]{9.7cm}
               
{\begin{center}Завдання {1}\end{center}}
\vspace{-6mm}
Укажіть ТИП значення виразу.\par\par
\begin{verbatim}
{y:y*y for y in range(-7,7)}
\end{verbatim}

\vspace{15mm}
У завданнях 2, 3 вказати, що буде виведено за виконання. Повідомлення інтерпретатора про помилку позначати ERROR

               
{\begin{center}Завдання {2}\end{center}}
\vspace{-6mm}

\begin{verbatim}
try:
    d = {78:2, 38:2, 19:5, 19:0}
    d[51] = 2
    for x in d :
        print(x, sep=' ', end=' ')
except:
    print('ERR')
\end{verbatim}


\end{minipage}
\vline \hspace{8mm}
\begin{minipage}[t]{8.0cm}                
               
{\begin{center}Завдання {3}\end{center}}
\vspace{-6mm}

\begin{verbatim}
d = 0

class D:
    d = 1
    
    def __init__(self):
        self.d = 2
        d = 3
        D.d = 4

obj = D()
D.d = 5
try:
    print(d, end=' ')
    print(obj.d, end=' ')
    print(D.d, end=' ')
    print(obj.d, end=' ')
    print(obj._D__d)
except:
    print('ERR')
\end{verbatim}

\end{minipage}

               
{\begin{center}Завдання {4}\end{center}}
\vspace{-6mm}
Задано множину \kw{x}, що складається з цілих чисел. Записати ВИРАЗ, значенням якого є множина, що складається  з квадратів елементів множини \kw{x}.\par

\vspace{10mm}
               
{\begin{center}Завдання {5}\end{center}}
\vspace{-6mm}
Змінні \kw{next} та \kw{prev} вузла списку позначають наступний та попередній за ним вузол відповідно. Починаючи з голови, у список записано послідовні цілі числа від~$-29$ до~$23$. Нехай змінна \kw{p} позначає вузол, що зберігає значення~$9$.
Яке число зберігається у вузлі \kw{p.prev.next.prev.next.prev} ?\par

               
{\begin{center}Завдання {6}\end{center}}
\vspace{-6mm}
Значенням змінної \kw{a} є цілочисловий масив типу $numpy.ndarray$ розмірів $5{\times}3$. На скільки збільшиться сума елементів масиву \kw{a} після виконання
\begin{verbatim}
a.T += 1
\end{verbatim}


Якщо код некоректний, то в якості відповіді зазначити ERROR.\par

               
{\begin{center}Завдання {10}\end{center}}
\vspace{-6mm}

Написати функцію $f$, яка для файлу, заданого шляхом, розв’язує задачу:\\ Кожен з двох текстових файлів містить послідовність чисел, записану у форматі одне число на рядок.  Перевірити, чи задають вони одну ту саму послідовність.

В усіх випадках розміри файлів такі, що їх вміст повністю завантажений у пам'ять бути не може. Якість коду також оцінюється.



\vspace{10mm}\par

\newpage
\begin{minipage}[t]{9.7cm}
               
{\begin{center}Завдання {7}\end{center}}
\vspace{-6mm}

\begin{center}
	\adjustimage{max size={0.8\linewidth}{0.8\paperheight}}
	{./15BinTree/trees_exp/108.png}
\end{center}
Указати послідовність міток вершин дерева, зображеного на рис., що утвориться в результаті обходу дерева в прямому порядку. Мітки записувати через пробіл.\par Алгоритм починає роботу з кореня (на рис. це верхня вершина).\par

\end{minipage}
\vline \hspace{8mm}
\begin{minipage}[t]{8.0cm}
               
{\begin{center}Завдання {8}\end{center}}
\vspace{-6mm}
Для графа на рисунку з вершини 3 запущено алгоритм пошуку в глибину, що будує кістякове дерево. Списки суміжності вершин переглядаються алгоритмом за зростанням міток вершин.\par
\begin{center}
	\adjustimage{max size={0.9\linewidth}{0.9\paperheight}}
	{./16Graphs/Traversals/332.png}
\end{center}
\par Які з наведених ребер входять до складу отриманого кістякового дерева?\par
1) (0, 1)\par
2) (2, 3)\par
3) (0, 4)\par
4) (5, 6)\par
5) (2, 5)\par
6) (0, 6)\par
{\vskip 15pt} \par

\end{minipage}

               
{\begin{center}Завдання {9}\end{center}}
\vspace{-6mm}
Файл \code{'1.txt'} містить журнал з рядками, в яких через кому записано поряд\-ко\-вий номер запису в журналі (ціле), код події (ціле), тип події (ERROR, \\WARNING, INFO, STATUS), ім'я користувача (складається з латиниці), час події (фор\-мат як у прикладі), опис події.

Білі символи навколо роздільників не допускаються.

Приклад рядка файлу



\begin{verbatim}
13,404,ERROR,username,2022-05-05 13:26,something happens
\end{verbatim}


Нижче наведено код, який має залишити у файлі в кожному рядку тільки код події та ім'я користувача. Рядок з прикладу має бути замінений на такий



\begin{verbatim}
404,username
\end{verbatim}




Код:



\begin{verbatim}
regex = re.compile(r'XXX')
with open('1.txt', encoding='utf8') as f:
    lst = f.readlines()
for i in range(len(lst)):
    lst[i] = re.sub(regex, YYY, lst[i])
with open('1.txt', 'w', encoding='utf8') as f:
    f.writelines(lst)
\end{verbatim}





Який рядок треба записати на місце \kw{XXX} ?
Який рядок треба записати на місце \kw{YYY} ?\par

%\end {large}
}
\newpage

{{23.05.2025 Контрольна робота, варіант  \No \textbf { 50 }} група К1{\parbox {0.5 cm} \dotfill}  ПІБ {\parbox {2.5 cm} \dotfill}\par}
{
%\begin {large}
\begin{minipage}[t]{9.7cm}
               
{\begin{center}Завдання {1}\end{center}}
\vspace{-6mm}
Укажіть ТИП значення виразу.\par\par
\begin{verbatim}
{(1,4), 8}
\end{verbatim}

\vspace{15mm}
У завданнях 2, 3 вказати, що буде виведено за виконання. Повідомлення інтерпретатора про помилку позначати ERROR

               
{\begin{center}Завдання {2}\end{center}}
\vspace{-6mm}

\begin{verbatim}
try:
    s = {62:7, 35:2, 54:6, 84:0, 62:1}
    s[67] = 5
    for x in s.keys():
        print(x, sep=' ', end=' ')
except:
    print('ERR')
\end{verbatim}


\end{minipage}
\vline \hspace{8mm}
\begin{minipage}[t]{8.0cm}                
               
{\begin{center}Завдання {3}\end{center}}
\vspace{-6mm}

\begin{verbatim}
__a = 0

class A:
    __a = 1
    
    def __init__(self):
        self.__a = 2
        __a = 3
        A.__a = 4

obj = A()
obj.__a = 5
try:
    print(__a, end=' ')
    print(obj.__a, end=' ')
    print(A.__a, end=' ')
    print(obj.__a, end=' ')
    print(obj._A__a)
except:
    print('ERR')
\end{verbatim}

\end{minipage}

               
{\begin{center}Завдання {4}\end{center}}
\vspace{-6mm}
Записати ВИРАЗ, значенням якого є словник, ключами якого є цілі числа від 16 до 2018 (включно), що не діляться на жодне з чисел 2, 3, а ключам відповідають їх квадратні корені\par

\vspace{10mm}
               
{\begin{center}Завдання {5}\end{center}}
\vspace{-6mm}
Змінні \kw{next} та \kw{prev} вузла списку позначають наступний та попередній за ним вузол відповідно. Починаючи з голови, у список записано послідовні цілі числа від~$-35$ до~$42$. Нехай змінна \kw{p} позначає вузол, що зберігає значення~$0$.
Яке число зберігається у вузлі \kw{p.next.prev.next.prev.next} ?\par

               
{\begin{center}Завдання {6}\end{center}}
\vspace{-6mm}
Значенням змінної \kw{a} є цілочисловий масив типу $numpy.ndarray$ розмірів $6{\times}4$. На скільки збільшиться сума елементів масиву \kw{a} після виконання
\begin{verbatim}
a[-2] += 3
\end{verbatim}


Якщо код некоректний, то в якості відповіді зазначити ERROR.\par

               
{\begin{center}Завдання {10}\end{center}}
\vspace{-6mm}

Написати функцію $f$, яка для файлу, заданого шляхом, розв’язує задачу:\\ Текстовий файл містить послідовність чисел $a_1$, $a_2$, $\ldots$, $a_t$, записану у форматі одне число на рядок. Знайти третє найменше значення, що записане у файлі. Кожне значення, що записане у файлі, враховується стільки разів, скільки входить у файл. Для послідовності 2, 4, 12, 5, 2, 0 таким значенням буде 2.

В усіх випадках розміри файлів такі, що їх вміст повністю завантажений у пам'ять бути не може. Якість коду також оцінюється.



\vspace{10mm}\par

\newpage
\begin{minipage}[t]{9.7cm}
               
{\begin{center}Завдання {7}\end{center}}
\vspace{-6mm}

\begin{center}
	\adjustimage{max size={0.8\linewidth}{0.8\paperheight}}
	{./15BinTree/trees_exp/79.png}
\end{center}
Указати послідовність міток вершин дерева, зображеного на рис., що утвориться в результаті обходу дерева в оберненому порядку. Мітки записувати через пробіл.\par Алгоритм починає роботу з кореня (на рис. це верхня вершина).\par

\end{minipage}
\vline \hspace{8mm}
\begin{minipage}[t]{8.0cm}
               
{\begin{center}Завдання {8}\end{center}}
\vspace{-6mm}
Для графа на рисунку з вершини 0 запущено алгоритм пошуку в ширину, що будує кістякове дерево. Списки суміжності вершин переглядаються алгоритмом за зростанням міток вершин.\par
\begin{center}
	\adjustimage{max size={0.9\linewidth}{0.9\paperheight}}
	{./16Graphs/Traversals/59.png}
\end{center}
\par Які з наведених ребер входять до складу отриманого кістякового дерева?\par
1) (3, 4)\par
2) (3, 6)\par
3) (3, 5)\par
4) (4, 5)\par
5) (0, 4)\par
6) (1, 4)\par
{\vskip 15pt} \par

\end{minipage}

               
{\begin{center}Завдання {9}\end{center}}
\vspace{-6mm}
Файл \code{'1.txt'} містить журнал з рядками, в яких через двокрапку записано поряд\-ко\-вий номер запису в журналі (ціле), тип події (ERROR, WARNING, INFO, STATUS), ім'я користувача (складається з латиниці), час події (фор\-мат як у прикладі), код події (ціле).

Білі символи навколо роздільників не допускаються.

Приклад рядка файлу



\begin{verbatim}
13:ERROR:username:2022-05-05 13-26:404
\end{verbatim}


Нижче наведено код, який має замінити у файлі імена користувачів на рядок \code{unknown}.



\begin{verbatim}
regex = re.compile(r'XXX')
with open('1.txt', encoding='utf8') as f:
    lst = f.readlines()
for i in range(len(lst)):
    lst[i] = re.sub(regex, YYY, lst[i])
with open('1.txt', 'w', encoding='utf8') as f:
    f.writelines(lst)
\end{verbatim}





Який рядок треба записати на місце \kw{XXX} ?
Який рядок треба записати на місце \kw{YYY} ?\par

%\end {large}
}
\newpage

{{23.05.2025 Контрольна робота, варіант  \No \textbf { 51 }} група К1{\parbox {0.5 cm} \dotfill}  ПІБ {\parbox {2.5 cm} \dotfill}\par}
{
%\begin {large}
\begin{minipage}[t]{9.7cm}
               
{\begin{center}Завдання {1}\end{center}}
\vspace{-6mm}
Укажіть ТИП значення виразу.\par\par
\begin{verbatim}
{x:x for x in range(7)}
\end{verbatim}

\vspace{15mm}
У завданнях 2, 3 вказати, що буде виведено за виконання. Повідомлення інтерпретатора про помилку позначати ERROR

               
{\begin{center}Завдання {2}\end{center}}
\vspace{-6mm}

\begin{verbatim}
try:
    t = {28:0, 17:0, 58:4, 58:3}
    t[38] = 2
    for x in t.items():
        print(*x, sep=' ', end=' ')
except:
    print('ERR')
\end{verbatim}


\end{minipage}
\vline \hspace{8mm}
\begin{minipage}[t]{8.0cm}                
               
{\begin{center}Завдання {3}\end{center}}
\vspace{-6mm}

\begin{verbatim}
_b = 0

class C:
    _b = 1
    
    def __init__(self):
        self._b = 2
        _b = 3
        C._b = 4

obj = C()
C._b = 5
try:
    print(_b, end=' ')
    print(obj._b, end=' ')
    print(C._b, end=' ')
    print(obj._b, end=' ')
    print(obj._C__b)
except:
    print('ERR')
\end{verbatim}

\end{minipage}

               
{\begin{center}Завдання {4}\end{center}}
\vspace{-6mm}
Нехай змінна \kw{s} позначає дуже довгу послідовність цілих, по якій можна ітерувати.

Написати вираз-генератор, що задає підпослідовність її елементів, що менші за задане число num.

Обмеження: усі елементи шуканої підпослідовності одночасно в пам'яті розміщені бути не можуть.\par

\vspace{10mm}
               
{\begin{center}Завдання {5}\end{center}}
\vspace{-6mm}
Змінні \kw{next} та \kw{prev} вузла списку позначають наступний та попередній за ним вузол відповідно. Починаючи з голови, у список записано послідовні цілі числа від~$-33$ до~$18$. Нехай змінна \kw{p} позначає вузол, що зберігає значення~$4$.
Яке число зберігається у вузлі \kw{p.next.prev.prev.prev.prev} ?\par

               
{\begin{center}Завдання {6}\end{center}}
\vspace{-6mm}
Значенням змінної \kw{a} є цілочисловий масив типу $numpy.ndarray$ розмірів $3{\times}4$. На скільки збільшиться сума елементів масиву \kw{a} після виконання
\begin{verbatim}
a.T[0] += 2
\end{verbatim}


Якщо код некоректний, то в якості відповіді зазначити ERROR.\par

               
{\begin{center}Завдання {10}\end{center}}
\vspace{-6mm}

Написати функцію $f$, яка для файлу, заданого шляхом, розв’язує задачу:\\ Текстовий файл містить послідовність чисел $a_1$, $a_2$, $\ldots$, $a_t$, записану у форматі одне число на рядок. Знайти середнє арифметичне чисел $a_1$, $a_4$, $a_7$, $\ldots$ (починаючи з першого через два).

В усіх випадках розміри файлів такі, що їх вміст повністю завантажений у пам'ять бути не може. Якість коду також оцінюється.



\vspace{10mm}\par

\newpage
\begin{minipage}[t]{9.7cm}
               
{\begin{center}Завдання {7}\end{center}}
\vspace{-6mm}

\begin{center}
	\adjustimage{max size={0.8\linewidth}{0.8\paperheight}}
	{./15BinTree/trees_exp/38.png}
\end{center}
Указати послідовність міток вершин дерева, зображеного на рис., що утвориться в результаті обходу дерева в прямому порядку. Мітки записувати через пробіл.\par Алгоритм починає роботу з кореня (на рис. це верхня вершина).\par

\end{minipage}
\vline \hspace{8mm}
\begin{minipage}[t]{8.0cm}
               
{\begin{center}Завдання {8}\end{center}}
\vspace{-6mm}
Для графа на рисунку з вершини 4 запущено алгоритм пошуку в глибину, що будує кістякове дерево. Списки суміжності вершин переглядаються алгоритмом за зростанням міток вершин.\par
\begin{center}
	\adjustimage{max size={0.9\linewidth}{0.9\paperheight}}
	{./16Graphs/Traversals/331.png}
\end{center}
\par Які з наведених ребер входять до складу отриманого кістякового дерева?\par
1) (0, 3)\par
2) (0, 6)\par
3) (4, 6)\par
4) (1, 6)\par
5) (2, 5)\par
6) (3, 4)\par
{\vskip 15pt} \par

\end{minipage}

               
{\begin{center}Завдання {9}\end{center}}
\vspace{-6mm}
Рядки відомості через кому містять порядковий номер (ціле), пріз\-ви\-ще студен\-та, ім'я студента, номер заліковки, оцінку в балах за 100 бальною шкалою.

Білі символи навколо роздільників не допускаються.

Приклад такого рядка присвоєно змінній \kw{s}.



\begin{verbatim}
s = '13,surname,name,123456,77'
\end{verbatim}


Нижче наведено код, за виконання якого для рядка \kw{s} значенням змінної \kw{out} має стати номер заліковки.

Для наведеного прикладу значенням змінної \kw{out} має стати рядок \code{123456} .



\begin{verbatim}
res = re.fullmatch(r'XXX', s)
s = ''
out = YYY
\end{verbatim}


Код має працювати не тільки для конкретного рядка з прикладу, але й для кожного рядка з відомості.


Який рядок треба записати на місце \kw{XXX} ?
Який рядок треба записати на місце \kw{YYY} ?\par

%\end {large}
}
\newpage

{{23.05.2025 Контрольна робота, варіант  \No \textbf { 52 }} група К1{\parbox {0.5 cm} \dotfill}  ПІБ {\parbox {2.5 cm} \dotfill}\par}
{
%\begin {large}
\begin{minipage}[t]{9.7cm}
               
{\begin{center}Завдання {1}\end{center}}
\vspace{-6mm}
Укажіть ТИП значення виразу.\par\par
\begin{verbatim}
{y:None for y in range(-8,3)}
\end{verbatim}

\vspace{15mm}
У завданнях 2, 3 вказати, що буде виведено за виконання. Повідомлення інтерпретатора про помилку позначати ERROR

               
{\begin{center}Завдання {2}\end{center}}
\vspace{-6mm}

\begin{verbatim}
try:
    s = {72:7, 36:3, 44:0, 49:5, 44:3}
    s[47] = 3
    for x, y in s.items():
        print(x, y, sep=' ', end=' ')
except:
    print('ERR')
\end{verbatim}


\end{minipage}
\vline \hspace{8mm}
\begin{minipage}[t]{8.0cm}                
               
{\begin{center}Завдання {3}\end{center}}
\vspace{-6mm}

\begin{verbatim}
b = 0
m = 1

class C:
    m = 2
    
    def __init__(self):
        global m
        self.b = 3
        m = 4
        C.m = 5

obj = C()
try:
    print(b, end=' ')
    print(m, end=' ')
    print(C.m, end=' ')
    print(obj.b)
except:
    print('ERR')
\end{verbatim}

\end{minipage}

               
{\begin{center}Завдання {4}\end{center}}
\vspace{-6mm}
Нехай змінна \kw{s} позначає дуже довгу послідовність, по якій можна ітерувати.

Написати вираз-генератор, що задає підпослідовність її елементів, що є числами типу \kw{float}.

Обмеження: усі елементи шуканої підпослідовності одночасно в пам'яті розміщені бути не можуть.\par

\vspace{10mm}
               
{\begin{center}Завдання {5}\end{center}}
\vspace{-6mm}
Змінні \kw{next} та \kw{prev} вузла списку позначають наступний та попередній за ним вузол відповідно. Починаючи з голови, у список записано послідовні цілі числа від~$-27$ до~$25$. Нехай змінна \kw{p} позначає вузол, що зберігає значення~$-3$.
Яке число зберігається у вузлі \kw{p.prev.next.next.next.next} ?\par

               
{\begin{center}Завдання {6}\end{center}}
\vspace{-6mm}
Значенням змінної \kw{a} є цілочисловий масив типу $numpy.ndarray$ розмірів $7{\times}5$. На скільки збільшиться сума елементів масиву \kw{a} після виконання
\begin{verbatim}
a[1] += 1
\end{verbatim}


Якщо код некоректний, то в якості відповіді зазначити ERROR.\par

               
{\begin{center}Завдання {10}\end{center}}
\vspace{-6mm}

Написати функцію $f$, яка для файлу, заданого шляхом, розв’язує задачу:\\ Текстовий файл містить послідовність чисел $a_1$, $a_2$, $\ldots$, $a_t$, записану у форматі одне число на рядок. Знайти третє найменше значення, що записане у файлі. Кожне значення, що записане у файлі, враховується тільки один раз. \par
Для послідовності 1, 4, 12, 5, 1, 0 таким значенням буде 4.

В усіх випадках розміри файлів такі, що їх вміст повністю завантажений у пам'ять бути не може. Якість коду також оцінюється.



\vspace{10mm}\par

\newpage
\begin{minipage}[t]{9.7cm}
               
{\begin{center}Завдання {7}\end{center}}
\vspace{-6mm}

\begin{center}
	\adjustimage{max size={0.8\linewidth}{0.8\paperheight}}
	{./15BinTree/trees_exp/49.png}
\end{center}
Указати послідовність міток вершин дерева, зображеного на рис., що утвориться в результаті обходу дерева в оберненому порядку. Мітки записувати через пробіл.\par Алгоритм починає роботу з кореня (на рис. це верхня вершина).\par

\end{minipage}
\vline \hspace{8mm}
\begin{minipage}[t]{8.0cm}
               
{\begin{center}Завдання {8}\end{center}}
\vspace{-6mm}
Для графа на рисунку з вершини 1 запущено алгоритм пошуку в глибину, що будує кістякове дерево. Списки суміжності вершин переглядаються алгоритмом за зростанням міток вершин.\par
\begin{center}
	\adjustimage{max size={0.9\linewidth}{0.9\paperheight}}
	{./16Graphs/Traversals/229.png}
\end{center}
\par Які з наведених ребер входять до складу отриманого кістякового дерева?\par
1) (4, 5)\par
2) (3, 5)\par
3) (1, 2)\par
4) (5, 6)\par
5) (2, 5)\par
6) (0, 6)\par
{\vskip 15pt} \par

\end{minipage}

               
{\begin{center}Завдання {9}\end{center}}
\vspace{-6mm}
Рядки відомості через двокрапку містять порядковий номер (ціле), назву предмету (знаків пунктуації не містить), прізвище студента, ім'я студента, оцінку в балах за 100 бальною шкалою.

Білі символи навколо роздільників не допускаються.

Приклад такого рядка присвоєно змінній \kw{s}.



\begin{verbatim}
s = '13:algebra:surname:name:77'
\end{verbatim}


Нижче наведено код, за виконання якого значенням змінної \kw{out} має стати прізвище студента.

Для наведеного прикладу значенням змінної \kw{out} має стати рядок \code{surname}



\begin{verbatim}
res = re.fullmatch(r'XXX', s)
s = ''
out = YYY
\end{verbatim}


Код має працювати не тільки для конкретного рядка з прикладу, але й для кожного рядка з відомості.


Який рядок треба записати на місце \kw{XXX} ?
Який рядок треба записати на місце \kw{YYY} ?\par

%\end {large}
}
\newpage

{{23.05.2025 Контрольна робота, варіант  \No \textbf { 53 }} група К1{\parbox {0.5 cm} \dotfill}  ПІБ {\parbox {2.5 cm} \dotfill}\par}
{
%\begin {large}
\begin{minipage}[t]{9.7cm}
               
{\begin{center}Завдання {1}\end{center}}
\vspace{-6mm}
Укажіть ТИП значення виразу.\par\par
\begin{verbatim}
(2, 7, -4)
\end{verbatim}

\vspace{15mm}
У завданнях 2, 3 вказати, що буде виведено за виконання. Повідомлення інтерпретатора про помилку позначати ERROR

               
{\begin{center}Завдання {2}\end{center}}
\vspace{-6mm}

\begin{verbatim}
try:
    d = {69:2, 22:3, 52:1, 22:1}
    d[65] = 9
    for x in d.values():
        print(x, sep=' ', end=' ')
except:
    print('ERR')
\end{verbatim}


\end{minipage}
\vline \hspace{8mm}
\begin{minipage}[t]{8.0cm}                
               
{\begin{center}Завдання {3}\end{center}}
\vspace{-6mm}

\begin{verbatim}
__c = 0

class B:
    __c = 1
    
    def __init__(self):
        self.__c = 2
        __c = 3
        B.__c = 4

obj = B()
obj.__c = 5
try:
    print(__c, end=' ')
    print(obj.__c, end=' ')
    print(B.__c, end=' ')
    print(obj.__c, end=' ')
    print(obj._B__c)
except:
    print('ERR')
\end{verbatim}

\end{minipage}

               
{\begin{center}Завдання {4}\end{center}}
\vspace{-6mm}

Записати ВИРАЗ, значенням якого є список, що складається з квадратів цілих чисел від 16 до 2001 (включно), що не діляться на жодне з чисел 2, 3.\par

\vspace{10mm}
               
{\begin{center}Завдання {5}\end{center}}
\vspace{-6mm}
Змінні \kw{next} та \kw{prev} вузла списку позначають наступний та попередній за ним вузол відповідно. Починаючи з голови, у список записано послідовні цілі числа від~$-25$ до~$19$. Нехай змінна \kw{p} позначає вузол, що зберігає значення~$-2$.
Яке число зберігається у вузлі \kw{p.prev.prev.next.next.next} ?\par

               
{\begin{center}Завдання {6}\end{center}}
\vspace{-6mm}
Значенням змінної \kw{a} є цілочисловий масив типу $numpy.ndarray$ розмірів $5{\times}7$. На скільки збільшиться сума елементів масиву \kw{a} після виконання
\begin{verbatim}
a.T[-1] += 3
\end{verbatim}


Якщо код некоректний, то в якості відповіді зазначити ERROR.\par

               
{\begin{center}Завдання {10}\end{center}}
\vspace{-6mm}

Написати функцію $f$, яка для файлу, заданого шляхом, розв’язує задачу:\\ Текстовий файл містить послідовність чисел $a_1$, $a_2$, $\ldots$, $a_t$, записану у форматі одне число на рядок. Знайти третє найбільше значення, що записане у файлі. Кожне значення, що записане у файлі, враховується тільки один раз. Для послідовності 1, 4, 12, 5, 1, 0 таким значенням буде 4.

В усіх випадках розміри файлів такі, що їх вміст повністю завантажений у пам'ять бути не може. Якість коду також оцінюється.



\vspace{10mm}\par

\newpage
\begin{minipage}[t]{9.7cm}
               
{\begin{center}Завдання {7}\end{center}}
\vspace{-6mm}

\begin{center}
	\adjustimage{max size={0.8\linewidth}{0.8\paperheight}}
	{./15BinTree/trees_exp/120.png}
\end{center}
Указати послідовність міток вершин дерева, зображеного на рис., що утвориться в результаті обходу дерева в оберненому порядку. Мітки записувати через пробіл.\par Алгоритм починає роботу з кореня (на рис. це верхня вершина).\par

\end{minipage}
\vline \hspace{8mm}
\begin{minipage}[t]{8.0cm}
               
{\begin{center}Завдання {8}\end{center}}
\vspace{-6mm}
Для графа на рисунку з вершини 3 запущено алгоритм пошуку в ширину, що будує кістякове дерево. Списки суміжності вершин переглядаються алгоритмом за зростанням міток вершин.\par
\begin{center}
	\adjustimage{max size={0.9\linewidth}{0.9\paperheight}}
	{./16Graphs/Traversals/94.png}
\end{center}
\par Які з наведених ребер входять до складу отриманого кістякового дерева?\par
1) (0, 4)\par
2) (1, 4)\par
3) (0, 1)\par
4) (2, 3)\par
5) (1, 2)\par
6) (3, 5)\par
{\vskip 15pt} \par

\end{minipage}

               
{\begin{center}Завдання {9}\end{center}}
\vspace{-6mm}
Рядки журналу через двокрапку містять ім'я користувача (складається з латиниці), час події (фор\-мат як у прикладі), тип події (ERROR, WARNING, INFO, STATUS), код події (ціле).

Білі символи навколо роздільників не допускаються.

Приклад такого рядка присвоєно змінній \kw{s}.



\begin{verbatim}
s = 'username:2022-05-05 13-26:ERROR:404'
\end{verbatim}


Нижче наведено код, за виконання якого значенням змінної \kw{out} має стати код події.

Для наведеного прикладу значенням змінної \kw{out} має стати рядок \code{404}



\begin{verbatim}
res = re.fullmatch(r'XXX', s)
s = ''
out = YYY
\end{verbatim}


Код має працювати не тільки для конкретного рядка з прикладу, але й для кожного рядка з журналу.


Який рядок треба записати на місце \kw{XXX} ?
Який рядок треба записати на місце \kw{YYY} ?\par

%\end {large}
}
\newpage

{{23.05.2025 Контрольна робота, варіант  \No \textbf { 54 }} група К1{\parbox {0.5 cm} \dotfill}  ПІБ {\parbox {2.5 cm} \dotfill}\par}
{
%\begin {large}
\begin{minipage}[t]{9.7cm}
               
{\begin{center}Завдання {1}\end{center}}
\vspace{-6mm}
Укажіть ТИП значення виразу.\par\par
\begin{verbatim}
[13,2,4]+[3,5]
\end{verbatim}

\vspace{15mm}
У завданнях 2, 3 вказати, що буде виведено за виконання. Повідомлення інтерпретатора про помилку позначати ERROR

               
{\begin{center}Завдання {2}\end{center}}
\vspace{-6mm}

\begin{verbatim}
try:
    t = {27:1, 26:4, 63:7, 50:2, 44:9}
    t[50] = 6
    for x in t.items():
        print(x, sep=' ', end=' ')
except:
    print('ERR')
\end{verbatim}


\end{minipage}
\vline \hspace{8mm}
\begin{minipage}[t]{8.0cm}                
               
{\begin{center}Завдання {3}\end{center}}
\vspace{-6mm}

\begin{verbatim}
c = 0
j = 1

class A:
    c = 2
    
    def __init__(self):
        self.j = 3
        c = 4
        A.c = 5

obj = A()
try:
    print(c, end=' ')
    print(j, end=' ')
    print(A.c, end=' ')
    print(obj.j)
except:
    print('ERR')
\end{verbatim}

\end{minipage}

               
{\begin{center}Завдання {4}\end{center}}
\vspace{-6mm}
Нехай змінна \kw{s} позначає дуже довгу послідовність цілих, по якій можна ітерувати.

Написати вираз-генератор, що задає підпослідовність її елементів, що є додатни\-ми.

Обмеження: усі елементи шуканої підпослідовності одночасно в пам'яті розміщені бути не можуть.\par

\vspace{10mm}
               
{\begin{center}Завдання {5}\end{center}}
\vspace{-6mm}
Змінні \kw{next} та \kw{prev} вузла списку позначають наступний та попередній за ним вузол відповідно. Починаючи з голови, у список записано послідовні цілі числа від~$-38$ до~$46$. Нехай змінна \kw{p} позначає вузол, що зберігає значення~$-8$.
Яке число зберігається у вузлі \kw{p.next.next.prev.prev.prev} ?\par

               
{\begin{center}Завдання {6}\end{center}}
\vspace{-6mm}
Значенням змінної \kw{a} є цілочисловий масив типу $numpy.ndarray$ розмірів $6{\times}6$. На скільки збільшиться сума елементів масиву \kw{a} після виконання
\begin{verbatim}
a[1] += 2
\end{verbatim}


Якщо код некоректний, то в якості відповіді зазначити ERROR.\par

               
{\begin{center}Завдання {10}\end{center}}
\vspace{-6mm}

Написати функцію $f$, яка для файлу, заданого шляхом, розв’язує задачу:\\ Текстовий файл містить послідовність чисел $a_1$, $a_2$, $\ldots$, $a_t$, записану у форматі одне число на рядок. Знайти третє найбільше значення, що записане у файлі. Кожне значення, що записане у файлі, враховується стільки разів, скільки входить у файл. Для послідовності 2, 4, 12, 5, 2, 0 таким значенням буде 2.

В усіх випадках розміри файлів такі, що їх вміст повністю завантажений у пам'ять бути не може. Якість коду також оцінюється.



\vspace{10mm}\par

\newpage
\begin{minipage}[t]{9.7cm}
               
{\begin{center}Завдання {7}\end{center}}
\vspace{-6mm}

\begin{center}
	\adjustimage{max size={0.8\linewidth}{0.8\paperheight}}
	{./15BinTree/trees_exp/24.png}
\end{center}
Указати послідовність міток вершин дерева, зображеного на рис., що утвориться в результаті обходу дерева в прямому порядку. Мітки записувати через пробіл.\par Алгоритм починає роботу з кореня (на рис. це верхня вершина).\par

\end{minipage}
\vline \hspace{8mm}
\begin{minipage}[t]{8.0cm}
               
{\begin{center}Завдання {8}\end{center}}
\vspace{-6mm}
Для графа на рисунку з вершини 2 запущено алгоритм пошуку в глибину, що будує кістякове дерево. Списки суміжності вершин переглядаються алгоритмом за зростанням міток вершин.\par
\begin{center}
	\adjustimage{max size={0.9\linewidth}{0.9\paperheight}}
	{./16Graphs/Traversals/310.png}
\end{center}
\par Які з наведених ребер входять до складу отриманого кістякового дерева?\par
1) (3, 4)\par
2) (0, 6)\par
3) (0, 3)\par
4) (0, 2)\par
5) (0, 5)\par
6) (3, 5)\par
{\vskip 15pt} \par

\end{minipage}

               
{\begin{center}Завдання {9}\end{center}}
\vspace{-6mm}
Файл \code{'1.txt'} містить відомість з рядками, в яких через кому записано поряд\-ковий номер (ціле), назву предмету (знаків пунктуації не містить), прізвище студента, ім'я студента, номер заліковки, оцінку в балах за 100 бальною шкалою.

Білі символи навколо роздільників не допускаються.

Приклад рядка файлу



\begin{verbatim}
13,algebra,surname,name,123456,77
\end{verbatim}


Нижче наведено код, який має в кожному рядку файлу переставити місцями назву предмету та номер заліковки.

Рядок з прикладу має бути замінений на такий



\begin{verbatim}
13,123456,surname,name,algebra,77
\end{verbatim}




Код



\begin{verbatim}
regex = re.compile(r'XXX')
with open('1.txt', encoding='utf8') as f:
    lst = f.readlines()
for i in range(len(lst)):
    lst[i] = re.sub(regex, YYY, lst[i])
with open('1.txt', 'w', encoding='utf8') as f:
    f.writelines(lst)
\end{verbatim}





Який рядок треба записати на місце \kw{XXX} ?
Який рядок треба записати на місце \kw{YYY} ?\par

%\end {large}
}
\newpage

{{23.05.2025 Контрольна робота, варіант  \No \textbf { 55 }} група К1{\parbox {0.5 cm} \dotfill}  ПІБ {\parbox {2.5 cm} \dotfill}\par}
{
%\begin {large}
\begin{minipage}[t]{9.7cm}
               
{\begin{center}Завдання {1}\end{center}}
\vspace{-6mm}
Укажіть ТИП значення виразу.\par\par
\begin{verbatim}
{x*x:x for x in range(-7,7)}
\end{verbatim}

\vspace{15mm}
У завданнях 2, 3 вказати, що буде виведено за виконання. Повідомлення інтерпретатора про помилку позначати ERROR

               
{\begin{center}Завдання {2}\end{center}}
\vspace{-6mm}

\begin{verbatim}
try:
    s = {87:0, 35:7, 54:9, 87:9}
    s[39] = 7
    for x in s.items():
        print(x, sep=' ', end=' ')
except:
    print('ERR')
\end{verbatim}


\end{minipage}
\vline \hspace{8mm}
\begin{minipage}[t]{8.0cm}                
               
{\begin{center}Завдання {3}\end{center}}
\vspace{-6mm}

\begin{verbatim}
d = 0
m = 1

class D:
    m = 2
    
    def __init__(self):
        global m
        self.d = 3
        m = 4
        D.m = 5

obj = D()
try:
    print(d, end=' ')
    print(m, end=' ')
    print(D.m, end=' ')
    print(obj.d)
except:
    print('ERR')
\end{verbatim}

\end{minipage}

               
{\begin{center}Завдання {4}\end{center}}
\vspace{-6mm}
Нехай змінна \kw{s} позначає дуже довгу послідовність, по якій можна ітерувати.

Написати вираз-генератор, що задає підпослідовність її елементів, що не є рядками типу \kw{str}.

Обмеження: усі елементи шуканої підпослідовності одночасно в пам'яті розміщені бути не можуть.\par

\vspace{10mm}
               
{\begin{center}Завдання {5}\end{center}}
\vspace{-6mm}
Змінні \kw{next} та \kw{prev} вузла списку позначають наступний та попередній за ним вузол відповідно. Починаючи з голови, у список записано послідовні цілі числа від~$-37$ до~$50$. Нехай змінна \kw{p} позначає вузол, що зберігає значення~$1$.
Яке число зберігається у вузлі \kw{p.prev.prev.next.next.prev} ?\par

               
{\begin{center}Завдання {6}\end{center}}
\vspace{-6mm}
Значенням змінної \kw{a} є цілочисловий масив типу $numpy.ndarray$ розмірів $4{\times}3$. На скільки збільшиться сума елементів масиву \kw{a} після виконання
\begin{verbatim}
a[-3] += 1
\end{verbatim}


Якщо код некоректний, то в якості відповіді зазначити ERROR.\par

               
{\begin{center}Завдання {10}\end{center}}
\vspace{-6mm}

Написати функцію $f$, яка для файлу, заданого шляхом, розв’язує задачу:\\ Кожен з двох текстових файлів містить послідовність чисел, записану у форматі одне число на рядок.  Перевірити, чи задають вони одну ту саму послідовність.

В усіх випадках розміри файлів такі, що їх вміст повністю завантажений у пам'ять бути не може. Якість коду також оцінюється.



\vspace{10mm}\par

\newpage
\begin{minipage}[t]{9.7cm}
               
{\begin{center}Завдання {7}\end{center}}
\vspace{-6mm}

\begin{center}
	\adjustimage{max size={0.8\linewidth}{0.8\paperheight}}
	{./15BinTree/trees_exp/64.png}
\end{center}
Указати послідовність міток вершин дерева, зображеного на рис., що утвориться в результаті обходу дерева в прямому порядку. Мітки записувати через пробіл.\par Алгоритм починає роботу з кореня (на рис. це верхня вершина).\par

\end{minipage}
\vline \hspace{8mm}
\begin{minipage}[t]{8.0cm}
               
{\begin{center}Завдання {8}\end{center}}
\vspace{-6mm}
Для графа на рисунку з вершини 3 запущено алгоритм пошуку в ширину, що будує кістякове дерево. Списки суміжності вершин переглядаються алгоритмом за зростанням міток вершин.\par
\begin{center}
	\adjustimage{max size={0.9\linewidth}{0.9\paperheight}}
	{./16Graphs/Traversals/288.png}
\end{center}
\par Які з наведених ребер входять до складу отриманого кістякового дерева?\par
1) (3, 6)\par
2) (0, 4)\par
3) (0, 5)\par
4) (0, 6)\par
5) (1, 2)\par
6) (1, 3)\par
{\vskip 15pt} \par

\end{minipage}

               
{\begin{center}Завдання {9}\end{center}}
\vspace{-6mm}
Рядки журналу через кому містять ім'я користувача (складається з латиниці), час події (фор\-мат як у прикладі), тип події (ERROR, WARNING, INFO, STATUS), код події (ціле).

Білі символи навколо роздільників не допускаються.

Приклад такого рядка присвоєно змінній \kw{s}.



\begin{verbatim}
s = 'username,2022-05-05 13:26,ERROR,404'
\end{verbatim}


Нижче наведено код, за виконання якого значенням змінної \kw{out} має стати код події.

Для наведеного прикладу значенням змінної \kw{out} має стати рядок \code{404}



\begin{verbatim}
res = re.fullmatch(r'XXX', s)
s = ''
out = YYY
\end{verbatim}


Код має працювати не тільки для конкретного рядка з прикладу, але й для кожного рядка з журналу.


Який рядок треба записати на місце \kw{XXX} ?
Який рядок треба записати на місце \kw{YYY} ?\par

%\end {large}
}
\newpage

{{23.05.2025 Контрольна робота, варіант  \No \textbf { 56 }} група К1{\parbox {0.5 cm} \dotfill}  ПІБ {\parbox {2.5 cm} \dotfill}\par}
{
%\begin {large}
\begin{minipage}[t]{9.7cm}
               
{\begin{center}Завдання {1}\end{center}}
\vspace{-6mm}
Укажіть ТИП значення виразу.\par\par
\begin{verbatim}
[{1,2}, 5]
\end{verbatim}

\vspace{15mm}
У завданнях 2, 3 вказати, що буде виведено за виконання. Повідомлення інтерпретатора про помилку позначати ERROR

               
{\begin{center}Завдання {2}\end{center}}
\vspace{-6mm}

\begin{verbatim}
try:
    t = {64:4, 67:8, 83:7, 50:8, 83:5}
    t[38] = 7
    for x in t.values():
        print(x, sep=' ', end=' ')
except:
    print('ERR')
\end{verbatim}


\end{minipage}
\vline \hspace{8mm}
\begin{minipage}[t]{8.0cm}                
               
{\begin{center}Завдання {3}\end{center}}
\vspace{-6mm}

\begin{verbatim}
a = 0
h = 1

class B:
    a = 2
    
    def __init__(self):
        self.h = 3
        a = 4
        B.a = 5

obj = B()
try:
    print(a, end=' ')
    print(h, end=' ')
    print(B.a, end=' ')
    print(obj.h)
except:
    print('ERR')
\end{verbatim}

\end{minipage}

               
{\begin{center}Завдання {4}\end{center}}
\vspace{-6mm}
Нехай змінна \kw{s} позначає дуже довгу послідовність цілих, по якій можна ітерувати.

Написати вираз-генератор, що задає підпослідовність її елементів, що більші за задане число num.

Обмеження: усі елементи шуканої підпослідовності одночасно в пам'яті розміщені бути не можуть.\par

\vspace{10mm}
               
{\begin{center}Завдання {5}\end{center}}
\vspace{-6mm}
Змінні \kw{next} та \kw{prev} вузла списку позначають наступний та попередній за ним вузол відповідно. Починаючи з голови, у список записано послідовні цілі числа від~$-17$ до~$38$. Нехай змінна \kw{p} позначає вузол, що зберігає значення~$-4$.
Яке число зберігається у вузлі \kw{p.next.next.prev.prev.next} ?\par

               
{\begin{center}Завдання {6}\end{center}}
\vspace{-6mm}
Значенням змінної \kw{a} є цілочисловий масив типу $numpy.ndarray$ розмірів $7{\times}4$. На скільки збільшиться сума елементів масиву \kw{a} після виконання
\begin{verbatim}
a.T[-1] += 2
\end{verbatim}


Якщо код некоректний, то в якості відповіді зазначити ERROR.\par

               
{\begin{center}Завдання {10}\end{center}}
\vspace{-6mm}

Написати функцію $f$, яка для файлу, заданого шляхом, розв’язує задачу:\\ Текстовий файл містить послідовність чисел $a_1$, $a_2$, $\ldots$, $a_t$, записану у форматі одне число на рядок. Знайти середнє арифметичне чисел $a_1$, $a_4$, $a_7$, $\ldots$ (починаючи з першого через два).

В усіх випадках розміри файлів такі, що їх вміст повністю завантажений у пам'ять бути не може. Якість коду також оцінюється.



\vspace{10mm}\par

\newpage
\begin{minipage}[t]{9.7cm}
               
{\begin{center}Завдання {7}\end{center}}
\vspace{-6mm}

\begin{center}
	\adjustimage{max size={0.8\linewidth}{0.8\paperheight}}
	{./15BinTree/trees_exp/123.png}
\end{center}
Указати послідовність міток вершин дерева, зображеного на рис., що утвориться в результаті обходу дерева в оберненому порядку. Мітки записувати через пробіл.\par Алгоритм починає роботу з кореня (на рис. це верхня вершина).\par

\end{minipage}
\vline \hspace{8mm}
\begin{minipage}[t]{8.0cm}
               
{\begin{center}Завдання {8}\end{center}}
\vspace{-6mm}
Для графа на рисунку з вершини 2 запущено алгоритм пошуку в глибину, що будує кістякове дерево. Списки суміжності вершин переглядаються алгоритмом за зростанням міток вершин.\par
\begin{center}
	\adjustimage{max size={0.9\linewidth}{0.9\paperheight}}
	{./16Graphs/Traversals/322.png}
\end{center}
\par Які з наведених ребер входять до складу отриманого кістякового дерева?\par
1) (2, 5)\par
2) (3, 4)\par
3) (1, 2)\par
4) (0, 5)\par
5) (5, 6)\par
6) (1, 5)\par
{\vskip 15pt} \par

\end{minipage}

               
{\begin{center}Завдання {9}\end{center}}
\vspace{-6mm}
Файл \code{'1.txt'} містить відомість з рядками, в яких через кому записа\-но поряд\-ковий номер (ціле), назву предмету (знаків пунктуації не містить), прізвище студента, ім'я студента, номер заліковки, оцінку в балах за 100 бальною шкалою.

Білі символи навколо роздільників не допускаються.

Приклад рядка файлу



\begin{verbatim}
13,algebra,surname,name,123456,77
\end{verbatim}


Нижче наведено код, який має залишити у файлі в кожному рядку тільки номер заліковки, назву предмету та оцінку в балах за 100-бальною шкалою (у заданій послі\-дов\-ності).

Рядок з прикладу має бути замінений на такий



\begin{verbatim}
123456,algebra,77
\end{verbatim}




Код



\begin{verbatim}
regex = re.compile(r'XXX')
with open('1.txt', encoding='utf8') as f:
    lst = f.readlines()
for i in range(len(lst)):
    lst[i] = re.sub(regex, YYY, lst[i])
with open('1.txt', 'w', encoding='utf8') as f:
    f.writelines(lst)
\end{verbatim}





Який рядок треба записати на місце \kw{XXX} ?
Який рядок треба записати на місце \kw{YYY} ?\par

%\end {large}
}
\newpage

{{23.05.2025 Контрольна робота, варіант  \No \textbf { 57 }} група К1{\parbox {0.5 cm} \dotfill}  ПІБ {\parbox {2.5 cm} \dotfill}\par}
{
%\begin {large}
\begin{minipage}[t]{9.7cm}
               
{\begin{center}Завдання {1}\end{center}}
\vspace{-6mm}
Укажіть ТИП значення виразу.\par\par
\begin{verbatim}
[7,2,3,4]
\end{verbatim}

\vspace{15mm}
У завданнях 2, 3 вказати, що буде виведено за виконання. Повідомлення інтерпретатора про помилку позначати ERROR

               
{\begin{center}Завдання {2}\end{center}}
\vspace{-6mm}

\begin{verbatim}
try:
    d = {39:5, 22:5, 70:3, 70:1}
    d[66] = 8
    for x in d :
        print(x, sep=' ', end=' ')
except:
    print('ERR')
\end{verbatim}


\end{minipage}
\vline \hspace{8mm}
\begin{minipage}[t]{8.0cm}                
               
{\begin{center}Завдання {3}\end{center}}
\vspace{-6mm}

\begin{verbatim}
b = 0
k = 1

class A:
    b = 2
    
    def __init__(self):
        global b
        self.k = 3
        b = 4
        A.b = 5

obj = A()
try:
    print(b, end=' ')
    print(k, end=' ')
    print(A.b, end=' ')
    print(obj.k)
except:
    print('ERR')
\end{verbatim}

\end{minipage}

               
{\begin{center}Завдання {4}\end{center}}
\vspace{-6mm}
Нехай змінна \kw{s} позначає послідовність цілих, по якій можна ітерувати.

Написати вираз, що перевіряє, чи правда, що всі елементи послідовності більші за задане число num.\par

\vspace{10mm}
               
{\begin{center}Завдання {5}\end{center}}
\vspace{-6mm}
Змінні \kw{next} та \kw{prev} вузла списку позначають наступний та попередній за ним вузол відповідно. Починаючи з голови, у список записано послідовні цілі числа від~$-30$ до~$30$. Нехай змінна \kw{p} позначає вузол, що зберігає значення~$-1$.
Яке число зберігається у вузлі \kw{p.prev.next.prev.prev.next} ?\par

               
{\begin{center}Завдання {6}\end{center}}
\vspace{-6mm}
Значенням змінної \kw{a} є цілочисловий масив типу $numpy.ndarray$ розмірів $6{\times}5$. На скільки збільшиться сума елементів масиву \kw{a} після виконання
\begin{verbatim}
a.T[-2] += 3
\end{verbatim}


Якщо код некоректний, то в якості відповіді зазначити ERROR.\par

               
{\begin{center}Завдання {10}\end{center}}
\vspace{-6mm}

Написати функцію $f$, яка для файлу, заданого шляхом, розв’язує задачу:\\ Текстовий файл містить послідовність чисел $a_1$, $a_2$, $\ldots$, $a_t$, записану у форматі одне число на рядок. Знайти третє найменше значення, що записане у файлі. Кожне значення, що записане у файлі, враховується тільки один раз. \par
Для послідовності 1, 4, 12, 5, 1, 0 таким значенням буде 4.

В усіх випадках розміри файлів такі, що їх вміст повністю завантажений у пам'ять бути не може. Якість коду також оцінюється.



\vspace{10mm}\par

\newpage
\begin{minipage}[t]{9.7cm}
               
{\begin{center}Завдання {7}\end{center}}
\vspace{-6mm}

\begin{center}
	\adjustimage{max size={0.8\linewidth}{0.8\paperheight}}
	{./15BinTree/trees_exp/143.png}
\end{center}
Указати послідовність міток вершин дерева, зображеного на рис., що утвориться в результаті обходу дерева в оберненому порядку. Мітки записувати через пробіл.\par Алгоритм починає роботу з кореня (на рис. це верхня вершина).\par

\end{minipage}
\vline \hspace{8mm}
\begin{minipage}[t]{8.0cm}
               
{\begin{center}Завдання {8}\end{center}}
\vspace{-6mm}
Для графа на рисунку з вершини 4 запущено алгоритм пошуку в ширину, що будує кістякове дерево. Списки суміжності вершин переглядаються алгоритмом за зростанням міток вершин.\par
\begin{center}
	\adjustimage{max size={0.9\linewidth}{0.9\paperheight}}
	{./16Graphs/Traversals/90.png}
\end{center}
\par Які з наведених ребер входять до складу отриманого кістякового дерева?\par
1) (0, 6)\par
2) (3, 6)\par
3) (0, 3)\par
4) (4, 5)\par
5) (1, 5)\par
6) (2, 5)\par
{\vskip 15pt} \par

\end{minipage}

               
{\begin{center}Завдання {9}\end{center}}
\vspace{-6mm}
Файл \code{'1.txt'} містить відомість з рядками, в яких через кому записа\-но назву предмету (знаків пунктуації не містить), прізвище студента, ім'я студента, номер залі\-ков\-ки, оцінку в балах за 100 бальною шкалою.

Білі символи навколо роздільників не допускаються.

Приклад рядка файлу



\begin{verbatim}
algebra,surname,name,123456,77
\end{verbatim}


Нижче наведено код, який має замінити у файлі усі номери заліковок на рядок \code{***}.



\begin{verbatim}
regex = re.compile(r'XXX')
with open('1.txt', encoding='utf8') as f:
    lst = f.readlines()
for i in range(len(lst)):
    lst[i] = re.sub(regex, YYY, lst[i])
with open('1.txt', 'w', encoding='utf8') as f:
    f.writelines(lst)
\end{verbatim}





Який рядок треба записати на місце \kw{XXX} ?
Який рядок треба записати на місце \kw{YYY} ?\par

%\end {large}
}
\newpage

{{23.05.2025 Контрольна робота, варіант  \No \textbf { 58 }} група К1{\parbox {0.5 cm} \dotfill}  ПІБ {\parbox {2.5 cm} \dotfill}\par}
{
%\begin {large}
\begin{minipage}[t]{9.7cm}
               
{\begin{center}Завдання {1}\end{center}}
\vspace{-6mm}
Укажіть ТИП значення виразу.\par\par
\begin{verbatim}
{x for x in range(7) if x%3==1}
\end{verbatim}

\vspace{15mm}
У завданнях 2, 3 вказати, що буде виведено за виконання. Повідомлення інтерпретатора про помилку позначати ERROR

               
{\begin{center}Завдання {2}\end{center}}
\vspace{-6mm}

\begin{verbatim}
try:
    d = {79:2, 87:7, 83:6, 47:7, 83:3}
    d[79] = 9
    for x in d.items():
        print(*x, sep=' ', end=' ')
except:
    print('ERR')
\end{verbatim}


\end{minipage}
\vline \hspace{8mm}
\begin{minipage}[t]{8.0cm}                
               
{\begin{center}Завдання {3}\end{center}}
\vspace{-6mm}

\begin{verbatim}
c = 0
n = 1

class C:
    n = 2
    
    def __init__(self):
        self.c = 3
        n = 4
        C.n = 5

obj = C()
try:
    print(c, end=' ')
    print(n, end=' ')
    print(C.n, end=' ')
    print(obj.c)
except:
    print('ERR')
\end{verbatim}

\end{minipage}

               
{\begin{center}Завдання {4}\end{center}}
\vspace{-6mm}
Нехай змінна \kw{s} позначає послідовність цілих, по якій можна ітерувати.

Написати вираз, що перевіряє, чи правда, що послідовність  містить хоча б одне додатнє число.\par

\vspace{10mm}
               
{\begin{center}Завдання {5}\end{center}}
\vspace{-6mm}
Змінні \kw{next} та \kw{prev} вузла списку позначають наступний та попередній за ним вузол відповідно. Починаючи з голови, у список записано послідовні цілі числа від~$-31$ до~$36$. Нехай змінна \kw{p} позначає вузол, що зберігає значення~$0$.
Яке число зберігається у вузлі \kw{p.next.prev.next.next.prev} ?\par

               
{\begin{center}Завдання {6}\end{center}}
\vspace{-6mm}
Значенням змінної \kw{a} є цілочисловий масив типу $numpy.ndarray$ розмірів $3{\times}7$. На скільки збільшиться сума елементів масиву \kw{a} після виконання
\begin{verbatim}
a[2] += 1
\end{verbatim}


Якщо код некоректний, то в якості відповіді зазначити ERROR.\par

               
{\begin{center}Завдання {10}\end{center}}
\vspace{-6mm}

Написати функцію $f$, яка для файлу, заданого шляхом, розв’язує задачу:\\ Текстовий файл містить послідовність чисел $a_1$, $a_2$, $\ldots$, $a_t$, записану у форматі одне число на рядок. Знайти третє найбільше значення, що записане у файлі. Кожне значення, що записане у файлі, враховується стільки разів, скільки входить у файл. Для послідовності 2, 4, 12, 5, 2, 0 таким значенням буде 2.

В усіх випадках розміри файлів такі, що їх вміст повністю завантажений у пам'ять бути не може. Якість коду також оцінюється.



\vspace{10mm}\par

\newpage
\begin{minipage}[t]{9.7cm}
               
{\begin{center}Завдання {7}\end{center}}
\vspace{-6mm}

\begin{center}
	\adjustimage{max size={0.8\linewidth}{0.8\paperheight}}
	{./15BinTree/trees_exp/96.png}
\end{center}
Указати послідовність міток вершин дерева, зображеного на рис., що утвориться в результаті обходу дерева в прямому порядку. Мітки записувати через пробіл.\par Алгоритм починає роботу з кореня (на рис. це верхня вершина).\par

\end{minipage}
\vline \hspace{8mm}
\begin{minipage}[t]{8.0cm}
               
{\begin{center}Завдання {8}\end{center}}
\vspace{-6mm}
Для графа на рисунку з вершини 0 запущено алгоритм пошуку в ширину, що будує кістякове дерево. Списки суміжності вершин переглядаються алгоритмом за зростанням міток вершин.\par
\begin{center}
	\adjustimage{max size={0.9\linewidth}{0.9\paperheight}}
	{./16Graphs/Traversals/179.png}
\end{center}
\par Які з наведених ребер входять до складу отриманого кістякового дерева?\par
1) (1, 6)\par
2) (0, 2)\par
3) (2, 4)\par
4) (4, 5)\par
5) (0, 4)\par
6) (0, 6)\par
{\vskip 15pt} \par

\end{minipage}

               
{\begin{center}Завдання {9}\end{center}}
\vspace{-6mm}
Файл \code{'1.txt'} містить відомість з рядками, в яких через кому записано поряд\-ковий номер (ціле), назву предмету (знаків пунктуації не містить), прізвище студента, ім'я студента, номер заліковки, оцінку в балах за 100 бальною шкалою.

Білі символи навколо роздільників не допускаються.

Приклад рядка файлу



\begin{verbatim}
13,algebra,surname,name,123456,77
\end{verbatim}


Нижче наведено код, який шукає усі предмети, який складав студент із заліковкою \code{777}, та зберігає їх типи в змінній \kw{codes}.



\begin{verbatim}
codes = set()
regex = re.compile(r'XXX')
with open('1.txt', encoding='utf8') as f:
    for s in f:
        res = re.fullmatch(regex, s)
        if not res: continue
        s = ''
        codes.add(YYY)
\end{verbatim}



Який рядок треба записати на місце \kw{XXX} ?
Який рядок треба записати на місце \kw{YYY} ?\par

%\end {large}
}
\newpage

{{23.05.2025 Контрольна робота, варіант  \No \textbf { 59 }} група К1{\parbox {0.5 cm} \dotfill}  ПІБ {\parbox {2.5 cm} \dotfill}\par}
{
%\begin {large}
\begin{minipage}[t]{9.7cm}
               
{\begin{center}Завдання {1}\end{center}}
\vspace{-6mm}
Укажіть ТИП значення виразу.\par\par
\begin{verbatim}
{(1,4), 8}
\end{verbatim}

\vspace{15mm}
У завданнях 2, 3 вказати, що буде виведено за виконання. Повідомлення інтерпретатора про помилку позначати ERROR

               
{\begin{center}Завдання {2}\end{center}}
\vspace{-6mm}

\begin{verbatim}
try:
    t = {15:5, 23:0, 22:3, 22:3}
    t[15] = 6
    for x in t.keys():
        print(x, sep=' ', end=' ')
except:
    print('ERR')
\end{verbatim}


\end{minipage}
\vline \hspace{8mm}
\begin{minipage}[t]{8.0cm}                
               
{\begin{center}Завдання {3}\end{center}}
\vspace{-6mm}

\begin{verbatim}
_c = 0

class D:
    _c = 1
    
    def __init__(self):
        self._c = 2
        _c = 3
        D._c = 4

obj = D()
D._c = 5
try:
    print(_c, end=' ')
    print(obj._c, end=' ')
    print(D._c, end=' ')
    print(obj._c, end=' ')
    print(obj._D__c)
except:
    print('ERR')
\end{verbatim}

\end{minipage}

               
{\begin{center}Завдання {4}\end{center}}
\vspace{-6mm}
Нехай змінна \kw{s} позначає послідовність цілих, по якій можна ітерувати.

Написати вираз, що перевіряє, чи правда, що всі елементи послідовності є від’ємни\-ми.\par

\vspace{10mm}
               
{\begin{center}Завдання {5}\end{center}}
\vspace{-6mm}
Змінні \kw{next} та \kw{prev} вузла списку позначають наступний та попередній за ним вузол відповідно. Починаючи з голови, у список записано послідовні цілі числа від~$-24$ до~$15$. Нехай змінна \kw{p} позначає вузол, що зберігає значення~$-3$.
Яке число зберігається у вузлі \kw{p.next.next.next.next.next} ?\par

               
{\begin{center}Завдання {6}\end{center}}
\vspace{-6mm}
Значенням змінної \kw{a} є цілочисловий масив типу $numpy.ndarray$ розмірів $5{\times}6$. На скільки збільшиться сума елементів масиву \kw{a} після виконання
\begin{verbatim}
a[0, 1] += 3
\end{verbatim}


Якщо код некоректний, то в якості відповіді зазначити ERROR.\par

               
{\begin{center}Завдання {10}\end{center}}
\vspace{-6mm}

Написати функцію $f$, яка для файлу, заданого шляхом, розв’язує задачу:\\ Текстовий файл містить послідовність чисел $a_1$, $a_2$, $\ldots$, $a_t$, записану у форматі одне число на рядок. Знайти третє найменше значення, що записане у файлі. Кожне значення, що записане у файлі, враховується стільки разів, скільки входить у файл. Для послідовності 2, 4, 12, 5, 2, 0 таким значенням буде 2.

В усіх випадках розміри файлів такі, що їх вміст повністю завантажений у пам'ять бути не може. Якість коду також оцінюється.



\vspace{10mm}\par

\newpage
\begin{minipage}[t]{9.7cm}
               
{\begin{center}Завдання {7}\end{center}}
\vspace{-6mm}

\begin{center}
	\adjustimage{max size={0.8\linewidth}{0.8\paperheight}}
	{./15BinTree/trees_exp/142.png}
\end{center}
Указати послідовність міток вершин дерева, зображеного на рис., що утвориться в результаті обходу дерева в оберненому порядку. Мітки записувати через пробіл.\par Алгоритм починає роботу з кореня (на рис. це верхня вершина).\par

\end{minipage}
\vline \hspace{8mm}
\begin{minipage}[t]{8.0cm}
               
{\begin{center}Завдання {8}\end{center}}
\vspace{-6mm}
Для графа на рисунку з вершини 4 запущено алгоритм пошуку в глибину, що будує кістякове дерево. Списки суміжності вершин переглядаються алгоритмом за зростанням міток вершин.\par
\begin{center}
	\adjustimage{max size={0.9\linewidth}{0.9\paperheight}}
	{./16Graphs/Traversals/11.png}
\end{center}
\par Які з наведених ребер входять до складу отриманого кістякового дерева?\par
1) (1, 5)\par
2) (2, 4)\par
3) (3, 5)\par
4) (0, 6)\par
5) (1, 4)\par
6) (2, 6)\par
{\vskip 15pt} \par

\end{minipage}

               
{\begin{center}Завдання {9}\end{center}}
\vspace{-6mm}
Файл \code{'1.txt'} містить журнал з рядками, в яких через кому записано поряд\-ко\-вий номер запису в журналі (ціле), тип події (ERROR, WARNING, INFO, STATUS), ім'я користувача (складається з латиниці), час події (фор\-мат як у прикладі), код події (ціле).

Білі символи навколо роздільників не допускаються.

Приклад рядка файлу



\begin{verbatim}
13,ERROR,username,2022-05-05 13:26,404
\end{verbatim}


Нижче наведено код, який шукає усі коди помилок (подія ERROR), що зустріча\-ють\-ся в журналі, та зберігає їх у змінній \kw{codes}.



\begin{verbatim}
codes = set()
regex = re.compile(r'XXX')
with open('1.txt', encoding='utf8') as f:
    for s in f:
        res = re.fullmatch(regex, s)
        if not res: continue
        s = ''
        codes.add(YYY)
\end{verbatim}



Який рядок треба записати на місце \kw{XXX} ?
Який рядок треба записати на місце \kw{YYY} ?\par

%\end {large}
}
\newpage

{{23.05.2025 Контрольна робота, варіант  \No \textbf { 60 }} група К1{\parbox {0.5 cm} \dotfill}  ПІБ {\parbox {2.5 cm} \dotfill}\par}
{
%\begin {large}
\begin{minipage}[t]{9.7cm}
               
{\begin{center}Завдання {1}\end{center}}
\vspace{-6mm}
Укажіть ТИП значення виразу.\par\par
\begin{verbatim}
{x:x for x in range(7)}
\end{verbatim}

\vspace{15mm}
У завданнях 2, 3 вказати, що буде виведено за виконання. Повідомлення інтерпретатора про помилку позначати ERROR

               
{\begin{center}Завдання {2}\end{center}}
\vspace{-6mm}

\begin{verbatim}
try:
    s = {81:9, 42:5, 32:0, 81:1}
    s[67] = 1
    for x, y in s.items():
        print(x, y, sep=' ', end=' ')
except:
    print('ERR')
\end{verbatim}


\end{minipage}
\vline \hspace{8mm}
\begin{minipage}[t]{8.0cm}                
               
{\begin{center}Завдання {3}\end{center}}
\vspace{-6mm}

\begin{verbatim}
b = 0

class C:
    b = 1
    
    def __init__(self):
        self.b = 2
        b = 3
        C.b = 4

obj = C()
obj.b = 5
try:
    print(b, end=' ')
    print(obj.b, end=' ')
    print(C.b, end=' ')
    print(obj.b, end=' ')
    print(obj._C__b)
except:
    print('ERR')
\end{verbatim}

\end{minipage}

               
{\begin{center}Завдання {4}\end{center}}
\vspace{-6mm}
Нехай змінна \kw{s} позначає дуже довгу послідовність, по якій можна ітерувати.

Написати вираз-генератор, що задає підпослідовність її елементів, що не є рядками типу \kw{str}.

Обмеження: усі елементи шуканої підпослідовності одночасно в пам'яті розміщені бути не можуть.\par

\vspace{10mm}
               
{\begin{center}Завдання {5}\end{center}}
\vspace{-6mm}
Змінні \kw{next} та \kw{prev} вузла списку позначають наступний та попередній за ним вузол відповідно. Починаючи з голови, у список записано послідовні цілі числа від~$-37$ до~$35$. Нехай змінна \kw{p} позначає вузол, що зберігає значення~$-7$.
Яке число зберігається у вузлі \kw{p.prev.prev.prev.prev.prev} ?\par

               
{\begin{center}Завдання {6}\end{center}}
\vspace{-6mm}
Значенням змінної \kw{a} є цілочисловий масив типу $numpy.ndarray$ розмірів $4{\times}3$. На скільки збільшиться сума елементів масиву \kw{a} після виконання
\begin{verbatim}
a.T[-3] += 2
\end{verbatim}


Якщо код некоректний, то в якості відповіді зазначити ERROR.\par

               
{\begin{center}Завдання {10}\end{center}}
\vspace{-6mm}

Написати функцію $f$, яка для файлу, заданого шляхом, розв’язує задачу:\\ Текстовий файл містить послідовність чисел $a_1$, $a_2$, $\ldots$, $a_t$, записану у форматі одне число на рядок. Знайти третє найбільше значення, що записане у файлі. Кожне значення, що записане у файлі, враховується тільки один раз. Для послідовності 1, 4, 12, 5, 1, 0 таким значенням буде 4.

В усіх випадках розміри файлів такі, що їх вміст повністю завантажений у пам'ять бути не може. Якість коду також оцінюється.



\vspace{10mm}\par

\newpage
\begin{minipage}[t]{9.7cm}
               
{\begin{center}Завдання {7}\end{center}}
\vspace{-6mm}

\begin{center}
	\adjustimage{max size={0.8\linewidth}{0.8\paperheight}}
	{./15BinTree/trees_exp/82.png}
\end{center}
Указати послідовність міток вершин дерева, зображеного на рис., що утвориться в результаті обходу дерева в прямому порядку. Мітки записувати через пробіл.\par Алгоритм починає роботу з кореня (на рис. це верхня вершина).\par

\end{minipage}
\vline \hspace{8mm}
\begin{minipage}[t]{8.0cm}
               
{\begin{center}Завдання {8}\end{center}}
\vspace{-6mm}
Для графа на рисунку з вершини 1 запущено алгоритм пошуку в глибину, що будує кістякове дерево. Списки суміжності вершин переглядаються алгоритмом за зростанням міток вершин.\par
\begin{center}
	\adjustimage{max size={0.9\linewidth}{0.9\paperheight}}
	{./16Graphs/Traversals/19.png}
\end{center}
\par Які з наведених ребер входять до складу отриманого кістякового дерева?\par
1) (0, 6)\par
2) (0, 4)\par
3) (3, 4)\par
4) (1, 3)\par
5) (1, 5)\par
6) (2, 5)\par
{\vskip 15pt} \par

\end{minipage}

               
{\begin{center}Завдання {9}\end{center}}
\vspace{-6mm}
Файл \code{'1.txt'} містить журнал з рядками, в яких через двокрапку записано поряд\-ко\-вий номер запису в журналі (ціле), тип події (ERROR, WARNING, INFO, STATUS), ім'я користувача (складається з латиниці), час події (фор\-мат як у прикладі), опис події.

Білі символи навколо роздільників не допускаються.

Приклад рядка файлу



\begin{verbatim}
13:ERROR:username:2022-05-05 13-26:something happens
\end{verbatim}


Нижче наведено код, який має залишити у файлі в кожному рядку тільки тип події та її опис.

Рядок з прикладу має бути замінений на такий



\begin{verbatim}
ERROR:something happens
\end{verbatim}




Код:



\begin{verbatim}
regex = re.compile(r'XXX')
with open('1.txt', encoding='utf8') as f:
    lst = f.readlines()
for i in range(len(lst)):
    lst[i] = re.sub(regex, YYY, lst[i])
with open('1.txt', 'w', encoding='utf8') as f:
    f.writelines(lst)
\end{verbatim}







Який рядок треба записати на місце \kw{XXX} ?
Який рядок треба записати на місце \kw{YYY} ?\par

%\end {large}
}
\newpage

%end
\end{document}

